%=====================================================================
%  CONCENTRICITY OVER THE OCTONIONS -- MASTER
%  Part 1: classical zeta facts (continuation, functional equation,
%          infinite Euler product, infinite Hadamard product) -- cited.
%  Part 2: slice-preserving theory, the ring R under the *-product, and the
%          equivalence/translation theorems (concentric zero-spheres <=> RH).
%  Part 3: the Concentricity Theorem. The base is
%          B = PGL(2,R) \ltimes OnePoint(R) and the sphere world SphereWorld.
%          Conditions C1--C3 determine one distinguished Mobius/exponential
%          element; orbit--stabilizer carries its whole O*-valued action
%          across the base. Its sphere projection is AsectionSlice, while
%          the value-bearing AsectionFunctor A : B -> Grpd has normalized
%          A-value states as fibres. Then T_A = Int_B A_A is formed. The
%          readout is pi_0(Int_B R_A) = colim_B(pi_0 . R_A). The colimit performs the
%          transport-generated identifications and yields one residue-C
%          real-value singleton {c}, c in R. Hence the residue-C
%          zero 6-spheres have common centre c and are concentric. Classical
%          translation and the numerical identification c=1/2 are downstream.
%=====================================================================

\documentclass[11pt]{article}

\usepackage[utf8]{inputenc}
\usepackage[margin=1in]{geometry}
\usepackage{amsmath,amssymb,amsthm,mathtools}
\usepackage{stmaryrd}   % \fatsemi — diagrammatic composition
\usepackage{amscd}
\usepackage{tikz-cd}
\usepackage{enumitem}
\usepackage{xcolor}
\usepackage[colorlinks=true,linkcolor=blue,citecolor=blue,urlcolor=blue]{hyperref}

% Vertical spacing.  Without \raggedbottom, a page that cannot fit an
% unbreakable display or tikzcd stretches the glue between its paragraphs to
% fill the height, which reads as a gap in the middle of the text.  This sends
% the slack to the foot of the page instead.
\raggedbottom
% Let displays sit closer to a page break before being pushed to the next page.
\predisplaypenalty=0
\postdisplaypenalty=0

%---------------------------------------------------------------------
%  Macros
%---------------------------------------------------------------------
\newcommand{\OO}{\mathbb{O}}
\newcommand{\CC}{\mathbb{C}}
\newcommand{\RR}{\mathbb{R}}
\newcommand{\HH}{\mathbb{H}}
\newcommand{\ZZ}{\mathbb{Z}}
\newcommand{\QQ}{\mathbb{Q}}
\newcommand{\NN}{\mathbb{N}}
\newcommand{\Ostar}{\OO^{*}}
\newcommand{\Cstar}{\CC^{*}}
\newcommand{\zetaO}{\zeta_{\Ostar}}
\newcommand{\zetaC}{\zeta_{\Cstar}}
\DeclareMathOperator{\Spec}{Spec}
\DeclareMathOperator{\Hom}{Hom}
\DeclareMathOperator{\Aut}{Aut}
\DeclareMathOperator{\re}{Re}
\DeclareMathOperator{\im}{Im}
\providecommand{\Gtwo}{G_2}

% Conspicuous markers (visible in compiled output):
\newcommand{\TODO}[1]{\medskip\noindent\textcolor{red}{\textbf{[TODO:} #1\textbf{]}}\medskip}
\newcommand{\PROGRAM}[1]{\medskip\noindent\textcolor{red}{\textbf{[PROGRAM:} #1\textbf{]}}\medskip}
\newcommand{\SOURCE}[1]{\textcolor{purple}{\small$\langle$source: #1$\rangle$}}
\newcommand{\stag}[1]{\href{https://stacks.math.columbia.edu/tag/#1}{\texttt{#1}}}
\providecommand{\Grpd}{\mathrm{Grpd}}
\providecommand{\SphereWorld}{\mathrm{SphereWorld}}

% Lean blueprint macros: no-ops when this file is compiled standalone; the
% leanblueprint package redefines them when the dependency graph and the Lean
% linkage are built. Keeping them here lets every statement carry its Lean
% declaration and its dependencies without any Lean names appearing in the prose.
\providecommand{\lean}[1]{}
\providecommand{\leanok}{}
\providecommand{\authorobject}[2]{}
\providecommand{\notready}{}
\providecommand{\mathlibok}{}
\providecommand{\uses}[1]{}
\providecommand{\proves}[1]{}

%---------------------------------------------------------------------
%  Theorem environments (document-wide numbering)
%---------------------------------------------------------------------
\theoremstyle{plain}
\newtheorem{theorem}{Theorem}[section]
\newtheorem{proposition}[theorem]{Proposition}
\newtheorem{lemma}[theorem]{Lemma}
\newtheorem{corollary}[theorem]{Corollary}
\theoremstyle{definition}
\newtheorem{definition}[theorem]{Definition}
\newtheorem{condition}[theorem]{Condition}
\theoremstyle{remark}
\newtheorem{remark}[theorem]{Remark}
\newtheorem{example}[theorem]{Example}

\title{Concentric Zero-Spheres of a class of sections of the commutative ring $\mathcal R$\\ of slice preserving functions
over the compactified Octonions\\[0.4em]
\large
with the Riemann Hypothesis as a corollary}
\author{Jesse Michael Paul\\[2pt]
{\normalsize PhD Candidate, Department of Mathematics \& Statistics}\\
{\normalsize University of North Carolina at Greensboro}}
\hypersetup{pdfauthor={Jesse Michael Paul}}
\date{June 2026}

\begin{document}
\maketitle

\begin{quote}\small\itshape
``\dots a `viewpoint' by itself remains fragmentary. It reveals to us one of the aspects of a
scenery or panorama, among a multiplicity of others which are equally valuable, equally `real'.
It is when complementary viewpoints of a common reality are conjugated, that is, when our `eyes'
are multiplied, that the gaze is able to penetrate further ahead in the reckoning of things.
\dots it also happens, sometimes, that a sheaf of viewpoints converging to a unique and vast
scenery gives rise to a novel thing; a thing which transcends each of the partial perspectives,
in the same way that a living being transcends each of its limbs and organs. This new thing could
be called a vision.''
\begin{flushright}\upshape --- A.\ Grothendieck, \emph{R\'ecoltes et Semailles}\end{flushright}
\end{quote}

\begin{quote}\small\itshape
``The measure of our success is whether what we do enables people to understand and think
more clearly and effectively about mathematics.''
\begin{flushright}\upshape --- W.\ Thurston, \emph{On Proof and Progress in Mathematics}\end{flushright}
\end{quote}



\begin{abstract}
Let $\mathcal R$ be the commutative ring of slice-preserving slice-regular functions on the
compactified octonions $\OO^{*}=S^8$, under the regular ($\ast$) product. Part~1 gathers the
classical facts about the Riemann zeta function. Part~2 develops the slice-preserving theory,
builds $\mathcal R$, places the octonionic zeta $\zetaO$ inside it, and proves the zero and
equivalence theorems: the classically nontrivial zeros of $\zeta$ are residue-$\CC$ zero
$6$-spheres, and the Riemann Hypothesis holds exactly when those spheres are concentric, sharing
a single real centre.

Part~3 proves the main result, the \emph{Concentricity Theorem}, for a section $A\in\mathcal R$
satisfying four conditions: \textup{(C1)} meromorphic continuation with a single simple pole,
whose pole-cancelled factor continues to a finite nonzero value at $p_A$;
\textup{(C2)} an infinite Euler product; \textup{(C3)} an infinite
slice-regular Weierstrass factorization whose spherical factors are the residue-$\CC$ zeros; and
\textup{(C4)} infinitely many such zeros. The theorem says the residue-$\CC$ zero $6$-spheres of
any such $A$ are concentric.

The proof builds a concentricity detector from the section itself, out of action groupoids,
matrix actions, and categorical homotopy theory. Conditions \textup{(C1)--(C3)} construct one
distinguished exponential M\"obius action on the projective base
$\mathcal B=PGL(2,\RR)\ltimes\operatorname{OnePoint}(\RR)$: the continued pole factor's
logarithm $\lambda_A$ is exponentiated in the diagonal matrix
$\operatorname{diag}(e^{\lambda_A},1)$ and conjugated into the slice M\"obius group. The
Euler product presents the action on the half-space; condition \textup{(C1)} continues the
pole-cancelled factor through the finite pole $p_A$, and the Weierstrass product presents that
same continued factor there. The slice-preserving normalization is $G_2$-equivariant on $\Ostar$, and
orbit--stabilizer proves that its object and arrow assignments define the sphere-world functor.
From then on that functor itself records the directions and concrete M\"obius matrix actions, while the A-section functor
$\mathcal A_A:\mathcal B\to\mathbf{Grpd}$ has the resulting normalized A-value states and their
transports as its fibres. Its Grothendieck construction is
$\mathcal T_A=\int_{\mathcal B}\mathcal A_A$.

The residue-$\CC$ system is the inverse-image action groupoid
$\mathcal R_A$ selected inside $\mathcal A_A$, with fully faithful inclusion
$\iota_A:\mathcal R_A\Rightarrow\mathcal A_A$, and hence a fully faithful inclusion of totals
$\int\iota_A:\int_{\mathcal B}\mathcal R_A\hookrightarrow\mathcal T_A$. For two arbitrary
objects of the inverse-image total, the A-specific projective/M\"obius transport compares their
positioned coordinates and the $G_2$ fibre arrow compares their slice directions. Hence this
inverse-image total is transitive, therefore connected.
Its component colimit is consequently a singleton. Applying $\pi_0$ fibrewise to
$\mathcal R_A$, the intrinsic real face of the same bound action descends to a natural
transformation $\pi_0\circ\mathcal R_A\Rightarrow\Delta(\RR)$. Its induced colimit map,
evaluated on the constructed zero-sphere representatives, reads one real value
$c\in\RR$ across the singleton component colimit. Thus every residue-$\CC$ zero
$6$-sphere has centre $c$. The theorem asserts the shared centre; naming the number is the task
of the corollaries.

Downstream corollaries translate residue-$\CC$ zeros to classically nontrivial zeros and, for
$\zetaO$, use its functional equation to identify $c=\tfrac12$: the Riemann Hypothesis.
\end{abstract}


\tableofcontents

%=====================================================================
\section*{Introduction}
%=====================================================================
The Riemann Hypothesis is usually met head-on, as a question about where the zeros of a single complex function fall along a line. This paper approaches it from the side. It asks instead whether a family of geometric objects, sitting a few dimensions up in the octonions, share a common centre, and shows that they must. The Riemann Hypothesis then follows, not as the thing we set out to prove, but as the shadow this larger fact casts back down onto the complex plane.

A nontrivial zero of the zeta function, seen from the complex line alone, is a lone point in a strip. Seen from the compactified octonions, that same zero opens into a six-dimensional sphere, and the critical-line question becomes a question of arrangement: are these spheres concentric, or scattered into different hyperplanes? It takes the whole continuum of imaginary directions in the octonions, read together, to see whether the alignment holds.

The paper is built to make that convergence visible and, at every step, checkable. Part 1 recalls the classical facts about the Riemann zeta function we will lean on. Part 2 develops the slice-preserving theory of functions on the octonions, builds the commutative ring $\mathcal R$ in which the octonionic zeta lives, and proves the reframing at the centre of everything: the classically nontrivial zeros are exactly the residue-$\CC$ zero six-spheres, and the Riemann Hypothesis holds precisely when those spheres are concentric. Part 3 proves the main result, the Concentricity Theorem, with a small piece of machinery I think of as a concentricity detector. From the section itself we build one distinguished action on a projective base; orbit--stabilizer proves its slice functor well defined, and the resulting concrete matrix actions let us read off, through categorical homotopy theory, whether the zero-spheres land in a single connected piece. They do; and for the octonionic zeta, that is the Riemann Hypothesis.

No one can picture an eight-dimensional graph in their head, but tools from categorical homotopy theory, like the total-object Grothendieck construction and standard colimit arguments, let us probe this structure and reason about it. The argument is also formalized in Lean, so that each step can be checked; the reader can not only verify the steps but, I hope, come to understand for themselves why the Concentricity Theorem is true and why the Riemann Hypothesis follows.

%=====================================================================
\part{Classical Background}
%=====================================================================
\section{The Riemann zeta function}

The Riemann zeta function is classical, and we take its basic theory as given. We recall only what
the construction uses: the meromorphic continuation and functional equation, the Euler product over
the primes, the Hadamard product over the zeros, and Hardy's theorem. We then read $\zeta$ in
compactified form, as a map to the Riemann sphere, which is the form in which it will extend to the
octonions.

\begin{theorem}[Meromorphic continuation and functional equation; Riemann \cite{Riemann1859}]\label{thm:riemann}
\lean{riemannZeta_meromorphicOn, riemannZeta_orderAt_one,
riemannZeta_conj, riemannZeta_one_sub_zero, nontrivialZero_re_mem_Ioo}
The Dirichlet series $\zeta(s)=\sum_{n\ge1}n^{-s}$, convergent for $\re s>1$, continues to a
meromorphic function on $\CC$: holomorphic as a $\CC$-valued function on
$\CC\setminus\{1\}$, with a single simple pole at $s=1$. Equivalently, read into the Riemann
sphere, the same function is a holomorphic map on the whole complex plane,
\[
  \zeta:\CC\longrightarrow\Cstar,\qquad \Cstar=\CC\cup\{\infty\}=S^2,\qquad \zeta(1)=\infty ,
\]
the pole becoming the regular value $\infty$. Holomorphy is on the domain $\CC$ in both
readings: with infinitely many zeros, $\zeta$ is not rational, so the domain point $\infty$
carries an essential singularity. (The later extension
of the \emph{domain} to $\Cstar$ assigns a marked value at $\infty$,
Definition~\ref{def:zeta-Cstar} with Remark~\ref{rmk:infty-marked}: an assignment of geometry
and value, made where the classical statement ends.)
The completed function $\xi(s)=\tfrac12 s(s-1)\pi^{-s/2}\Gamma(s/2)\zeta(s)$ is entire and satisfies
\[
  \xi(s)=\xi(1-s),\qquad \xi(\bar s)=\overline{\xi(s)} .
\]
Consequently, if $\rho$ is a nontrivial zero of $\zeta$ then so are $\bar\rho$, $1-\rho$, and
$1-\bar\rho$. The trivial zeros of $\zeta$ are at $s=-2,-4,\dots$; the nontrivial zeros lie in the
strip $0<\re s<1$. The Riemann Hypothesis asserts that every nontrivial zero has $\re s=\tfrac12$.
\end{theorem}

\begin{theorem}[Infinite Euler product; {\cite[Ch.~1]{Titchmarsh86}}]\label{thm:euler}
\lean{zetaEulerLog, zetaC2_summable, zetaC2_zero_free, zetaC2_euler}
For $\re s>1$,
\[
  \zeta(s)=\prod_{p\ \mathrm{prime}}\bigl(1-p^{-s}\bigr)^{-1}
         =\exp\!\Bigl(\sum_{p}\sum_{k\ge1}\tfrac1k\,p^{-ks}\Bigr),
\]
the product taken over the infinitely many primes; it converges absolutely and is zero-free on
$\re s>1$.
\end{theorem}

\begin{theorem}[Infinite Hadamard product; {\cite[Ch.~2]{Titchmarsh86}}]\label{thm:hadamard}
\lean{xi_entire, xi_zeros_eq_nontrivialZeros, zeta_hadamard,
riemannZeta_nontrivialZeros_infinite}
$\xi$ is entire of order $1$ with infinitely many zeros, which are exactly the nontrivial zeros
$\{\rho\}$ of $\zeta$, and admits the Hadamard factorization
\[
  \xi(s)=\xi(0)\prod_{\rho}\Bigl(1-\tfrac{s}{\rho}\Bigr)e^{s/\rho},
\]
the product taken over the infinitely many nontrivial zeros.
\end{theorem}

\begin{corollary}[Infinitude of the nontrivial zeros]\label{cor:hadamard-infinitude}
\lean{riemannZeta_nontrivialZeros_infinite}
\uses{thm:hadamard}
$\zeta$ has infinitely many nontrivial zeros.
\end{corollary}
\begin{proof}
Immediate from Theorem~\ref{thm:hadamard}: its zeros are infinitely many and are exactly
the nontrivial zeros of $\zeta$.
\end{proof}

\begin{theorem}[Hardy \cite{Hardy14}]\label{thm:hardy}
\lean{zeta_criticalLine_zeros_infinite}
Infinitely many nontrivial zeros of $\zeta$ lie on the line $\re s=\tfrac12$.
\end{theorem}

\subsection*{The compactified Riemann zeta on \texorpdfstring{$\Cstar$}{C*}}

For the slice construction, $\zeta$ is taken in compactified form, with values on the
Riemann sphere $\Cstar=\CC\cup\{\infty\}=S^2$ (the one-point compactification of $\CC$
\cite{Munkres00}). The classical facts above are retained as stated.

\begin{definition}[Compactified classical zeta]\label{def:zeta-Cstar}
\lean{zetaC, zetaC_coe, zetaC_one, zetaC_infty}
\uses{thm:riemann}
$\zetaC:\Cstar\to\Cstar$ is the meromorphic continuation of $\zeta$ regarded as a map to
the Riemann sphere: $\zetaC(s)=\zeta(s)$ for $s\in\CC\setminus\{1\}$; $\zetaC(1)=\infty$
(the simple pole, Theorem~\ref{thm:riemann}); and $\zetaC(\infty):=1$, an assigned
marked value (Remark~\ref{rmk:infty-marked}), selected by the one genuine limit
$\zeta(s)\to1$ as $\re(s)\to+\infty$ (the $p$-test: only the $n=1$ term of
$\sum_n n^{-s}$ survives). It is holomorphic into $\Cstar$ on $\CC\setminus\{1\}$,
its sole pole being $s=1$; at the domain point $\infty$ no analytic claim is made.
The functional equation, the conjugate symmetry
$\zetaC(\bar s)=\overline{\zetaC(s)}$, and Hardy's infinitude (Theorem~\ref{thm:hardy}) hold
verbatim for $\zetaC$, as they concern values on $\CC\setminus\{1\}$ where $\zetaC=\zeta$.
\end{definition}

\begin{remark}[The value at infinity is a marked point, not an extension]\label{rmk:infty-marked}
\uses{def:zeta-Cstar}
No analytic claim is made, or could be made, at the domain point $\infty$. A function
holomorphic on the whole Riemann sphere is rational, and $\zeta$ is not: it has an
essential singularity at $\infty$, tending to $1$ as $\re s\to+\infty$ while passing
through the trivial zeros along the negative real axis, so no assigned value makes
$\zetaC$ so much as continuous there. Definition~\ref{def:zeta-Cstar} therefore
performs two distinct acts, and only the first is analytic. Regarding the meromorphic
$\zeta$ as a holomorphic map into the target sphere, with $\zetaC(1)=\infty$, is the
classical sphere-valued reading of a meromorphic function. Assigning
$\zetaC(\infty):=1$ supplies the value the compactified setting consumes as data: the
point $\infty=N$ is shared by every slice Riemann sphere
(Definition~\ref{def:slices}), it is where the pole's value lands, and it is an
object of the projective base the transport machinery of Part~3 runs on
(Definition~\ref{def:base}), so compactified value paths are read at it. The
assignment is not free. Slice preservation forces the value to lie in every slice
sphere at once, and the intersection of the slice spheres is $\RR\cup\{N\}$; the
limit along the one region through which the construction approaches $\infty$ ---
the right half-plane, where $\zeta\to1$ genuinely --- then selects $1$. The
essential singularity is visible only along directions the construction never
travels. What is consumed downstream is the value, never smoothness: no statement
in this paper uses continuity or differentiability of $\zetaC$ at $\infty$. The
same discipline governs the whole class of Part~3: an A-section has infinitely many
zeros, hence is not rational, hence has an essential singularity at $N$; accordingly
every analytic hypothesis of Definition~\ref{def:A-section} is stated at finite
points, and the value at $N$ is carried as bare data (in the Lean development, the
field \texttt{valueAtInfinity}). This is the marked-extension device of
Remark~\ref{rmk:slice-pres-compact}: the compactification supplies geometry --- the
pole's value becomes a point of the target sphere, and all slice spheres share the
single point $N$ --- never an extension.
A companion distinction completes the picture. At $N$, Part~3 positions a
group action carrying analytic content: the slice-preserving content
the hypotheses supply over the octonions. There $N$ serves as an object
of the projective base, as the point of each value sphere at which the
distinguished M\"obius element arrives in Weierstrass form, and as the shared
point of the slice spheres. The analysis that
produces the element read there happens at the finite pole --- continuation
through the simple pole --- and its output is positioned at $N$ by the orbit
representative. No step of the construction evaluates the section analytically at the
domain point $\infty$: the position and the marked point are one geometric point playing
two roles, as a group element's position and as a function's domain point, and
only the first is ever active in the construction.
\end{remark}

\begin{lemma}[Compactification preserves the zero set]\label{lem:zero-Cstar}
\lean{zetaC_coe_eq_zero_iff, zetaC_zero_iff}
\uses{def:zeta-Cstar}
$Z(\zetaC)=Z(\zeta)$: the zeros of $\zetaC$ in $\Cstar$ are exactly the zeros of $\zeta$
in $\CC$; in particular the nontrivial zeros coincide.
\end{lemma}
\begin{proof}
On $\CC\setminus\{1\}$ the two functions agree, so they share all zeros there. The two
adjoined values lie outside the zero set: $\zetaC(1)=\infty$ is a pole and
$\zetaC(\infty)=1$. Hence compactification preserves the zero set exactly.
\end{proof}

%=====================================================================
\part{Slice-Preserving Theory, the Ring $\mathcal R$, and the Equivalence}
%=====================================================================

\section{Octonions and slices}

The octonions are an $8$-dimensional real division algebra, the largest of the normed division
algebras. This section fixes what we need from them: the algebra itself; Artin's theorem, that any
two octonions generate an associative subalgebra, so that powers and slicewise complex analysis
make sense; and the complex slices $\CC_v$, one for each imaginary direction $v\in S^6$, together
with their compactifications, the slice Riemann spheres. All the slices share the real axis and a
single point at infinity, the north pole $N$; that shared geometry is what the later construction
turns around.

The octonions $\OO$ have basis $\{1,e_1,\dots,e_7\}$ with each $e_i^2=-1$.

\begin{definition}\label{def:octonions}
\lean{Octonion, Octonion.normSq_mul, Octonion.alt_left, Octonion.alt_right}
The octonions are: (i) a \emph{division algebra}: every nonzero element has a
multiplicative inverse; (ii) a \emph{normed algebra}: $|ab|=|a||b|$; (iii)
\emph{non-associative}: $(ab)c\neq a(bc)$ in general; (iv) \emph{alternative}:
$(aa)b=a(ab)$ and $a(bb)=(ab)b$.
\end{definition}

An octonion $s\in\OO$ decomposes as $s=\sigma+w$ with $\sigma=\re(s)\in\RR$ and
$w=\im(s)\in\im(\OO)\cong\RR^7$; the unit imaginary octonions form $S^6\subset\im(\OO)$.

\begin{theorem}[Artin \cite{Schafer66}]\label{thm:artin}
In an alternative algebra, the subalgebra generated by any two elements is
associative.
\end{theorem}

\begin{corollary}\label{cor:powers}
\uses{thm:artin}
Powers $s^n$ are well-defined for any $s\in\OO$.
\end{corollary}

\begin{definition}[Complex slices and slice Riemann spheres]\label{def:slices}
\lean{Octonion.unitImaginarySphere, Octonion.sliceEmbed,
Octonion.sliceEmbed_mul, Octonion.sliceSphere}
\uses{def:octonions}
For any unit imaginary octonion $v\in S^6$, the \emph{complex slice} is
$\CC_v=\{a+bv:a,b\in\RR\}\subset\OO$. Since $v^2=-1$, this is a subalgebra isomorphic to
$\CC$ via the $\RR$-algebra isomorphism $\varphi_v:\CC\to\CC_v$, $i\mapsto v$. All slices
share the real axis: $\CC_v\cap\CC_w=\RR$ for $v\neq\pm w$. Its one-point compactification
is the \emph{slice Riemann sphere} $\CC_v^{*}:=\CC_v\cup\{\infty\}=S^2_v$, and $\varphi_v$
extends to an isomorphism of Riemann spheres $\varphi_v:\Cstar\xrightarrow{\sim}\CC_v^{*}$
with $\varphi_v(\infty)=\infty$ (GPV \cite{VS}, Prop.~4.1/Thm.~4.2). All slice spheres
share the single point at infinity, as well as the real axis.
\end{definition}

\section{Slice-preserving functions}
The class we use is the \emph{slice-preserving} slice-regular functions. We first recall slice
regularity \cite{GS07,GP11,CSS12,GSS13}, holomorphicity on each complex slice separately, as
the ambient notion, then isolate the slice-preserving subclass, on which classical complex methods
apply slicewise (by Artin's theorem each slice $\CC_v$ is an associative subalgebra where classical
complex analysis applies unchanged).

\begin{definition}[Slice regularity; {\cite[Def.~5.1.1]{CSS12}}]\label{def:slice-regular}
\uses{def:slices}
Let $\Omega\subseteq\OO$ be open. A function $f:\Omega\to\OO$ is \emph{slice regular}
if for each $v\in S^6$, $\varphi_v^{-1}\circ f\circ\varphi_v$ satisfies the
Cauchy--Riemann equations on $\varphi_v^{-1}(\Omega\cap\CC_v)\subseteq\CC$. A domain
$\Omega$ is \emph{axially symmetric} if $\sigma+\gamma v\in\Omega$ implies
$\sigma+\gamma w\in\Omega$ for all $w\in S^6$.
\end{definition}

\begin{theorem}[Representation Formula; {\cite[Thm.~5.1.7]{CSS12}}]\label{thm:rep-formula}
\uses{def:slice-regular}
Let $f$ be slice regular on an axially symmetric slice domain $\Omega$. Fix
$v_0\in S^6$ and write $f(\sigma+\gamma v_0)=\alpha(\sigma,\gamma)+v_0\cdot\beta(\sigma,\gamma)$
with $\alpha,\beta:\RR^2\to\RR$. Then for every $v\in S^6$,
$f(\sigma+\gamma v)=\alpha(\sigma,\gamma)+v\cdot\beta(\sigma,\gamma)$, where
$\alpha=\tfrac12[f(\sigma+\gamma v_0)+f(\sigma-\gamma v_0)]$ and
$\beta=\tfrac12 v_0^{-1}[f(\sigma-\gamma v_0)-f(\sigma+\gamma v_0)]$.
\end{theorem}

\begin{theorem}[Identity Theorem; {\cite[Cor.~5.1.9]{CSS12}}]\label{thm:identity}
\lean{stem_identity}
\uses{def:slice-regular}
Let $f$ be slice regular on an axially symmetric slice domain $\Omega$. If $f$
vanishes on a subset of a single slice $\Omega\cap\CC_{v_0}$ with an accumulation
point in $\Omega\cap\CC_{v_0}$, then $f\equiv0$ on all of $\Omega$.
\end{theorem}

\begin{theorem}[Extension Theorem; {\cite[Thm.~5.1.5]{CSS12}}]\label{thm:extension}
\uses{def:slice-regular}
Let $U\subseteq\CC$ be a domain symmetric under conjugation and $g:U\to\CC$
holomorphic. There is a unique slice regular $\mathrm{ext}(g):\Omega_U\to\OO$ on
$\Omega_U=\{\sigma+\gamma v:\sigma+i\gamma\in U,\ v\in S^6\}$ with
$\mathrm{ext}(g)|_{\CC_{v_0}}=\varphi_{v_0}\circ g\circ\varphi_{v_0}^{-1}$ for every
$v_0$.
\end{theorem}

\begin{definition}[Slice-preserving functions; \protect{\cite[Def.~2.7, Rem.~2.8]{AdF}; \cite{VS,SeriesExp,Wang}}]\label{def:slice-preserving}
\lean{IsIntrinsic, StemRing.isIntrinsic, StemRing.real_on_real}
\uses{def:slice-regular, thm:rep-formula}
A slice-regular function $f$ on an axially symmetric domain $\Omega$ is \emph{slice
preserving} if it maps every slice into itself: $f(\Omega\cap\CC_v)\subseteq\CC_v$ for all
$v\in S^6$. For octonions ($\#S^6>2$) the following are equivalent characterisations:
\begin{enumerate}[leftmargin=2.4em]
\item[(i)] $f$ is \emph{intrinsic}: $f(q^{c})=f(q)^{c}$ for all $q\in\Omega$, where $q^{c}$
  is octonionic conjugation (on a slice, $a+bv\mapsto a-bv$);
\item[(ii)] its representation-formula components are $\RR$-valued: writing
  $f(\sigma+\gamma v)=\alpha(\sigma,\gamma)+v\,\beta(\sigma,\gamma)$
  (Theorem~\ref{thm:rep-formula}), one has $\alpha,\beta:\RR^2\to\RR$, equivalently the
  inducing stem function $F=F_1+\iota F_2$ has $\RR$-valued components $F_1,F_2$;
\item[(iii)] $f$ preserves two distinct slices; equivalently $f(\Omega\cap\RR)\subseteq\RR$
  when $\Omega$ is a slice domain.
\end{enumerate}
Slice-preserving functions are the more restrictive subclass of slice-regular functions on
which classical complex methods apply slicewise. They are precisely the sections of
$\mathcal R$ (Definition~\ref{def:R}), and on them the regular product is the ordinary
pointwise product, so the product is commutative (Theorem~\ref{thm:wang}; Ghiloni--Perotti--
Stoppato \cite[Rem.~1.8]{SeriesExp}).
\end{definition}

\begin{remark}[Compactified extension]\label{rmk:slice-pres-compact}
Definition~\ref{def:slices} through Definition~\ref{def:slice-preserving}, and
Theorems~\ref{thm:rep-formula}--\ref{thm:extension}, are stated in the literature on
uncompactified domains $\Omega\subseteq\OO$ (and $U\subseteq\CC$). We use them on the
compactified $\Ostar=S^8$ and on the slice Riemann spheres $\CC_v^{*}=S^2_v$ (GPV's slice
Riemann manifolds \cite{VS}). The passage $\OO\rightsquigarrow\Ostar$,
$\CC_v\rightsquigarrow\CC_v^{*}$ adjoins only the shared point $\infty=N$; it is recorded as
a marked extension node (in the Lean development: cite the uncompactified statement and
extend through the one-point compactification \texttt{OnePoint}).
\end{remark}

\section{The octonionic zeta function \texorpdfstring{$\zetaO$}{zeta\_O}}

Slice-preserving theory already gives the tools to build the octonionic zeta function. On each slice
$\CC_v$ we read the compactified classical $\zeta$ through the isomorphism
$\varphi_v:\Cstar\to\CC_v^{*}$, and the representation formula glues these slice copies into a single
function on all of $\Ostar$. The one thing to check is that the two identifications of a slice with
$\CC$ agree, so that $\zetaO$ is well-defined; they do, because conjugating the input and reversing
the identification cancel. The pole of $\zeta$ at $s=1$ is carried to the north pole $N$.
At the shared point $N$ itself the function and the geometry are defined but nothing more
is claimed: the value there is the assigned marked value $1$ of
Remark~\ref{rmk:infty-marked} --- as real $s\to+\infty$, $\zeta(s)\to1$, the one genuine
limit --- and neither continuity nor differentiability at $N$ is asserted or used by the
argument.

\begin{definition}[Octonionic zeta]\label{def:zeta_O}
\lean{zetaO, zetaO_infty, zetaO_one, zetaO_sliceEmbed,
zetaO_sliceEmbed_of_im_neg, zetaO_ofReal}
\uses{def:zeta-Cstar, def:slices, thm:extension}
Here $N$ denotes the \emph{north pole} of $S^8$: the single point at infinity $\infty$
adjoined in the one-point compactification $\Ostar=\OO\cup\{\infty\}\cong S^8$ (inverse
stereographic projection; GPV \cite{VS}). It is one point with two names, so $\infty=N$;
it is the only point at infinity, shared by $\Ostar$ and by every slice Riemann sphere
$\CC_v^{*}=S^2_v$.

The \emph{octonionic zeta function} $\zetaO:\Ostar\to\Ostar$ is defined slicewise from the
compactified classical zeta (Definition~\ref{def:zeta-Cstar}) through the slice
identifications $\varphi_v:\Cstar\xrightarrow{\sim}\CC_v^{*}=S^2_v$ (the $\RR$-algebra
isomorphism $\CC\to\CC_v$ of Definition~\ref{def:slices}, extended to the Riemann spheres
by $\varphi_v(\infty)=N$; GPV \cite{VS}, Prop.~4.1/Thm.~4.2):
\begin{itemize}[leftmargin=2.2em]
\item[(i)] for $s\in\OO\setminus\RR$, with $w=\im(s)$, $v=w/|w|\in S^6$, $\gamma=|w|>0$,
  $\sigma=\re(s)$, so $s=\sigma+\gamma v$:\quad $\zetaO(s):=\varphi_v(\zetaC(\sigma+i\gamma))$;
\item[(ii)] for $s\in\RR\setminus\{1\}$:\quad $\zetaO(s):=\zetaC(s)=\zeta(s)\in\RR$;
\item[(iii)] for $s=1$ (the simple pole, Definition~\ref{def:zeta-Cstar}):\quad
  $\zetaO(1):=\infty=N$;
\item[(iv)] for $s=\infty=N$:\quad $\zetaO(\infty):=\zetaC(\infty)=1$, the assigned
  marked value of Remark~\ref{rmk:infty-marked}: as real $s\to+\infty$,
  $\zeta(s)\to1$, and the definition asserts the value only --- neither continuity
  nor differentiability of $\zetaO$ at $N$ is claimed here or used anywhere in the
  argument.
\end{itemize}
\end{definition}

\begin{proposition}[Well-definedness]\label{prop:well-defined}
\lean{zetaO_sliceEmbed', riemannZeta_conj}
\uses{def:zeta_O, thm:riemann}
$\zetaO:\Ostar\to\Ostar$ is well-defined.
\end{proposition}
\begin{proof}
For $s\in\OO\setminus\RR$ the decomposition $s=\sigma+\gamma v$ with $\gamma>0$,
$v\in S^6$ is forced ($\gamma=|\im s|$, $v=\im s/|\im s|$), so the only ambiguity is the
identification of the slice $\CC_v=\CC_{-v}$ with $\Cstar$. That slice carries two such
identifications, $\varphi_v$ ($i\mapsto v$) and $\varphi_{-v}$ ($i\mapsto-v$), under
which $s$ reads as $\sigma+i\gamma$ resp.\ $\sigma-i\gamma$; we must check
$\varphi_v(\zetaC(\sigma+i\gamma))=\varphi_{-v}(\zetaC(\sigma-i\gamma))$. By conjugate
symmetry $\zetaC(\sigma-i\gamma)=\overline{\zetaC(\sigma+i\gamma)}$
(\cite[Ch.~2]{Titchmarsh86}; Definition~\ref{def:zeta-Cstar}), writing
$\zetaC(\sigma+i\gamma)=a+bi$ gives $\varphi_{-v}(a-bi)=a+(-b)(-v)=a+bv=\varphi_v(a+bi)$:
the two sign reversals, one from conjugating the input, one from the reversed
identification, cancel. The remaining cases carry no ambiguity: for real $s\neq1$,
$\zetaO(s)=\zeta(s)\in\RR\subseteq\CC_v^{*}$ for every $v$; at $s=1$ the value is the
single point $\infty=N$; and at $s=\infty=N$ the value is $\zetaC(\infty)=1\in\RR$, again
the same for every $v$ since each $\varphi_v$ fixes $\RR$ pointwise.
Well-definedness at $N$ is a statement about the value, not about regularity: the
assigned marked value $1$ lies in $\RR$, which every slice sphere contains, so it is
unambiguous across all slices, and the function and the geometry are thereby defined
at $N$ while continuity and differentiability there are neither claimed nor used
(Remark~\ref{rmk:infty-marked}).
\end{proof}

\section{\texorpdfstring{$\mathcal R$}{R} is a commutative ring}\label{sec:R-ring}

The slice-preserving sections form a commutative ring $\mathcal R$, with $\zetaO$ among its elements.
Its multiplication is the regular ($\ast$) product, which descends from the Cayley--Dickson
construction of $\OO$ and from slice regularity; on slice-preserving functions it restricts on each
slice to the ordinary pointwise product (Wang), and is therefore commutative. The ring axioms, and
the integral-domain property, follow slicewise.

We work on $\Ostar$; write $\Omega_v:=\Ostar\cap\CC_v=\CC_v^{*}=S^2_v$ for the slice
Riemann sphere (Definition~\ref{def:slices}).

\begin{definition}[The ring of slice-preserving functions under the regular product]\label{def:R}
\lean{StemRing}
\uses{def:slice-preserving, def:slices}
$\mathcal R:=\mathcal O_{\Ostar}$ is the set of all slice-preserving functions
$f:\Ostar\to\Ostar$ (Definition~\ref{def:slice-preserving}),
$f(\Omega_v)\subseteq\CC_v^{*}$ for all $v\in S^6$, slice regular away from their
marked points --- the poles, where the value is $\infty$, and the shared point $N$
(Remark~\ref{rmk:infty-marked}) --- equipped with pointwise addition and the regular
($\ast$) product.
\end{definition}

Commutativity of $\mathcal R$ rests on a single slice-preserving fact from the literature: on each
slice the regular product is ordinary multiplication.

\begin{theorem}[Regular product on slice-preserving functions; {\cite[Rem.~2.11]{Wang}}]\label{thm:wang}
\uses{thm:artin, def:R}
For slice-preserving $f,g$ (with $f(\Omega_v)\subseteq\CC_v$), the regular product
restricts on each slice to the pointwise product, $(f\ast g)|_{\Omega_v}(w)=f(w)\cdot
g(w)\in\CC_v$, and $f\ast g=g\ast f$, $(f\ast g)\ast h=f\ast(g\ast h)$ (via Artin,
Theorem~\ref{thm:artin}).
\end{theorem}

\begin{proposition}\label{prop:R-comm-ring}
\lean{StemRing}
\uses{thm:wang, thm:identity, def:R}
$(\mathcal R,+,\ast)$ is a commutative ring with identity $1\in\mathcal R$.
\end{proposition}
\begin{proof}
Closure under $+$ is pointwise. For $\ast$: fixing a slice $\CC_K$, the product
reduces to pointwise multiplication on $\CC_K$ by Theorem~\ref{thm:wang}, which is
commutative because $\CC_K$ is a commutative associative subalgebra; associativity
and distributivity are inherited slicewise, and a slice-regular function is
determined by its restriction to any one slice (Theorem~\ref{thm:identity}). The
constant $1$ is the identity.
\end{proof}

\begin{proposition}[$\mathcal R$ is an integral domain]\label{prop:R-domain}
\lean{StemRing.eq_zero_or_eq_zero_of_mul_eq_zero}
\uses{thm:wang, thm:identity}
If $f,g\in\mathcal R$ and $f\ast g=0$, then $f=0$ or $g=0$.
\end{proposition}
\begin{proof}
$(f\ast g)|_{\Omega_K}(w)=f(w)\cdot g(w)$ in $\CC_K\cong\CC$
(Theorem~\ref{thm:wang}). The restrictions $f|_{\Omega_K}$, $g|_{\Omega_K}$ are
holomorphic on the connected slice; by the classical identity principle the pointwise
product vanishes identically iff one factor does. By Theorem~\ref{thm:identity} the
corresponding global function vanishes.
\end{proof}

\section{The slice-preserving analytic toolkit on the compactified octonions}\label{sec:toolkit}

We collect the slice-preserving analytic facts used in Part~3 as cited statements: the exponential,
its logarithm manifold, its degenerate fibre, the slice/stem square and $I$-independent lift, and the
winding lift. Each is established in the literature on the uncompactified $\OO$ and used here on
$\Ostar=S^8$ and the slice Riemann spheres $\CC_v^{*}=S^2_v$ (the passage is recorded in
Remark~\ref{rmk:slice-pres-compact}; the compactified slice-manifold structure itself is GPV's \cite{VS}).

\begin{theorem}[The slice exponential; \protect{\cite[\S3]{GPVwind}; \cite[Prop.~5.1]{VS}; \cite[Rem.~2.23]{AdF}}]\label{thm:slice-exp}
\lean{Octonion.exp, Octonion.exp_sliceEmbed, Octonion.norm_exp,
Octonion.exp_ne_zero}
\uses{def:slice-preserving, cor:powers}
The exponential $\exp:\OO\to\OO\setminus\{0\}$, $\exp(q)=\sum_{k\ge0}q^{k}/k!$ (powers well defined by Corollary~\ref{cor:powers}), is a
slice-regular, slice-preserving entire function (Definition~\ref{def:slice-preserving}),
given on each slice by
\[
\exp(x+Iy)=e^{x}(\cos y+I\sin y),\qquad x,y\in\RR,\ I\in S^6 .
\]
Consequently $|\exp q|=e^{\re q}$ depends only on the real part, and $\exp$ is nowhere zero.
\end{theorem}

\begin{theorem}[The logarithm manifold $E^{+}_{\OO}$; \protect{\cite[Prop.~5.1, Rem.~5.2, Def.~5.3, Prop.~5.4, Def.~5.5]{VS}; \cite[\S3]{GPVwind}}]\label{thm:log-manifold}
\lean{Octonion.Eexp, Octonion.logManifold, Octonion.Llog,
Octonion.Llog_Eexp, Octonion.Eexp_Llog, Octonion.exp_Llog}
\uses{thm:slice-exp}
Let $T:\OO\to\OO\times\im\OO$, $T(x+Iy)=(\sinh x\cos y+I\sinh x\sin y,\ Iy)$, and
$\OO^{+}=\{q:\re q>0\}$. The \emph{logarithm manifold} is $E^{+}_{\OO}:=T(\OO^{+})$, the
semi-helicoidal hypercomplex $8$-manifold; equivalently it is the image of the exponential
graph, parametrised by the \emph{$E^{+}_{\OO}$-exponential}
\[
E:\OO\longrightarrow E^{+}_{\OO}\subset\OO\times\im\OO,\qquad
E(x+Iy)=(\exp(x+Iy),\,Iy)=(e^{x}\cos y+Ie^{x}\sin y,\ Iy),
\]
an immersion and a diffeomorphism, with inverse the \emph{$E^{+}_{\OO}$-logarithm}
$L:E^{+}_{\OO}\to\OO$, $L(q,p)=\log|q|+p$ (here $p=\mathrm{Arg}(q)$). Writing $\pi$ for the
projection to the first factor, $\pi\circ E=\exp$ (Theorem~\ref{thm:slice-exp}). The manifold
$E^{+}_{\OO}$ is an adapted blow-up of $\OO$ along the real axis. Its projection
$\pi:E^{+}_{\OO}\to\OO\setminus\{0\}$ has the degenerate real fibres responsible for the
failure of covering and openness. The blow-up unwraps the multivalued logarithm into the
single-valued $L$, and a branch of the
hypercomplex logarithm is $\log_{\OO}=L\circ(\pi|_{\Omega})^{-1}$ on any path-connected
$\Omega\subset E^{+}_{\OO}$ on which $\pi|_{\Omega}$ is injective.
\end{theorem}

\begin{lemma}[The degenerate set of the exponential; its fibre over a real value]\label{lem:exp-degenerate}
\lean{Octonion.exp_fibre_neg_real, Octonion.exp_eq_neg_ofReal_iff,
Octonion.logManifold_fibre_neg_real, Octonion.degenerate_level_readout}
\uses{thm:slice-exp, thm:log-manifold}
With $\pi$ the first-factor projection, $\pi\circ E=\exp$
\textup{(\cite[Rem.~5.2(a)]{VS}, the commuting triangle)}. Its spherical
\emph{degenerate set} makes the projection an adapted blow-up with failures of covering and
openness \textup{(\cite[Rem.~5.2(b)]{VS})}. The direction $I(q)$ becomes singular at every
point of $\RR$ \textup{(\cite[Rem.~2.1]{GPVwind})}.

\emph{Proved from the slice form.} For $r>0$,
\[
\exp^{-1}(-r)=\bigl\{\,\log r+I(2k{+}1)\pi:\ I\in S^6,\ k\in\ZZ\,\bigr\}:
\]
by Theorem~\ref{thm:slice-exp}, $\exp(x+Iy)=e^{x}\cos y+Ie^{x}\sin y=-r$ forces
$e^{x}\sin y=0$, hence $y\in\pi\ZZ$; negativity of the value forces $y=(2k{+}1)\pi$ and
$x=\log r$; the direction $I$ is unconstrained. The fibre is thus indexed by the single real
\emph{level} $\log r=\log|{-}r|$ and, within it, by the odd winding indices $2k{+}1$ carried as
band data (the winding lift, Theorem~\ref{thm:winding-lift}; \cite[Cor.~5.21]{GPVwind}).
The fibre formula also appears as Preface motivation in \cite{VS}
(p.~972); the slice-form derivation above proves it.
\end{lemma}

\begin{definition}[Slice restriction and the $I$-independent lift; \protect{\cite[Rem.~2.11]{Wang}; \cite[\S3.3,~\S3.7]{BisiWinkelmann}; \cite[Def.~2.1, Prop.~2.2]{SeriesExp}}]\label{def:slice-squares}
\lean{IsIntrinsic, ASection.realize, ASection.realize_mem_sliceSphere,
ASection.realize_equivariant}
A slice-preserving section $A$ carries each slice Riemann sphere into itself,
\[
  A(\CC_v^{*})\subseteq\CC_v^{*}=S^2_v\qquad(v\in S^6),
\]
so its restriction to a slice is a holomorphic self-map of that slice sphere,
\[
  A_v:=\varphi_v^{-1}\circ A\circ\varphi_v:\ \Cstar\longrightarrow\Cstar,
\]
through the identification $\varphi_v:\Cstar\xrightarrow{\sim}\CC_v^{*}$ of
Definition~\ref{def:slices}. The \emph{$I$-independent lift} is the fact that this holomorphic stem is
the \emph{same} for every $v$: there is one holomorphic $A_\circ:\Cstar\to\Cstar$ with
\[
  A|_{\CC_v^{*}}=\varphi_v\circ A_\circ\circ\varphi_v^{-1}\qquad\text{for all }v\in S^6
\]
(Wang \cite[Rem.~2.11]{Wang}; Bisi--Winkelmann \cite[\S3.7]{BisiWinkelmann}). Equivalently, $A$ is
\emph{intrinsic}, $A(q^c)=A(q)^c$, and is induced ($A=\mathcal I(F)$) by a single stem function
$F=F_1+\iota F_2:D\to\OO_{\CC}$ on the symmetric $D\subset\CC$ whose components $F_1,F_2$ are
$\RR$-valued (the stem-function formalism, Ghiloni--Perotti \cite[Def.~2.1, Prop.~2.2]{SeriesExp}):
\[
  A(x+Iy)=F_1(x+iy)+I\,F_2(x+iy),\qquad F_1,F_2\ \RR\text{-valued}.
\]
It is this single $I$-independent slice stem that makes a slice-preserving section
$\Gtwo$-equivariant, hence blind to the $\Gtwo$-orbit direction
(Definition~\ref{def:section-map}).
\end{definition}

\begin{theorem}[The winding lift; \protect{\cite[Def.~3.4, Def.~4.1, Prop.~4.2, Def.~5.11]{GPVwind}}]\label{thm:winding-lift}
\lean{exists_log_continuation, winding_lift_unique,
Octonion.lift_iff_continuation}
\uses{thm:log-manifold}
Let $\gamma:[a,b]\to\OO\setminus\{0\}$ be a path \textup{(}on each slice $\CC_v$ this is the
classical complex continuation\textup{)}. A \emph{continuation of the hypercomplex logarithm along
$\gamma$} is a path $\tilde\gamma:[a,b]\to\OO$ with $\exp\circ\tilde\gamma=\gamma$, and a \emph{lift
of $\gamma$} to the logarithm manifold $E^{+}_{\OO}$ \textup{(}Theorem~\ref{thm:log-manifold}\textup{)}
is a path $\Gamma:[a,b]\to E^{+}_{\OO}$ with $\mathrm{pr}_1\circ\Gamma=\gamma$:
\[
\begin{tikzcd}[row sep=1.6em, column sep=2.6em]
& \OO\\
{[a,b]} \arrow[ur,"\tilde\gamma"] \arrow[r,"\gamma"'] & \OO\setminus\{0\} \arrow[u,"\exp"']
\end{tikzcd}
\qquad\qquad
\begin{tikzcd}[row sep=1.6em, column sep=2.6em]
& E^{+}_{\OO} \arrow[d,"\mathrm{pr}_1"]\\
{[a,b]} \arrow[ur,"\Gamma"] \arrow[r,"\gamma"'] & \OO\setminus\{0\}
\end{tikzcd}
\]
The two data are \emph{equivalent}: a lift $\Gamma$ exists if and only if a continuation
$\tilde\gamma$ exists, and they correspond by
\[
\tilde\gamma=L\circ\Gamma,\qquad \Gamma=(\exp\circ\tilde\gamma,\ \im\tilde\gamma),
\]
where $L$ is the $E^{+}_{\OO}$-logarithm of Theorem~\ref{thm:log-manifold}
\textup{(}\cite[Prop.~4.2]{GPVwind}\textup{)}. For a \emph{loop} $\gamma:S^1\to\OO\setminus\{0\}$ the
lift is taken on the universal cover $\exp:i\RR\to S^1$, as a continuous $\Gamma:i\RR\to E^{+}_{\OO}$
with $\mathrm{pr}_1\circ\Gamma=\gamma\circ\exp$ \textup{(}\cite[Def.~5.11]{GPVwind}\textup{)}:
\[
\begin{tikzcd}[row sep=1.6em, column sep=2.6em]
i\RR \arrow[r,"\Gamma"] \arrow[d,"\exp"'] & E^{+}_{\OO} \arrow[d,"\mathrm{pr}_1"]\\
S^1 \arrow[r,"\gamma"'] & \OO\setminus\{0\}
\end{tikzcd}
\]
\end{theorem}

\begin{proposition}[Tameness, existence, and the winding number; \protect{\cite[Def.~4.20, Cor.~5.21, Cor.~5.22]{GPVwind}}]\label{prop:winding-signature}
\lean{winding_lift_unique, winding_loop_defect}
\uses{thm:winding-lift}
The lift of Theorem~\ref{thm:winding-lift} is governed by the companion structure of $\gamma$.
\begin{enumerate}[label=\textup{(\roman*)},leftmargin=2.6em]
\item \emph{Uniqueness (tameness).} $\gamma$ is \emph{tame} when it admits a \emph{unique} companion
  imaginary-unit field; its lift is then unique \textup{(}\cite[Def.~4.20]{GPVwind}\textup{)}.
\item \emph{Existence and closure.} A lift of a loop $\gamma$ exists if and only if the signature
  satisfies $\sigma(\gamma|_{[\xi_l,\xi_{l+1}]})\in\{0,-1\}$ on each obstruction interval; and
  \emph{if it exists, the lift is itself a loop} \textup{(}\cite[Cor.~5.13]{GPVwind}; pinned
  against the arXiv text; the earlier citation to Cor.~5.22 is superseded\textup{)}.
\item \emph{Winding number.} For an untwisted tame loop with even circular signature $\sigma^c$, the
  winding number is $\omega=|\sigma^c|/2$ \textup{(}\cite[Cor.~5.21]{GPVwind}\textup{)}.
\end{enumerate}
\end{proposition}

\section{Slice preservation and symmetries}

With the ring in hand, we check that $\zetaO$ actually lives in it, and record the symmetry that will
organize its zeros. Slice preservation says $\zetaO$ keeps each slice inside itself; together with
slice regularity this makes it a section of $\mathcal R$. The symmetry is the action of the
automorphism group $\Gtwo=\Aut(\OO)$, which fixes the real axis and rotates the imaginary directions
$S^6$, and $\zetaO$ is equivariant under it. This forces every non-real zero to fill a whole
$6$-sphere, the geometry the next section reads off.

\begin{theorem}[Slice preservation]\label{thm:slice-pres}
\lean{zetaO_mem_sliceSphere}
\uses{def:zeta_O, def:slice-preserving}
$\zetaO$ is slice preserving (Definition~\ref{def:slice-preserving}):
$\zetaO(\CC_v^{*})\subseteq\CC_v^{*}$ for every $v\in S^6$.
\end{theorem}
\begin{proof}
Fix $v\in S^6$ and $s\in\CC_v^{*}$. If $s\in\CC_v\setminus\RR$ then $s=\sigma\pm\gamma v$
and $\zetaO(s)=\varphi_v(\zetaC(\sigma\pm i\gamma))\in\varphi_v(\Cstar)=\CC_v^{*}$
(Definition~\ref{def:zeta_O}(i)). If $s\in\RR\setminus\{1\}$ then
$\zetaO(s)=\zeta(s)\in\RR\subseteq\CC_v^{*}$; while $\zetaO(1)=\infty\in\CC_v^{*}$ and
$\zetaO(\infty)=1\in\CC_v^{*}$. In every case the value remains in $\CC_v^{*}$.
\end{proof}

\begin{theorem}[$\zetaO$ is a section of $\mathcal R$]\label{thm:zeta-in-R}
\lean{zetaSection}
\uses{thm:slice-pres, thm:extension, def:R, prop:well-defined}
$\zetaO\in\mathcal R$.
\end{theorem}
\begin{proof}
On each slice the restriction is the compactified classical zeta read through the
isomorphism $\varphi_v$, namely $\zetaO|_{\CC_v^{*}}=\varphi_v\circ\zetaC\circ\varphi_v^{-1}$,
holomorphic into $\CC_v^{*}$ away from the single pole.
Equivalently $\zetaO=\mathrm{ext}(\zeta)$ is the slice-regular extension of the
holomorphic, conjugation-symmetric $\zeta$ on $\CC\setminus\{1\}$
(Theorem~\ref{thm:extension}, \cite[Thm.~5.1.5]{CSS12}); hence $\zetaO$ is slice regular
on $\Ostar\setminus\{1,\infty\}$ (slice regularity is slicewise holomorphy,
Definition~\ref{def:slice-regular}). By Theorem~\ref{thm:slice-pres} it is slice
preserving, $\zetaO(\CC_v^{*})\subseteq\CC_v^{*}$. A slice-regular, slice-preserving
function $\Ostar\to\Ostar$ is by definition a section of $\mathcal R=\mathcal O_{\Ostar}$
(Definition~\ref{def:R}); the regular product on such sections restricts on each slice to
the ordinary product (Theorem~\ref{thm:wang}, Wang \cite[Rem.~2.11]{Wang}). Its zeros
\textup{(}Theorem~\ref{thm:zero-spheres}\textup{)} make $\zetaO$ a nonunit.
\end{proof}

\begin{definition}[The automorphism group $\Gtwo$]\label{def:G2}
\lean{G2, G2.smul_one, G2.smul_ofReal}
$\Gtwo=\Aut(\OO)$ is the $14$-dimensional compact, connected, simply connected Lie group of
$\RR$-algebra automorphisms of $\OO$. Each $g\in\Gtwo$ fixes $1$, hence the real axis pointwise, and
acts on $\im\OO\cong\RR^7$ as an element of $SO(7)$; thus $g\cdot(\sigma+\gamma v)=\sigma+\gamma\,g(v)$.
\end{definition}

\begin{theorem}[$\Gtwo$-equivariance]\label{thm:G2-equiv}
\lean{zetaO_equivariant}
\uses{def:G2, def:zeta_O}
$\zetaO(g\cdot s)=g\cdot\zetaO(s)$ for all $g\in\Gtwo$, $s\in\OO$.
\end{theorem}
\begin{proof}
For $s\in\RR$, $g\cdot s=s$ and $\zetaO(s)\in\RR$ is fixed, so both sides equal $\zetaO(s)$. For a
non-real $s=\sigma+\gamma v$ ($v\in S^6$, $\gamma>0$), $g$ carries the slice $\CC_v$ to the slice
$\CC_{g(v)}$ by the isometric relabelling $\varphi_{g(v)}=g\circ\varphi_v$
(Definition~\ref{def:G2}). Since $\zetaO$ is defined slicewise by the same $\zetaC$ read through
$\varphi_v$ (Definition~\ref{def:zeta_O}(i)),
\[
  \zetaO(g\cdot s)=\varphi_{g(v)}\bigl(\zetaC(\sigma+i\gamma)\bigr)
                 =g\bigl(\varphi_{v}(\zetaC(\sigma+i\gamma))\bigr)=g\cdot\zetaO(s).
\]
\end{proof}

\begin{theorem}[\protect{\cite{Baez02}}]\label{thm:G2-S6}
\lean{G2.exists_smul_eq_of_mem_unitImaginarySphere}
$\Gtwo$ acts transitively on $S^6$ with stabilizer $\mathrm{SU}(3)$:
$\mathrm{SU}(3)\hookrightarrow\Gtwo\twoheadrightarrow S^6$.
\end{theorem}

\begin{remark}[The $\Gtwo$-action extends to $\Ostar=S^8$]\label{rmk:G2-compact}
\uses{def:G2, def:carrier}
The action above is stated on $\OO$, but $\mathcal H_1$ (Part~3) runs on the compactified
$\Ostar=S^8$. Each $g\in\Gtwo$ fixes the real axis pointwise (an automorphism fixes $1$, hence all
of $\RR$) and acts on $\im\OO=\RR^7$ as an element of $SO(7)$; in particular $g$ is a linear
isometry of $\RR^8=\OO$, hence proper, and therefore extends uniquely to a homeomorphism of the
one-point compactification $\Ostar=\OO\cup\{\infty\}=S^8$ fixing the point at infinity $\infty=N$ (the one-point
compactification is functorial on proper maps; in Lean, \texttt{OnePoint}). Thus
$\Gtwo\curvearrowright\Ostar$ fixes the compactified real axis $\RR\cup\{\infty\}$, a great circle
through $N$, pointwise, and acts on each direction-sphere $S^6$ exactly as in
Theorem~\ref{thm:G2-S6}. This is the action on which $\mathcal H_1$ is built, with $\mathcal R$ its
ring of invariants over $\Ostar$.
\end{remark}

\section{Zero geometry and the equivalence}

This section is the bridge the whole paper is built to cross. It shows that the zeros of $\zetaO$ are
exactly the classical nontrivial zeros, now spread out: each nontrivial zero $\rho$ of $\zeta$
becomes a whole $6$-sphere of zeros of $\zetaO$, its residue-$\CC$ zero-sphere, sitting at real part
$\re\rho$. The Riemann Hypothesis, that every $\rho$ has real part $\tfrac12$, becomes the statement
that all these $6$-spheres share the same real part, that they are concentric. From here on we work
with the spheres.

\begin{theorem}[Zero Equivalence Theorem]\label{thm:zero-equivalence}
\lean{sliceEmbed_eq_zero_iff, zetaO_zero_iff}
\uses{def:zeta_O, thm:slice-pres, lem:zero-Cstar}
Let $\rho=\sigma+i\gamma\in\Cstar$. For every $v\in S^6$,
\[
\zetaO(\sigma+\gamma v)=0\quad\Longleftrightarrow\quad\zetaC(\rho)=0 .
\]
Equivalently: a non-real $s\in\Ostar$ is a zero of $\zetaO$ if and only if its slice
coordinate $\varphi_v^{-1}(s)$ is a zero of $\zetaC$; by Lemma~\ref{lem:zero-Cstar} these
are exactly the nontrivial zeros of the classical $\zeta$.
\end{theorem}
\begin{proof}
By Definition~\ref{def:zeta_O}(i), $\zetaO(\sigma+\gamma v)=\varphi_v(\zetaC(\sigma+i\gamma))$.
The slice identification $\varphi_v:\Cstar\to\CC_v^{*}=S^2_v$ is an isomorphism of slice
Riemann spheres (the $\RR$-algebra isomorphism $\CC\to\CC_v$ of
Definition~\ref{def:slices}, extended by $\varphi_v(\infty)=N$; GPV \cite{VS}). An
isomorphism is injective and sends $0$ to $0$, so $\varphi_v(z)=0\iff z=0$; with
$z=\zetaC(\rho)$ this gives both implications. (Slice preservation,
Theorem~\ref{thm:slice-pres}, is what places the value in $\CC_v^{*}$ so that
$\varphi_v^{-1}$ applies on the non-real side.) The final clause is
Lemma~\ref{lem:zero-Cstar}.
\end{proof}

\begin{theorem}[Nontrivial zeros are infinitely many disjoint $6$-spheres]\label{thm:zero-spheres}
\lean{zeroSphere, zeroSphere_eq_orbit, zeroSphere_zeros,
zeroSphere_disjoint, zeroSpheres_infinite}
\uses{thm:zero-equivalence, thm:G2-S6, thm:G2-equiv, cor:hadamard-infinitude}
For a nontrivial zero $\rho=\sigma+i\gamma$ of $\zetaC$ (so $\gamma>0$) set
\[
S_\rho=\{\sigma+\gamma v:v\in S^6\}=\sigma+\gamma\,S^6\subset\OO .
\]
Then:
\begin{enumerate}[leftmargin=2.4em]
\item[(i)] $S_\rho$ \emph{is} a $6$-sphere, the round sphere of radius $\gamma$ centred
  at the real point $\sigma$, equivalently the $\Gtwo$-orbit of any one of its points
  $\sigma+\gamma v_0$ (Theorem~\ref{thm:G2-equiv}; $\Gtwo$ acts transitively on $S^6$ with
  stabiliser $\mathrm{SU}(3)$, Theorem~\ref{thm:G2-S6}, Baez \cite{Baez02}).
\item[(ii)] every point of $S_\rho$ is a zero of $\zetaO$, and conjugate zeros give the
  same sphere, $S_\rho=S_{\bar\rho}$;
\item[(iii)] the nontrivial-zero set is the disjoint union
  $Z^{\mathrm{nt}}_{\Ostar}=\bigsqcup_{[\rho]}S_\rho$ over conjugate pairs
  $[\rho]=\{\rho,\bar\rho\}$, with distinct pairs giving disjoint spheres;
\item[(iv)] there are infinitely many such spheres (Corollary~\ref{cor:hadamard-infinitude};
  also Hardy \cite{Hardy14}, Theorem~\ref{thm:hardy}).
\end{enumerate}
\end{theorem}
\begin{proof}
(i) The map $v\mapsto\sigma+\gamma v$ is the affine map $x\mapsto\sigma+\gamma x$ of
$\im\OO\cong\RR^7$ restricted to $S^6$; with $\gamma>0$ its image is the round $6$-sphere
of radius $\gamma$ centred at $\sigma$. By $\Gtwo$-equivariance
(Theorem~\ref{thm:G2-equiv}) and transitivity of $\Gtwo$ on $S^6$
(Theorem~\ref{thm:G2-S6}), $\Gtwo\cdot(\sigma+\gamma v_0)=\{\sigma+\gamma\,g(v_0):g\in\Gtwo\}=S_\rho$.
So $S_\rho$ \emph{is} $S^6$, scaled by $\gamma$, translated to $\sigma$, realized as
a single $\Gtwo$-orbit; no embedding or diffeomorphism is invoked.

(ii) For any $v\in S^6$, $\zetaO(\sigma+\gamma v)=\varphi_v(\zetaC(\rho))=\varphi_v(0)=0$
by the Zero Equivalence Theorem~\ref{thm:zero-equivalence}. And
$S_{\bar\rho}=\{\sigma+\gamma(-v):v\in S^6\}=S_\rho$, since $v\mapsto-v$ is a bijection of
$S^6$.

(iii) For ``$\subseteq$'': a non-real zero $s=\sigma+\gamma v$ of $\zetaO$ forces
$\zetaC(\sigma+i\gamma)=0$ by Theorem~\ref{thm:zero-equivalence}, so $s\in S_\rho$ with
$\rho=\sigma+i\gamma$ nontrivial (Lemma~\ref{lem:zero-Cstar}); ``$\supseteq$'' is (ii).
Disjointness: if $\sigma+\gamma v=\sigma'+\gamma'w$, uniqueness of the real/imaginary
parts gives $\sigma=\sigma'$ and $\gamma v=\gamma'w$, whence $\gamma=\gamma'$ (taking
norms, $|v|=|w|=1$) and $\{\rho',\bar\rho'\}=\{\rho,\bar\rho\}$; so two spheres coincide or
are disjoint.

(iv) There are infinitely many spheres because $\zeta$ has infinitely many nontrivial zeros
(Corollary~\ref{cor:hadamard-infinitude}; also Hardy, Theorem~\ref{thm:hardy}).
\end{proof}

% [Conjugate pairs, Disjointness, and Zero-Set Decomposition are folded into
%  Theorem~\ref{thm:zero-spheres} above (items (ii)--(iv) and its proof).]

\begin{theorem}[The equivalence: concentricity $\Leftrightarrow$ RH]\label{thm:rh-equiv}
\lean{riemannHypothesis_iff_concentric}
\uses{thm:zero-spheres, thm:riemann}
The following are equivalent:
\begin{enumerate}[leftmargin=2.4em]
\item[(a)] the Riemann Hypothesis;
\item[(b)] the nontrivial-zero $6$-spheres $\{S_\rho\}$ of $\zetaO$
  (Theorem~\ref{thm:zero-spheres}) are \emph{concentric}: they share a single common
  centre, necessarily a point of the real axis.
\end{enumerate}
When they hold, the common centre is $\tfrac12$.
\end{theorem}
\begin{proof}
Each $S_\rho$ is centred at the real point $\sigma_\rho$, where
$\rho=\sigma_\rho+i\gamma_\rho$ (Theorem~\ref{thm:zero-spheres}(i)).

$(a)\Rightarrow(b)$: under RH every $\sigma_\rho=\tfrac12$, so all spheres share the
centre $\tfrac12$.

$(b)\Rightarrow(a)$: suppose the spheres share a common centre $c\in\RR$. By the functional
equation (Theorem~\ref{thm:riemann}), if $\rho$ is a nontrivial zero then so is
$1-\bar\rho=(1-\sigma_\rho)+i\gamma_\rho$, whose sphere $S_{1-\bar\rho}$ is centred at
$1-\sigma_\rho$. Concentricity forces $\sigma_\rho=c=1-\sigma_\rho$, hence
$\sigma_\rho=\tfrac12$; as $\rho$ was arbitrary this is RH, and the common centre is $\tfrac12$.
\end{proof}

%=====================================================================
\part{The Categorical Construction and the Concentricity Theorem}\label{part:stack}
%=====================================================================

Part 2 revealed that $\zeta$ carries a special geometry: spread onto the octonions, its zeros are
concentric exactly when the Riemann Hypothesis holds. That raises the question this part answers.
What general properties of a section of the ring $\mathcal R$ of slice-preserving functions force
its residue-$\CC$ zeros to be concentric? Following that question leads to four conditions,
\textup{C1--C4}, and to the whole class of functions satisfying them, which we call the A-sections.
The main result, the Concentricity Theorem, is that the residue-$\CC$ zeros of any A-section are
concentric. The octonionic zeta is one A-section among many, and the Riemann Hypothesis is one of
its corollaries.

From the section's own data we read one distinguished exponential M\"obius action on the
Riemann sphere, presented by the Euler product on the half-space and by the Weierstrass product
after continuation through the finite pole $p_A$;
orbit--stabilizer proves once that the resulting object and arrow assignments define a functor
to $\SphereWorld$. Thereafter its concrete projective/M\"obius matrix actions, together with the
same action on normalized A-value states, assemble the functor
$\mathcal A_A:\mathcal B\to\mathbf{Grpd}$. Its Grothendieck
construction $\mathcal T_A=\int_{\mathcal B}\mathcal A_A$ gathers every state and every
value-transport into one total object. The zero-sphere system is the inverse-image subdiagram
$\mathcal R_A$, included by $\iota_A:\mathcal R_A\Rightarrow\mathcal A_A$; its action morphisms,
the concrete projective/M\"obius matrices in the base and the $G_2$ matrices in the fibres, make
$\int_{\mathcal B}\mathcal R_A$ transitive and hence connected.
The component colimit therefore collapses to a singleton, and its real-value readout is
one centre $c\in\RR$. The concentricity proof takes place in this residue-$\CC$ inverse-image
system inside the ambient total.

\section{The geometric worlds and the categorical construction}\label{sec:tower-stack}

This section assembles the parts the detector runs on. Part~2 established
that exponentials are slice-preserving and live in the ring $\mathcal R$;
this section develops a theory for working with these functions, one that
shows how they \emph{act}, from the point of view of action-groupoid theory.
The key categories, the ones that capture the analysis and the geometry of
these actions, are three: the octonionic world $\mathcal H_1$; the
projective base $\mathcal B$, by the slice-preserving literature the home of
the unique tame continuous lift of the degenerate GPV fibre; and a world of
Riemann spheres, a whole continuum of $\CC_I^{*}$, each with its own
M\"obius self-maps. The general exponential slice-preserving action, and the
A-section-specific action the proof uses, naturally carry value states of
$\OO^{*}$ to value states of $\OO^{*}$; the theory requires slice-preserving
normalization, so a functor on $\OO^{*}$, the A-section equivariant functor
below, assembles what such a total value system looks like: the total
Grothendieck construction. To get there, we first explain the theory
minimally, for one exponential action: orbit--stabilizer makes the main
slice-preserving functor well defined, and the one construction is
simultaneously a value state, a group element, an action, and a functor. The
resulting functor then supplies the concrete matrix actions used in the total
Grothendieck construction and in the transitivity proof for the A-section's
residue-$\CC$ subsystem.

The categorical inventory is small, and we lay it out once. Five standing
groupoids are introduced: the octonionic world $\mathcal H_1$, the sphere
world, the projective base $\mathcal B$, the total $\mathcal T_A$, and the
residue-$\CC$ total $\int_{\mathcal B}\mathcal R_A$, with the $G_2$ fibre
groupoids $\mathcal A_A(b)$ and $\mathcal R_A(b)$ standing over the base
inside the two totals. Three functors run between them: the equivariant
realization $A_{\OO}$, the slice functor $A^{\mathrm{slice}}$, and the
A-section functor $\mathcal A_A$. The slice functor depends on the disk action,
and orbit--stabilizer proves its assignments well defined, so the paper first
explains how that slice-preserving functor is built. There is a
general theory here, slice-preserving ring action groupoids over $\Ostar$,
one for every element of $\mathcal R$; this paper builds the A-section
instance, because the hypotheses select one element. The point of building
functors is the tools they unlock: the total Grothendieck construction is
what this proof uses, and other tools of action-groupoid theory and
categorical homotopy stand ready. Thomason's theorem \cite{Thomason79}
places the Grothendieck construction in its homotopy-colimit context. The
proof developed here runs through ordinary connected components: the hypotheses assemble the
A-section-specific slice-preserving total, its inverse residue-$\CC$ system is
transitive, transitivity makes it connected, connectedness makes $\pi_0$ a
singleton. Applying the Grothendieck-components theorem to the constructed
functor $\mathcal R_A$ identifies that singleton with
$\operatorname*{colim}_{\mathcal B}(\pi_0\circ\mathcal R_A)$. The intrinsic
real face of the bound A-specific action descends through the fibrewise
components to a natural transformation
$\pi_0\circ\mathcal R_A\Rightarrow\Delta(\RR)$, whose induced map from that
singleton colimit reads the one real number.

What would the graph of a slice-preserving section look like? For an ordinary function one
draws the set of pairs (input, value). Here input and value both live on the compactified
octonions, so the graph sits in sixteen real dimensions, and no one will ever see it. But
slice preservation organizes it completely: the value at a point lies on that point's own
slice sphere, and one holomorphic stem is repeated in every direction. The graph is one
object in the octonions, and each direction $I\in S^6$ shows it as an ordinary graph: the
slice Riemann sphere $\CC_I^{*}$ with the pairing (input coordinate, value coordinate) of
the restricted stem, the same picture in every direction, glued along the shared real
circle and the shared point $N$. The construction of this section is that description made
categorical, and it is the situation the Grothendieck construction was made for: since the
fibered categories of SGA1 it has held one total object as compatibly glued fibres over a
base. A state will record an input coordinate together with its positioned image and the
value realized there, a point of the graph made into an object, and the groupoids will
record the symmetries that carry graph points to graph points: projective transformations
of the shared real circle, M\"obius motions within each sphere, and $G_2$ rotations of the
direction. A reader who keeps ``one graph in the octonions, seen one slice sphere at a
time'' in mind has the right picture for everything below.

\subsection*{The compactified octonions and the point at infinity}

\begin{definition}[The compactified octonions $\OO^{*}=S^8$]\label{def:carrier}
\lean{Octonion}
Let $\OO^{*}=\OO\cup\{\infty\}=S^8$ be the Alexandroff one-point compactification of $\OO=\RR^8$;
the slice Riemann-sphere structure on $S^m$, $m=\dim_{\RR}K\in\{4,8\}$, is GPV \cite{VS},
Prop.~4.1 and Thm.~4.2. The adjoined point is the north pole $N$; this is the point at
infinity. In Lean the underlying compactification is
\texttt{Mathlib.\-Topology.\-Compactification.\-OnePoint}, $\OO^{*}=\texttt{OnePoint }\OO$ with
$\infty=\texttt{OnePoint.infty}$, and the homeomorphism $\OO^{*}\cong S^8$ is
\texttt{onePoint\-Equiv\-Sphere\-Of\-Finrank\-Eq} applied to $\dim_{\RR}\OO=8$
(\texttt{Mathlib.\-Topology.\-Compactification.\-OnePoint.\-Sphere}). The compactification is what makes
the value of a section at its real pole a well-defined point of $\OO^{*}$, the point at
infinity $\infty=N$. The groupoids $\mathcal H_1$ and $\SphereWorld$ are built on
these compactified values (Definition~\ref{def:two-worlds}); the distinct projective base
$\mathcal B$ and direction-indexed sphere world used in the construction below are introduced in
Definition~\ref{def:base}.
\end{definition}

There is one geometric north pole $N=\infty$, and it plays a role in each
groupoid. In $\mathcal H_1$ this point of $\OO^{*}$ is a $\Gtwo$-fixed
object: its own connected component, with $\Aut(N)=\Gtwo$. In
$\SphereWorld$ it is the common point at infinity
of the slice Riemann spheres, fixed by every direction arrow. The section
relates the two roles through \textup{C1}: the pole of $A$ has value
$\infty$, so the section carries the pole point to $N$ on every slice
sphere at once. The residue field of a closed zero sorts the picture:
residue-$\RR$ (the $\Gtwo$-fixed locus $\RR\cup\{N\}$, classically trivial) versus residue-$\CC$
(the $6$-sphere orbits, classically nontrivial; the classification is Part~2 and
Theorem~\ref{thm:zero-spheres}).
The later total construction uses this one shared compactified point $N$, the same for every
zero index.

Looking forward: everything the construction does at $N$ happens in action groupoids:
$N$ is an object, and the arrows are given by group elements, the distinguished
M\"obius action and $G_2$. $N$
serves as a base object, as the shared point of every slice sphere, and as the value of
the section at its pole; the section's own value at the domain point $N$ enters as
the marked datum of Remark~\ref{rmk:infty-marked}. The infinite Euler
family of \textup{C2} with the infinite divisor of
\textup{C3}--\textup{C4} makes every A-section essentially singular at the domain point
$N$; what Part~3 reads at $N$ is what the construction supplies there --- the geometry
of the point and the marked value.

\subsection*{What a slice-preserving section does}

\begin{definition}[The section's slice data and equivariance]\label{def:section-map}
\lean{ASection.realize, ASection.realize_mem_sliceSphere,
ASection.realize_equivariant, ASection.sliceCoord_smul_invariant}
\uses{def:slice-squares, def:slice-preserving, thm:rep-formula}
Let $A\in\mathcal R$. By the $I$-independent lift (Definition~\ref{def:slice-squares}; Wang
\cite[Rem.~2.11]{Wang}, Bisi--Winkelmann \cite[\S3.7]{BisiWinkelmann}) there is one holomorphic
stem with $\RR$-valued components,
\[
  A(x+Iy)=F_1(x+iy)+I\,F_2(x+iy),\qquad F_1,F_2\ \RR\text{-valued},
\]
the same for every $I\in S^6$. Three consequences used below:
\begin{enumerate}[label=\textup{(\roman*)},leftmargin=2.6em]
\item \emph{Slice preservation of values.} $A(\CC_I^{*})\subseteq\CC_I^{*}$ for every $I$: the
  value at a point lies on that point's own slice sphere (Definition~\ref{def:R}).
\item \emph{$\Gtwo$-equivariance.} For $g\in\Gtwo$ and non-real $q=x+Iy$,
  \[
    A(g\cdot q)=F_1+g(I)\,F_2=g\cdot\bigl(F_1+I\,F_2\bigr)=g\cdot A(q),
  \]
  since $g\cdot(x+Iy)=x+g(I)y$ (Definition~\ref{def:G2}) and $F_1,F_2$ are real; for
  $q\in\RR\cup\{N\}$ both sides are $\Gtwo$-fixed, because $A(\RR)\subseteq\RR$ and
  $A(N)\in\bigcap_I\CC_I^{*}=\RR\cup\{N\}$ by \textup{(i)}.
\item \emph{Blindness to the sphere direction.} On a $\Gtwo$-orbit $S_{(\sigma,\gamma)}$ the stem
  coordinates $(F_1,F_2)(\sigma+i\gamma)$ are constant; in particular if $A$ vanishes at one
  point of the orbit it vanishes on the whole orbit (Lemma~\ref{lem:residue-spheres}).
\end{enumerate}
A slice-preserving section carries one holomorphic stem per slice, and the direction $S^6$ is
traversed only by the $\Gtwo$-morphisms already present in both worlds.
\end{definition}

\subsection*{The octonionic slice-preserving action groupoids}

Three groups act on this picture, and they are the only three. $PGL(2,\RR)$ acts on the
compactified real axis all slices share --- projective transformations of the one great
circle. $G_2=\Aut(\OO)$ acts on the directions: it fixes the real axis pointwise and rotates
$S^6$ transitively. And each slice Riemann sphere carries the M\"obius transformations of its
own compactified plane. A group action is naturally a groupoid, and this is the only
categorical device the paper needs: the \emph{translation groupoid} of $G\curvearrowright X$
has the points of $X$ as objects and the elements $g\in G$ satisfying $g\cdot x=y$ as its
morphisms $x\to y$, so an arrow \emph{is} one group element $g$ with
$g\cdot x=y$. The name earns its keep for one reason:
$\pi_0(G\ltimes X)=X/G$ recovers the orbits (Quillen \cite[\S1]{Quillen73}; in Lean,
\texttt{CategoryTheory.ActionCategory}), and the readout of the main
theorem uses exactly this. Every carrier below is a groupoid of one
of these three actions.  The sphere world $\SphereWorld$ is the
direction-indexed continuum of those action groupoids: one Riemann sphere
per direction, each carrying its own M\"obius action, with the $G_2$ arrows
moving between sphere directions. This structure is defined explicitly in the next
subsection: $\SphereWorld$ holds a continuum of Riemann-sphere pictures side
by side, one for every direction --- the multiplied viewpoints of the epigraph,
and the form in which an eight-dimensional graph lets itself be seen. Concretely, the
projective action will be read in this direction-indexed sphere-world groupoid: its objects
are directions $I\in S^6$ equipped with their compactified slice spheres, its morphisms carry
one $G_2$ automorphism and one M\"obius transformation, and a normalized value state pairs such a direction with a
point of its slice.

\begin{definition}[The octonionic world and the slice decomposition]\label{def:two-worlds}
\lean{H1, SphereWorld}
\uses{def:carrier, def:G2, thm:G2-S6, rmk:G2-compact, def:slices}
The octonionic world is
\[
  \mathcal H_1=\Gtwo\ltimes\OO^{*}.
\]
Its components are the $\Gtwo$-orbits: the compactified real fixed locus $\RR\cup\{N\}$ and,
for each $\sigma\in\RR$ and $\gamma>0$, the residue-$\CC$ sphere
$S_{(\sigma,\gamma)}=\{\sigma+\gamma v:v\in S^6\}$, labelled by its real centre $\sigma$ and
radius $\gamma$ (Baez \cite{Baez02}; Theorem~\ref{thm:G2-S6}).

$\SphereWorld$ organizes the same compactified values by their slice
decomposition. Every slice contains the real axis, and
\[
  \OO^{*}=\bigcup_{I\in S^6}\CC_I^{*},\qquad
  \CC_I^{*}\cap\CC_J^{*}=\RR\cup\{N\}\quad(J\neq\pm I).
\]
Its objects are the directions $I\in S^6$ with their compactified slice
spheres --- one sphere per direction, each with the M\"obius action of
its own compactified plane, the $G_2$ arrows relabelling which sphere one
stands on; the points $0$ and $N$ lie on every slice at once. M\"obius
transformations thus enter the construction twice: complex ones as each
slice sphere's own action in $\SphereWorld$, and real ones as the
projective base below, $\mathrm{PGL}(2,\RR)$ acting on the one
compactified real line all slices share.

Slice preservation itself determines what must be tracked: the
projective great circle every slice shares --- the base $\mathcal B$
below; the continuum of Riemann spheres --- $\SphereWorld$; and the
octonionic world $\mathcal H_1$, where the image of a slice-preserving
function assembles them.
\end{definition}

\subsection*{The projective base}

The base is the projective action groupoid
\[
  \mathcal B
  =PGL(2,\RR)\ltimes\operatorname{OnePoint}(\RR).
\]
Its objects are points of the compactified real projective line, the real great circle of
$\OO^{*}$: the slice through the real axis that every slice sphere shares. Its arrows are
projective elements: an arrow $h:x\to y$ is the class of an
invertible real matrix acting on the compactified line,
\[
  h=\left[\begin{pmatrix}a&b\\[2pt]c&d\end{pmatrix}\right],\qquad
  h\cdot x=\frac{ax+b}{cx+d}=y ,
\]
one element of $\mathrm{PGL}(2,\RR)$ carrying $x$ to $y$. We call these projective
elements throughout.

\subsection*{The sphere world}

The sphere world is the companion category to the base
\textup{(}\texttt{SphereWorld}\textup{)}, and it is the one category of this
paper with a \emph{continuum} of groups. Its objects are a continuum of
Riemann spheres inside $\OO^{*}$: one compactified slice sphere
$\CC_I^{*}\subset\OO^{*}$ for each unit imaginary direction
$I\in S^6\subset\OO$, each carrying the M\"obius group of its own
compactified plane: the invertible complex matrices modulo the scalar
matrices,
\[
  \operatorname{M\ddot ob}(\CC^{*})
  =\mathrm{GL}(2,\CC)\big/\bigl\{\lambda\mathbf 1:\lambda\in\CC^{\times}\bigr\}
  =\mathrm{PGL}(2,\CC),
\]
with $\mathbf 1$ the identity matrix.
An arrow is a pair: one
automorphism in $G_2$, relabelling which sphere one stands on, and one
M\"obius transformation, acting on the sphere; a normalized value state pairs a
direction with a point of its slice. Beside the base the picture is
complete: the base is the real M\"obius world of the one line all slices
share, and the sphere world is the continuum of complex M\"obius worlds, one
for each unit imaginary direction.

\subsection*{The disk action: a general slice-preserving action from the
base to the sphere world}

The aim is to understand slice-preserving maps on $\OO^{*}$ in terms of
action groupoids, and the two groupoids are already on the table: the
projective base $\mathcal B$ just defined, whose arrows are classes of real
matrices, and the sphere world just introduced, whose spheres each carry
their own complex M\"obius group. This subsection builds the generic material of an
assignment from the first to the second: one complex M\"obius element,
produced by exponentiation and read onto every slice sphere by conjugation
--- a slice-preserving action. Everything here is general, with no A-section
hypothesis used; the hypotheses are stated in the next subsection, and they
specialize this general functor to the A-section's slice functor
$A^{\mathrm{slice}}$ below. The object being built is a functor from $\mathcal B$ to the sphere
world --- an assignment of objects to objects and arrows to arrows,
respecting sources and targets, that carries identities to identities
and composites to composites; these two equalities are its
\emph{functoriality} --- and the values of a section already live in the
sphere world: each value sits on its point's own slice sphere. Two groups are already chosen, and
they are the only two. $\mathrm{PGL}(2,\RR)$ is the classes of invertible
real matrices: the arrows of the base, acting on the compactified real line
that every slice shares. $\mathrm{GL}(2,\CC)$ is the invertible complex
matrices: its classes are the M\"obius transformations, one group on each
slice sphere, acting on that sphere's compactified plane. The real classes
are already the real part of each M\"obius group, and the disk
automorphisms, the Blaschke maps, are their Cayley conjugates. This is what
lets the main general slice-preserving exponential functor be built from a
general \emph{analytic} ring element equipped with its
degenerate GPV base. The construction uses two matrices, the Cayley matrix
and the exponential:
\[
 C=\begin{pmatrix}1&-i\\[2pt]1&i\end{pmatrix},\qquad
 M(z)=\begin{pmatrix}e^z&0\\[2pt]0&1\end{pmatrix}:
\]
$C$ is the map $w\mapsto(w-i)/(w+i)$, carrying the base's compactified real
line onto the unit circle of a slice sphere; $M$ takes a complex number $z$
to the diagonal matrix whose entry is $e^z$, exponentiation in matrix form.
An invertible matrix acts on the Riemann sphere by
\[
  \begin{pmatrix}a&b\\[2pt]c&d\end{pmatrix}\in\mathrm{GL}(2,\CC),
  \qquad
  \begin{pmatrix}a&b\\[2pt]c&d\end{pmatrix}\cdot w=\frac{aw+b}{cw+d}.
\]
The class of the matrix up to scalars acts, and this is the M\"obius
map itself: a scalar multiple, all four entries scaled at once, enters
numerator and denominator together and cancels,
\[
  \begin{pmatrix}\lambda a&\lambda b\\[2pt]\lambda c&\lambda d\end{pmatrix}\cdot w
    =\frac{\lambda a\,w+\lambda b}{\lambda c\,w+\lambda d}
    =\frac{\lambda\,(aw+b)}{\lambda\,(cw+d)}
    =\frac{aw+b}{cw+d},\qquad\lambda\in\CC^{\times}.
\]
The action is total on the compactified plane: numerator and denominator
never vanish at the same $w$, for if both did, then
$a(cw+d)-c(aw+b)=ad-bc$, the determinant, would vanish; where the
denominator alone vanishes the value is the point at infinity $N$.
And it is a group action: substituting one fraction into the other
multiplies the matrices,
\[
\begin{aligned}
  \begin{pmatrix}a&b\\[2pt]c&d\end{pmatrix}\cdot
  \left(\begin{pmatrix}a'&b'\\[2pt]c'&d'\end{pmatrix}\cdot w\right)
  &=\frac{a\,\frac{a'w+b'}{c'w+d'}+b}{c\,\frac{a'w+b'}{c'w+d'}+d}
   =\frac{a(a'w+b')+b(c'w+d')}{c(a'w+b')+d(c'w+d')}\\[2pt]
  &=\frac{(aa'+bc')\,w+(ab'+bd')}{(ca'+dc')\,w+(cb'+dd')}
   =\left(\begin{pmatrix}a&b\\[2pt]c&d\end{pmatrix}
   \begin{pmatrix}a'&b'\\[2pt]c'&d'\end{pmatrix}\right)\cdot w,
\end{aligned}
\]
the middle step clearing $c'w+d'$ from numerator and denominator. Acting
twice is acting by the matrix product, and every later appeal to the group
action is an appeal to this computation.

The functor being built must act on any slice sphere, at any
normalization, so the first fact to establish is that the M\"obius action
reaches every point of a sphere.

\begin{lemma}[The M\"obius action is transitive]\label{lem:mobius-transitive}
\uses{def:two-worlds}
Any two points of a slice Riemann sphere are joined by one M\"obius
transformation: the action of the classes of $\mathrm{GL}(2,\CC)$ on the
compactified plane is transitive.
\end{lemma}

\begin{proof}
The
translation to $t$ and the inversion are the classes
\[
  \begin{pmatrix}1&t\\[2pt]0&1\end{pmatrix}\cdot w=w+t,
  \qquad
  s=\left[\begin{pmatrix}0&1\\[2pt]1&0\end{pmatrix}\right],
  \quad s\cdot w=\frac{1}{w}.
\]
The translation by $t$ carries $0$ to $t$: one matrix for each finite
value. The inversion at $0$ has numerator $0\cdot 0+1=1$ and denominator
$1\cdot 0+0=0$: the denominator alone vanishes, and where that happens
the value is the point at infinity (the totality rule above). So
$s\cdot 0=N$. The
action at $N$ is then computed by the displayed group-action calculation:
\[
\begin{aligned}
  \begin{pmatrix}a&b\\[2pt]c&d\end{pmatrix}\cdot N
   &=\begin{pmatrix}a&b\\[2pt]c&d\end{pmatrix}\cdot(s\cdot 0)
   =\left(\begin{pmatrix}a&b\\[2pt]c&d\end{pmatrix}
     \begin{pmatrix}0&1\\[2pt]1&0\end{pmatrix}\right)\cdot 0\\
   &=\begin{pmatrix}b&a\\[2pt]d&c\end{pmatrix}\cdot 0
   =\frac{a}{c},
\end{aligned}
\]
the first equality by $N=s\cdot 0$, the second by the displayed group-action
calculation, the
third by the matrix product --- the point $N$ itself when $c=0$. So the action is transitive: any two
points of the sphere are joined by one class, a composite of the displayed
translations and inversion.
\end{proof}

The real projective action supplies the more specific comparison used
later.  Let $a,b\in\RR$ and let $u,v$ lie in the same open half-plane.
The real M\"obius transformation
\[
 z\longmapsto b+
 \left(
 \frac{\operatorname{Im}(1/(v-b))}
      {\operatorname{Im}(1/(u-a))}\,\frac1{z-a}
 +\operatorname{Re}(1/(v-b))
 -\frac{\operatorname{Im}(1/(v-b))}
       {\operatorname{Im}(1/(u-a))}\,
       \operatorname{Re}(1/(u-a))
 \right)^{-1}                                      \tag{PGL}\label{eq:real-projective-transitivity}
\]
is represented by an invertible real matrix.  Its positive real ratio is
well defined because the two imaginary parts have the same sign.  At
$z=a$ the inner expression is the point at infinity, so the displayed
transformation has value $b$.  At $z=u$ the inner expression is exactly
$1/(v-b)$, so its value is $v$.  Thus one concrete element of
$\mathrm{PGL}(2,\RR)$ simultaneously carries the base point $a$ to $b$
and the half-plane point $u$ to $v$.  Conjugation by
$\operatorname{cayleyProjective}$ gives the corresponding matrix action
on every slice Riemann sphere.  This calculation will be cited when the
residue total is proved transitive.

The diagonal matrix
acts by
\[
  M(z)\cdot w=e^{z}w=e^{x}\bigl(e^{iy}w\bigr),\qquad z=x+iy,
\]
the modulus scaled by $e^{x}$ and the argument advanced by $y$.

Now the orbit--stabilizer construction. Take a value $w$ somewhere on a
slice Riemann sphere; by the transitivity just shown it lies in the orbit
of the M\"obius action, and in particular $C$ is itself an invertible
class, so there is exactly one point $w_0$
with
\[
  w=C\cdot w_0,\qquad w_0=C^{-1}\cdot w,
\]
the value standing as the $C$-positioned image of $w_0$. The matrix $M(z)$
acts at $w_0$ by $M(z)\cdot w_0=e^{z}w_0$; to act on the
value $w$ where it stands, conjugate, and group multiplication computes the
conjugated element doing exactly that:
\[
  (C\,M(z)\,C^{-1})\cdot w
    =(C\,M(z)\,C^{-1})\cdot(C\cdot w_0)
    =C\cdot\bigl(M(z)\cdot w_0\bigr),
\]
the inner $C^{-1}C$ cancelling by associativity of the group action. This is orbit--stabilizer,
and it is the positioning pattern of the whole paper: to act at a positioned
location, conjugate the action by the element that does the positioning ---
it runs over the base below at every object and arrow. The operation just
performed gets the name it keeps.

\begin{definition}[The Cayley conjugation and the disk action]\label{def:cayley-disk}
\lean{GreatCircle.cayleyGL, GreatCircle.cayleyProjective,
GreatCircle.diagonalGL, GreatCircle.expUnit, GreatCircle.diskExpAction}
\uses{def:two-worlds}
The passage from a complex number to a slice M\"obius map is two steps,
exponentiate then conjugate:
\[
  z\;\longmapsto\;M(z)\;\longmapsto\;
  D(z)=\operatorname{cayleyProjective}\bigl(M(z)\bigr)
      =\bigl[C\,M(z)\,C^{-1}\bigr]
      \in\operatorname{M\ddot ob}(\CC^*),
\]
and $\operatorname{cayleyProjective}=\bigl[\,C\,(\cdot)\,C^{-1}\,\bigr]$
is this same conjugation on every invertible matrix of
$\mathrm{GL}(2,\CC)$.
\end{definition}

$\operatorname{cayleyProjective}$ takes any invertible matrix, because
the positioned real matrices of the base pass through the same map at
the orbit--stabilizer factorization below
(Lemma~\ref{lem:orbit-stab}).

\begin{lemma}[Homomorphism identities]\label{lem:disk-homomorphism}
\lean{GreatCircle.cayleyProjective_mul, GreatCircle.diagonalGL_mul,
GreatCircle.diagonalGLHom, GreatCircle.expUnit_add,
GreatCircle.diskExpAction_add}
\uses{def:cayley-disk}
$\operatorname{cayleyProjective}$ is a group homomorphism, and
\[
  M(z+w)=M(z)M(w),\qquad D(z+w)=D(z)D(w).
\]
\end{lemma}

\begin{proof}
The homomorphism is the same cancellation that ran the conjugation: for
any two invertible matrices $M_1,M_2$,
\[
  C\,(M_1M_2)\,C^{-1}
    =C\,M_1\,(C^{-1}C)\,M_2\,C^{-1}
    =(C\,M_1\,C^{-1})(C\,M_2\,C^{-1}).
\]
Matrix multiplication with $e^{z+w}=e^ze^w$ gives $M(z+w)=M(z)M(w)$,
and applying the homomorphism $\operatorname{cayleyProjective}$ to that
equation gives $D(z+w)=D(z)D(w)$.
\end{proof}

Every later appearance of $\operatorname{cayleyProjective}$, positioning
the distinguished element over the base and acting by a stabilizer
factor, is this same map applied to a different matrix. Thus
exponentiation is the homomorphism turning addition of logarithmic lifts into composition
of the actual slice M\"obius maps --- recorded once, because every GPV lift below inhabits it:
\[
\begin{tikzcd}[column sep=huge]
\Gamma(t) \arrow[r,"\exp"] \arrow[d,"D"'] &
  v(t)\in\CC^{\times}
  \arrow[d,"u\mapsto\operatorname{cayleyProjective}(\operatorname{diag}(u{,}1))"] \\
D(\Gamma(t)) \arrow[r,equal] &
  \operatorname{cayleyProjective}\big(\operatorname{diag}(v(t){,}1)\big)
\end{tikzcd}
\]
The top row is the logarithmic lift; the bottom row is the M\"obius map it generates; the
square commutes at every instant, and adding a winding $2\pi ik$ moves only the top row,
$e^{\Gamma+2\pi ik}=e^{\Gamma}$. The Euler product will arrive as one completed logarithmic
sum, and this square is how one sum acts as one group element at every instant. In the Lean
source these are the chain \texttt{diagonalGL}, \texttt{diagonalGLHom}, \texttt{expUnit}, and
\texttt{diskExpAction}, with their multiplication identities.

Slice preservation is itself a group action, and already one exponentiation carries its
GPV degenerate base. For $u\in\CC^{\times}$ the logarithms of $u$ form the exponential fibre
\[
  \exp^{-1}(u)=\{\lambda+2\pi ik:\ k\in\ZZ\},\qquad
  M(\lambda+2\pi ik)=M(\lambda),
\]
the branches carried by the logarithm manifold of Theorem~\ref{thm:log-manifold} (the
commuting triangle $\pi\circ E=\exp$): a lift lives in this fibre, the winding is the
integer $k$ separating its elements, and $M$ carries the whole fibre to the one matrix.

The Euler product enters the group through its logarithmic coordinate.  There
the infinite product is an infinite \emph{sum}, and local normal
convergence makes the completed sum a single complex number at every point where it is
stated. (\emph{Locally normally} means: every point has a neighbourhood on which the sum of
suprema $\sum_p\sup\lvert\ell_p\rvert$ converges. This is the sense in which the cited Euler
and Weierstrass products converge, Theorem~\ref{thm:euler} and
Proposition~\ref{prop:weierstrass}; it gives uniform convergence on compacts, hence
continuity of the completed sum, and it is carried in the development as the summability and
majorant hypotheses of the A-section interface.) Exponentiation acts once, on the completed sum.

The slice facts of Part~2 make this passage well defined, and they act in
order. First the commuting triangle of Theorem~\ref{thm:log-manifold}, the
GPV logarithm manifold: the $E^{+}_{\OO}$-exponential $E$ and the
first-factor projection $\pi$ close over the exponential,
\[
\begin{tikzcd}[column sep=large,row sep=large]
\OO \arrow[r,"E"] \arrow[dr,"\exp"'] &
E^{+}_{\OO} \arrow[d,"\pi"]\\
& \OO\setminus\{0\},
\end{tikzcd}
\qquad \pi\circ E=\exp.
\]
The target is $\OO\setminus\{0\}$ because the exponential is nowhere
zero (Theorem~\ref{thm:slice-exp}); the cited spaces are the
uncompactified ones of the literature, read on $\OO^{*}$ as everywhere
through Remark~\ref{rmk:slice-pres-compact}. The blow-up unwraps the
multivalued logarithm once and for all: the fibre of $\pi$ over a value
$v$ is the exponential fibre $\exp^{-1}(v)$, its elements the branches,
carried into the manifold by $E$. And the fibre forces one real number.
A branch is any $\lambda$ with $e^{\lambda}=v$; taking norms, the slice
exponential gives $\lVert v\rVert=e^{\operatorname{Re}\lambda}$
(Theorem~\ref{thm:slice-exp}), so
\[
  \operatorname{Re}\lambda=\log\lVert v\rVert
\]
for every branch at once: two branches differ by the purely imaginary
$2\pi ik$, and the real coordinate is a function of the value alone.
This is the unique GPV real fibre, taken up in detail when the
A-section functors are built below. Along a path $\delta$, a continuation of the
logarithm is a lift $\Gamma$ with $e^{\Gamma(t)}=v(t)$; it exists whenever the value stays
away from zero along the path, and tameness makes it unique once its initial value $\Gamma(0)$ is fixed in the exponential fibre
(Theorem~\ref{thm:winding-lift}, Proposition~\ref{prop:winding-signature}). The winding is
the integer: two lifts of the same value path differ by one constant $2\pi ik$, and
exponentiation is blind to it, $e^{\Gamma+2\pi ik}=e^{\Gamma}$, so every branch
generates the same matrix at every instant while the real part
$\operatorname{Re}\Gamma(t)=\log\lVert v(t)\rVert$ is the one real value fixed by the
degenerate fibre. This is the general well-definedness: for any slice-preserving section
and any pole-avoiding, zero-free path, the completed logarithmic data determines one group
element and one real value at each instant, independent of every choice made along the
way.

\subsection*{The A-section hypotheses}

Here the hypotheses enter --- their first working appearance, so we state them informally at
once; the formal statement is Definition~\ref{def:A-section}. An \emph{A-section} is a
section $A\in\mathcal R$ with: \textup{(C1)} a meromorphic continuation with exactly one pole
--- simple and real --- whose pole-cancelled factor continues through that finite point with a
finite nonzero value; \textup{(C2)} an infinite prime-indexed Euler
product on a right half-space, each term a power series in $p^{-q}$ for its own prime ---
one logarithmic base;
\textup{(C3)} an infinite slice-regular
Weierstrass factorization whose sphere factors enumerate the residue-$\CC$ zeros;
\textup{(C4)} infinitely many such zeros. The hypotheses manufacture one number: C1 continues the pole-cancelled factor through
the simple pole, C2 and C3 present that same factor on a punctured neighbourhood of the pole, and its value at
the pole is the multiplier $u_A\neq0$ --- the infinite structure completed and evaluated,
which $M$ turns into one matrix.  No evaluation at the projective north object is part of
this analytic continuation.

The general machine now specializes to what the hypotheses supply. Shown at
the A-section itself, on the zero-free half-space of \textup{C2}:
\[
  \Gamma_A(z)=\sum_{p}\ell_p(z),\qquad
  \ell_p(z)=\sum_{k\ge1}\tfrac{a_{p,k}}{k}\,p^{-kz},\qquad
  e^{\Gamma_A(z)}=A(z),\qquad
  M\big(\Gamma_A(z)\big)=\operatorname{diag}\big(A(z),1\big):
\]
the infinite prime-power degenerate base completes to one sum, exponentiation
acts on it once --- $M(z+w)=M(z)M(w)$ is the homomorphism from addition to
composition --- and the infinite family enters the group as one matrix.
In this display and in every later analytic display, $z$ runs on the one
compactified plane every slice identifies with, and $A(z)$ is the
A-section's one holomorphic stem read there
(Definition~\ref{def:slice-squares}): one stem evaluation is the
section's evaluation on every slice at once. This
is also why the half-space is zero-free: on it $A=e^{\Gamma_A}$ is a value
of the exponential, and the exponential is nowhere zero
(Theorem~\ref{thm:slice-exp}). The
unique tame lift re-shows the same construction along any path: a
continuation $\Gamma$ with $e^{\Gamma(t)}=A(\delta(t))$ and fixed initial
value is unique (Theorem~\ref{thm:winding-lift}), so $M(\Gamma(t))$ is the
matrix of the value at every instant, the same for every branch of the
fibre.

\subsection*{The slice projection $A^{\mathrm{slice}}:\mathcal B\to\SphereWorld$}

Conditions \textup{C1--C3} construct one element, and this subsection builds
it and the slice functor it defines.

\begin{definition}[The multiplier and the distinguished disk action $D_A$]\label{def:DA}
\lean{ASection.distinguishedDiskAction,
ASection.distinguishedDiskAction_eq_fullMultiplier,
ASection.distinguishedPoleFactor_euler,
ASection.distinguishedPoleFactor_weierstrass, ASection.eulerPrimeSum}
\uses{def:A-section, def:cayley-disk, lem:disk-homomorphism}
Write $p_A$ for the pole of $A$ and $g_A=\texttt{distinguishedPoleFactor\,A}$ for the
pole-cancelled factor. On a punctured neighbourhood of the pole, conditions \textup{C2} and \textup{C3}
present this one factor twice:
\[
\begin{aligned}
 g_A(z)
   &=(z-p_A)
     \exp\!\left(\sum_{p}\ell_p(z)\right),\\[2pt]
 g_A(z)
   &=z^{m_A}R_{\mathrm{fac},A}(z)e^{h_A(z)}
     \prod_{n\geq0}\mathcal E\bigl(z;q_{A,n}\bigr).
\end{aligned}
\]
Condition \textup{C1} continues this same factor through the simple pole --- the finite
point $p_A$, never the domain point $\infty$ (Remark~\ref{rmk:infty-marked}); the sum
diverges at $p_A$ itself, the factor does not --- and evaluates it there:
\[
  u_A=g_A(p_A)\neq0.
\]
This is the one number the hypotheses manufacture: the infinite structure completed and
evaluated. Its logarithms form the exponential fibre of the previous subsection,
$\exp^{-1}(u_A)=\{\lambda_A+2\pi ik:\ k\in\ZZ\}$, and the function/group-element passage
is the commuting square
\[
\begin{tikzcd}[column sep=large,row sep=large]
z \arrow[r,"\exp"] \arrow[d,mapsto] & e^z\in\CC^\times \arrow[d,mapsto]\\
M(z)=\operatorname{diag}(e^z,1)
  \arrow[r,"\operatorname{cayleyProjective}"'] &
D(z)\in\operatorname{M\ddot ob}(\CC^*).
\end{tikzcd}
\]
through which $M$ collapses the fibre to the one matrix
\[
  M(\lambda_A)=\operatorname{diag}\bigl(e^{\lambda_A},1\bigr)=\operatorname{diag}(u_A,1),
\]
so the element the hypotheses construct is
\[
 D_A=D(\lambda_A)
   =\operatorname{cayleyProjective}\bigl(\operatorname{diag}(u_A,1)\bigr).
\]
Here $\lambda_A$ is the logarithm and $M$ is what undoes it: the entry of
$M(\lambda_A)$ is $e^{\lambda_A}=u_A$, the multiplier itself, so $D_A$ is an
exponential action and not a logarithmic one. The logarithm is the coordinate
in which the lift moves; the element it generates lives one exponential above
it. Every element of the fibre gives this same matrix and hence this same $D_A$: the
logarithm carries the winding, and the exponential action it generates is one element.
This is what lets the lift run through the whole base while $D_A$
stays one element. The Euler and Weierstrass expressions displayed above are the two
presentations of this same factor.
\end{definition}

The element is now fixed. To position it over the projective base, we first
construct the projective elements and the unique stabilizer factor that the
slice functor uses.

\begin{definition}[The orbit representatives]\label{def:orbit-reps}
\lean{GreatCircle.orbitRep, GreatCircle.orbitRep_spec}
\uses{def:carrier}
For each point $b$ of the base fix one projective element $o_b$ with
$o_b\cdot N=b$, once --- the development's \texttt{GreatCircle.orbitRep}
--- with $o_N=1$ at the north object.
\end{definition}

\begin{lemma}[North stabilizer uniqueness]\label{lem:orbit-stab}
\lean{GreatCircle.stabilizerPart, GreatCircle.orbit_stabilizer_factor,
GreatCircle.stabilizerPart_unique, GreatCircle.stabilizerPart_id,
GreatCircle.stabilizerPart_comp}
\uses{def:orbit-reps}
Every arrow $h:a\to b$ of $\mathcal B$ factors through the north
stabilizer $\operatorname{Stab}(N)=\{r:r\cdot N=N\}$ in exactly one way,
\[
  h=o_b\,r_h\,o_a^{-1},\qquad
  r_h=o_b^{-1}h\,o_a\in\operatorname{Stab}(N),
\]
and the stabilizer factor obeys the identity and composition equalities
\[
  r_{1_a}=1,\qquad r_{h_1\fatsemi h_2}=r_{h_2}\,r_{h_1}
\]
for every composable pair $h_1:a\to b$, $h_2:b\to c$.
Groupwise every symbol lives in the one group: $h$, $o_a$, $o_b$, and
$r_h$ are elements of $\mathrm{PGL}(2,\RR)$, classes of invertible real
matrices, and $r_h$ lies in the subgroup $\operatorname{Stab}(N)$; the
factorization is an equation in $\mathrm{PGL}(2,\RR)$.
\end{lemma}

\begin{proof}
Take an arbitrary arrow of the base, $h:a\to b$: one element of
$\mathrm{PGL}(2,\RR)$ with $h\cdot a=b$. Its source $a$ has the fixed
presentation $o_a\cdot N=a$, its target $b$ has $o_b\cdot N=b$. Form the
product $o_b^{-1}h\,o_a$ and let it act on $N$, one factor at a time, the
right factor first:
\[
  (o_b^{-1}h\,o_a)\cdot N
   = o_b^{-1}\cdot\bigl(h\cdot(o_a\cdot N)\bigr)
   = o_b^{-1}\cdot(h\cdot a)
   = o_b^{-1}\cdot b
   = N .
\]
The first equality is the group-action calculation; the second is the defining equation
of $o_a$; the third is $h\cdot a=b$; and the last inverts the defining
equation of $o_b$, from $o_b\cdot N=b$ to $o_b^{-1}\cdot b=N$. So
$r_h=o_b^{-1}h\,o_a$ lies in $\operatorname{Stab}(N)$, and
$h=o_b\,r_h\,o_a^{-1}$. The factor is unique: if $h=o_b\,r\,o_a^{-1}$
for any stabilizer element $r$, multiplying on the left by $o_b^{-1}$
and on the right by $o_a$ gives
\[
  r=o_b^{-1}h\,o_a=r_h .
\]
For the identity, the identity arrow $1_a:a\to a$ has source and target
$a$, so the formula $r_h=o_b^{-1}h\,o_a$ reads with $o_a$ on both sides,
and the representatives cancel:
\[
  r_{1_a}=o_a^{-1}\,1\,o_a=1.
\]
For composition take a composable pair --- the codomain of $h_1$
is the domain of $h_2$ --- $h_1:a\to b$ and $h_2:b\to c$. Their
composite $h_1\fatsemi h_2:a\to c$ is composition in diagrammatic order,
$h_1$ first, then $h_2$, and as matrices it is the product $h_2h_1$. Its
ends are $a$ and $c$, so the formula reads
$r_{h_1\fatsemi h_2}=o_c^{-1}(h_2h_1)\,o_a$, and inserting
$o_b\,o_b^{-1}=1$ between the two matrices splits it into the two
stabilizer factors:
\[
  r_{h_1\fatsemi h_2}=o_c^{-1}(h_2 h_1)\,o_a
    =(o_c^{-1}h_2\,o_b)(o_b^{-1}h_1\,o_a)=r_{h_2}\,r_{h_1}.
\]
\end{proof}

The chosen transporter from the reference object $N$ to each base point $b$
is therefore the fixed $o_b$, and what an arbitrary arrow adds is exactly
its uniquely determined stabilizer factor $r_h$.  This is a normalization
inside the displayed action matrices, not a projection functor and not an
analytic evaluation at $N$.

\begin{definition}[The slice functor $A^{\mathrm{slice}}$]\label{def:aslice}
\lean{ASection.AsectionSlice}
\uses{def:DA, def:orbit-reps, lem:orbit-stab}
The direction-and-M\"obius projection is
\[
  A^{\mathrm{slice}}:\mathcal B\longrightarrow\SphereWorld,
\]
written \texttt{AsectionSlice\,A} in the Lean source. It is a functor, and
its two assignments are built from three constructed pieces: one
projective element $o_b$ with $o_b\cdot N=b$, fixed for each object at
Definition~\ref{def:orbit-reps}, through which every arrow $h$
from $a$ to $b$ factors as $h=o_b\,r_h\,o_a^{-1}$
(Lemma~\ref{lem:orbit-stab}); the disk action $D_A$; and
$\operatorname{cayleyProjective}$, carrying the base's
matrices onto the slice spheres. On objects, $A^{\mathrm{slice}}(b)$ is
the normalized slice-sphere object at the already constructed position
\[
  \operatorname{cayleyProjective}(o_b)\,D_A.
\]
On arrows: $h$ goes to the sphere-world arrow whose $G_2$ automorphism is
the identity and whose M\"obius transformation is the factorization read
through the same builds,
\[
  \bigl(\operatorname{cayleyProjective}(o_b)\,D_A\bigr)\,
  \operatorname{cayleyProjective}(r_h)\,
  \bigl(\operatorname{cayleyProjective}(o_a)\,D_A\bigr)^{-1}.
\]
What makes the two assignments well defined
is $D_A$: it is the one factor either formula takes from the A-section,
and every branch of the exponential fibre gives that same element
(Definition~\ref{def:DA}). Functoriality is proved at
Proposition~\ref{prop:aslice-functor}, read through the disk action:
$D(z+w)=D(z)D(w)$ turns sums into compositions, and
$\operatorname{cayleyProjective}$ carries $r_{1_a}=1$ and
$r_{h_1\fatsemi h_2}=r_{h_2} r_{h_1}$ onto the M\"obius transformations ($\fatsemi$ is composition
in diagrammatic order, the left factor first); equal completed data gives
equal transformations.
\end{definition}

\begin{proposition}[$A^{\mathrm{slice}}$ is a functor]\label{prop:aslice-functor}
\lean{ASection.AsectionSlice, ASection.projectiveObjectFrame,
ASection.projectiveArrowElement, ASection.projectiveObjectFrame_north,
ASection.projectiveArrowElement_id, ASection.projectiveArrowElement_comp}
\uses{def:aslice, lem:orbit-stab}
The two assignments of Definition~\ref{def:aslice} --- the positioned
disk action at each object and the sphere-world arrow at each arrow ---
define a functor
$A^{\mathrm{slice}}:\mathcal B\to\SphereWorld$.
\end{proposition}

\begin{proof}
At an object $b$, the representative $o_b$ is one element of
$\mathrm{PGL}(2,\RR)$; $\operatorname{cayleyProjective}$ carries it onto
the slice spheres, and the action at $b$ is
\[
  \operatorname{cayleyProjective}(o_b)\,D_A .
\]
At the north position of the projective base, the chosen representative is
the identity, so the action placed in that fibre is $D_A$.  This is a
statement about projective positioning; it does not evaluate the A-section
at the domain point $N$. At an arrow $h:a\to b$, the
stabilizer factor $r_h$ of Lemma~\ref{lem:orbit-stab} gives the M\"obius
transformation
\[
  \bigl(\operatorname{cayleyProjective}(o_b)\,D_A\bigr)\,
  \operatorname{cayleyProjective}(r_h)\,
  \bigl(\operatorname{cayleyProjective}(o_a)\,D_A\bigr)^{-1}.
                              \tag{OS}\label{eq:orbit-stabilizer}
\]
Read from right to left, the inverse removes the source position,
$\operatorname{cayleyProjective}(r_h)$ applies the factor contributed by
$h$, and the left factor places the result at the target.

For identities, $r_{1_a}=1$, so the M\"obius transformation is
\[
  \bigl(\operatorname{cayleyProjective}(o_a)\,D_A\bigr)\,
  \operatorname{cayleyProjective}(1)\,
  \bigl(\operatorname{cayleyProjective}(o_a)\,D_A\bigr)^{-1}=1;
\]
its $G_2$ component is also the identity. For $h_1:a\to b$ and
$h_2:b\to c$, the identity
$r_{h_1\fatsemi h_2}=r_{h_2}r_{h_1}$ and the homomorphism
$\operatorname{cayleyProjective}$ give
\[
\begin{aligned}
  &\bigl(\operatorname{cayleyProjective}(o_c)\,D_A\bigr)\,
  \operatorname{cayleyProjective}(r_{h_2})\,
  \operatorname{cayleyProjective}(r_{h_1})\,
  \bigl(\operatorname{cayleyProjective}(o_a)\,D_A\bigr)^{-1}\\
  &\quad=\Bigl(\bigl(\operatorname{cayleyProjective}(o_c)\,D_A\bigr)\,
  \operatorname{cayleyProjective}(r_{h_2})\,
  \bigl(\operatorname{cayleyProjective}(o_b)\,D_A\bigr)^{-1}\Bigr)\\
  &\quad\phantom{=}\;\Bigl(\bigl(\operatorname{cayleyProjective}(o_b)\,D_A\bigr)\,
  \operatorname{cayleyProjective}(r_{h_1})\,
  \bigl(\operatorname{cayleyProjective}(o_a)\,D_A\bigr)^{-1}\Bigr).
\end{aligned}
\]
The left side is the transformation assigned to
$h_1\fatsemi h_2$, and the right side is the composite of the
transformations assigned to $h_1$ and $h_2$. Hence
$A^{\mathrm{slice}}$ preserves identities and composition.
\end{proof}

The orbit--stabilizer calculation has now completed its one role: it made
the object and arrow assignments of $A^{\mathrm{slice}}$ well defined and
proved that they preserve identities and composites
(Proposition~\ref{prop:aslice-functor}).  From this point onward the proof
uses the resulting functor itself.  Its internal factorization is not
reopened in the constructions of $\mathcal A_A$, $\mathcal T_A$, or
$\mathcal R_A$.

\begin{lemma}[Normalized pole-fibre transitivity]\label{lem:normalized-pole-transitive}
\uses{lem:mobius-transitive, prop:aslice-functor}
For any two normalized coordinates $u$ and $v$ in the same open
half-plane, there is a concrete invertible matrix morphism
$k:p_A\to p_A$ in $\mathcal B$ such that
\[
 A^{\mathrm{slice}}(p_A)^{-1}\mathbin{\cdot}
 \bigl(A^{\mathrm{slice}}(k)\mathbin{\cdot}
 (A^{\mathrm{slice}}(p_A)\mathbin{\cdot}u)\bigr)=v.
                                                        \tag{PT}\label{eq:normalized-pole-transitivity}
\]
\end{lemma}

\begin{proof}
The concrete real projective matrix in
\eqref{eq:real-projective-transitivity}, with the base source and target
both $p_A$, fixes $p_A$ and carries any prescribed point of the open
half-plane to any other.  Reading that calculation in the normalized
coordinate at $p_A$ conjugates its matrix action by the fixed invertible
action $A^{\mathrm{slice}}(p_A)$.  Conjugation preserves products,
inverses, and transitivity, and gives exactly the left side of
\eqref{eq:normalized-pole-transitivity}.  The resulting matrix is the
stated morphism $k:p_A\to p_A$.
\end{proof}

$A^{\mathrm{slice}}$ records the spheres and the positioned
M\"obius motion, and the equivariant realization $A_{\OO}$ below realizes
the section's value on them. The full A-section functor
below retains, in addition, the normalized point and the value produced on
that sphere.

Exponentiating the prime-power sum --- the A-section action whose
pole-cancelled continuation has the C3 Weierstrass presentation at the
finite point $p_A$ --- is why there is a fixed real-value transport: the path from
$\delta(0)$ to $\delta(1)$ determines the family of matrix actions
$D(\Gamma(t))$, together with its real level at every instant. The GPV
transport of the next subsection records exactly this.

\subsection*{The GPV transport}

The GPV lift records how $g_A$ varies between its presentations.

\begin{definition}[The GPV transport]\label{def:gpv-transport}
\lean{ASection.GpvTransport}
\uses{thm:winding-lift, prop:winding-signature, def:A-section}
Let
$a,b\in\mathcal B$. By the $I$-independent lift
(Definition~\ref{def:slice-squares}), evaluating the A-section off the
real axis is evaluating its one holomorphic stem on the one compactified
plane every slice identifies with; a projective-base transport from $a$
to $b$ records that stem evaluation along a path. It carries
three synchronized data, each a continuous function of one parameter
$t\in[0,1]$, the instant of the transport: a domain path $\delta^*$ ---
at each instant $t$, the point of that plane at which the stem is
evaluated, running from the point of $a$ to the point of $b$; the value
$v$ the section takes there; and a logarithmic lift $\Gamma$, the
continuous logarithm of that value.
The instant parametrizes the transport from $a$ to $b$ and belongs to the transport alone; the
geometry keeps its own coordinates, and the three data are read together at each instant.
The hypotheses are what make the instant necessary. C1 is a continuation, and continuation
happens along a path: the one factor is continued to the pole, and its value at
each instant is the continued value. The logarithm of C2's completed sum becomes
single-valued along a path: the unique tame lift fixes every branch from the initial one,
which is the GPV mechanism itself. And the argument reads what the path carries: the one
pole-cancelled factor from its Euler presentation on the half-space to its
Weierstrass presentation at the finite pole $p_A$, one real value carried
the whole way. What varies with the instant is the prime-power
sum: the degenerate base of C2 is evaluated along the domain path, and the matrix it
generates varies with it. The matrix acts throughout on one fixed input point
of its locus. The beginning and arrival of the path are determined by the hypotheses: it
begins in the Euler half-space of C2, where the prime-power degenerate base converges, and
ends at the pole of C1; the instants $t=0$ and
$t=1$ are this beginning and this arrival. Here the critical strip's classical mystery
takes its present form: the strip lies beyond the half-space where the prime-power base
converges, and the zeros stand there as fixed coordinates; the path is how the arithmetic
of the primes reaches them, as the matrix produced at the pole. The value at the pole
generates that matrix, and the matrix acts on the standing input: the prime-power sum
varies in the construction of the matrix, and the input it acts on is fixed.
The section takes its values on the compactified sphere, and the exponential
fibres live one exponential below, over $\CC^{\times}$; so the value is carried twice at once, as the
finite nonzero $v(t)\in\CC^{\times}$ that the exponential reaches, and as its reading
$v^{*}(t)$ on the compactified sphere, the same value seen where the section takes it. Their defining equations are
\[
\begin{aligned}
 \delta^*(0)&=\operatorname{complexPoint}(a),
 &\delta^*(1)&=\operatorname{complexPoint}(b),\\
 v^*(t)&=A\bigl(\delta^*(t)\bigr),
 &e^{\Gamma(t)}&=v(t),\\
 \Gamma(1)-\Gamma(0)&=2\pi i k,
 &&k\in\ZZ\quad\text{when the value path is a circuit}:
\end{aligned}
\]
the domain path runs from $a$ to $b$; the value, read on the sphere, is the section's
value along it; and the lift exponentiates to the value where the exponential lives.
The last equation is not asserted for a general path: it records the circuit case,
when the source and target values agree, and $k$ is its winding integer. These are precisely the fields of
\texttt{ASection.GpvTransport}.
\end{definition}

At each $t$ the lift equation
passes through the matrix square:
\[
\begin{tikzcd}[column sep=huge,row sep=large]
\Gamma(t) \arrow[r,"\exp"] \arrow[d,"D"'] &
v(t)\in\CC^\times
  \arrow[d,"{u\mapsto\operatorname{cayleyProjective}(\operatorname{diag}(u,1))}"]\\
D(\Gamma(t)) \arrow[r,equal] &
\operatorname{cayleyProjective}\bigl(\operatorname{diag}(v(t),1)\bigr).
\end{tikzcd}
\]
Thus every point of the lift is already an action of the same kind as $D_A$.

\begin{lemma}[The real level]\label{lem:real-level}
\lean{ASection.GpvTransport.level, ASection.GpvTransport.continuous_level,
ASection.GpvTransport.lift_unique, ASection.GpvTransport.level_independent,
ASection.GpvTransport.lift_endpoint_re_eq}
\uses{def:gpv-transport, thm:slice-exp, thm:winding-lift}
Along a GPV transport the real level
$\operatorname{Re}\Gamma(t)=\log\lVert v(t)\rVert$ is continuous, depends
only on the value along the path, and is the same for every lift of that
value path.
\end{lemma}

\begin{proof}
The lift equation is $e^{\Gamma(t)}=v(t)$. Take norms: the slice exponential
has $\lVert e^{\Gamma(t)}\rVert=e^{\operatorname{Re}\Gamma(t)}$
(Theorem~\ref{thm:slice-exp}), so
\[
  e^{\operatorname{Re}\Gamma(t)}=\lVert v(t)\rVert,
  \qquad
  \operatorname{Re}\Gamma(t)=\log\lVert v(t)\rVert:
\]
the real level is continuous and depends only on the value along the path.
Uniqueness and branch-independence are the cited GPV facts
(Theorem~\ref{thm:winding-lift}, Proposition~\ref{prop:winding-signature}):
the tame lift is unique once its initial value $\Gamma(0)$ is fixed, and any
two lifts of the same value path differ by one purely imaginary constant
$2\pi ik$, so every lift carries the same real part
\textup{(}\cite[Cor.~4.13]{GPVwind}\textup{)}.
\end{proof}

At each instant the
action is $D(\Gamma(t))$, moving with the value along the domain path: one
group element and one real value at every instant, the branch ambiguity
spent fibrewise on the degenerate fibre. The winding field records the
circuit case, the path whose value returns, and its integer counts the
branches crossed. This is the fixed real-value transport: the path from
$\delta(0)$ to $\delta(1)$ determines the matrix actions, and their real
levels define the transported real value.

Condition \textup{C2} now supplies the canonical arithmetic instance of this construction. Along
an Euler half-space loop $\delta$,
\[
  \Gamma_A(t)=\sum_{p}\ell_p(\delta(t)),
  \qquad e^{\Gamma_A(t)}=A(\delta(t)),
\]
with the prime index retained inside the complete summation. Local normal convergence makes
$\Gamma_A$ continuous, the zero-free half-space makes its value lie in $\CC^\times$, and the
GPV uniqueness theorem determines the whole lift from its initial value
$\Gamma_A(0)$ in the exponential fibre. The sum is
complete at every instant; only the evaluation point moves along the domain path. Since $\delta(0)=
\delta(1)$, the completed sum is evaluated at the same point at both ends, so the checked
winding is $k=0$: the loop returns to the same element of the fibre, $\Gamma_A(1)=\Gamma_A(0)$, level and
matrix with it.

The same construction continues to the finite pole.  This analytic passage
uses the finite continued multiplier and its exponential fibre; it does not
introduce a new projective object or a new value of the A-section.

\begin{lemma}[Finite-pole arrival]\label{lem:finite-pole-arrival}
\lean{ASection.GpvTransport.diskExpAction_eq_value,
ASection.GpvTransport.lift_endpoint_re_eq}
\uses{def:gpv-transport, def:DA, lem:real-level}
Run the domain path of a GPV transport of the continued factor to
$\delta(1)=p_A$. Then
\[
  v(1)=g_A(p_A)=u_A,\qquad
  D(\Gamma(1))=D(\lambda_A)=D_A,\qquad
  \operatorname{Re}\Gamma(1)=\operatorname{Re}\lambda_A,
\]
and two such runs differ by one whole winding,
$\Gamma_2(1)-\Gamma_1(1)=2\pi ik$ with $k\in\ZZ$.
\end{lemma}

\begin{proof}
The value along the run is $g_A(\delta(t))$: condition \textup{C1}
continues $g_A$ through the simple pole, so the value is nonzero the
whole way, entering from the Euler half-space where
$g_A=(z-p_A)\,e^{\Gamma_A}$ (condition \textup{C2}) and arriving at
$v(1)=g_A(p_A)=u_A$, where condition \textup{C3} presents the
multiplier $u_A$ in Weierstrass form. The arriving action is $D_A$ as
Definition~\ref{def:DA} constructed it,
$D_A=D(\lambda_A)=\operatorname{cayleyProjective}\bigl(\operatorname{diag}(u_A,1)\bigr)$,
and its level is $\operatorname{Re}\lambda_A$ by
Lemma~\ref{lem:real-level}. Every finite-pole run therefore produces this
same matrix action and the same real value.  The two lifts differ by one
purely imaginary constant $2\pi ik$, the whole winding; exponentiation
removes that difference, while taking real parts removes it as well.
\end{proof}

The prime sum, its value under $A$, its logarithmic level and winding, and the matrix action
it generates are the data of one element, $D(\Gamma(t))$ at each instant, travelling together.

Condition \textup{C4}, infinitely many residue-$\CC$ zeros, matches the class to its
infinite structure --- the infinite Euler family of \textup{C2}, the infinite Weierstrass
divisor of \textup{C3} --- and populates the degenerate fibre of the transport with an
infinite family; Theorem~\ref{thm:concentricity} reads one real value across the whole
population.

\subsection*{The equivariant realization $A_{\OO}:\mathcal H_1\to\mathcal H_1$}

An A-section is a ring element whose value at each point belongs to that
point's own slice sphere. The slice functor just built records the
positioned sphere-world action. The total Grothendieck construction also
uses its point-valued realization: input in $\OO^{*}$, output in
$\OO^{*}$, and the same $G_2$ element on each equivariant arrow. This
subsection defines that functor.

\begin{definition}[The equivariant realization]\label{def:equivariant}
\lean{ASection.AsectionEquivariant}
\uses{def:aslice, def:section-map, def:two-worlds}
The realization is a functor between two groupoids already built, from the
octonionic world to itself:
\[
  A_{\OO}:\mathcal H_1\longrightarrow\mathcal H_1,
  \qquad \mathcal H_1=\Gtwo\ltimes\OO^{*},
\]
written \texttt{AsectionEquivariant} in the Lean source. It is the
point-valued realization of the same slice-preserving action whose
positioned sphere-world projection is $A^{\mathrm{slice}}$. On an object
$q=x+Iy$ it uses the stem already constructed in
Definition~\ref{def:section-map}:
\[
  A_{\OO}(x+Iy)=F_1(x+iy)+I\,F_2(x+iy),
\]
with the compactified value prescribed at $N$. Thus no second section or
independent value map is introduced. On arrows, an arrow of $\mathcal H_1$
is one element $g\in G_2$ with $g\cdot q=q'$, and $A_{\OO}$ sends it to
that same element,
\[
  A_{\OO}\bigl(g:q\to q'\bigr)\;=\;
  g:A_{\OO}(q)\to A_{\OO}(q').
\]
\end{definition}

\begin{proposition}[$A_{\OO}$ is a functor]\label{prop:equivariant-functor}
\uses{def:equivariant}
The two assignments of Definition~\ref{def:equivariant} define a functor
$A_{\OO}:\mathcal H_1\to\mathcal H_1$.
\end{proposition}

\begin{proof}
The target is an arrow of $\mathcal H_1$ because equivariance
(Definition~\ref{def:section-map}(ii)) gives
$g\cdot A_{\OO}(q)=A_{\OO}(g\cdot q)=A_{\OO}(q')$. Preservation of identities
and composition is the group's own: the identity $1:q\to q$ goes to
$1:A_{\OO}(q)\to A_{\OO}(q)$,
and a composite $hg$ goes to $hg$, the same group elements on both sides.
\end{proof}

$A^{\mathrm{slice}}$ and $A_{\OO}$ are two functors of the same slice-preserving
action: the first records the projective direction-and-M\"obius transport, the
second the value of the section, retaining the same $G_2$ arrow by
equivariance. Condition \textup{C2} identifies the element acting in both with
the exponentiated log sum of prime powers --- the Euler product --- of the
given A-section; after the continuation of \textup{C1} through the finite
pole $p_A$, condition \textup{C3} presents the same pole-cancelled multiplier
in Weierstrass form and selects the residue-$\CC$ locus; and \textup{C4}
makes the selected population infinite. The A-section functor of the next
subsection uses both.

\subsection*{The unique GPV real fibre}

The projective group theory and the functoriality of
$A^{\mathrm{slice}}$ are already proved. The remaining analytic fact is the
branch-independence supplied by the GPV lift.

\begin{remark}[The unique GPV real fibre]\label{rmk:gpv-real-fibre}
\uses{thm:winding-lift, lem:exp-degenerate}
The real fibre is what is exponentiated. Over each base object, the
branches of the prime-power lift differ by $2\pi ik$. Consequently they
have the same real part, and $M$ sends every branch to the same matrix:
\[
  \operatorname{Re}(\lambda+2\pi ik)=\operatorname{Re}\lambda,
  \qquad M(\lambda+2\pi ik)=M(\lambda)
\]
(Definition~\ref{def:DA}). Local normal convergence makes the completed
prime-power sum one complex number at each point where it is defined; its
exponential is one value, and its logarithms form one exponential fibre.
At the pole $p_A$, condition C3 presents the completed multiplier $u_A$ in
Weierstrass form, and
\[
  \exp^{-1}(u_A)=\{\lambda_A+2\pi ik:\ k\in\ZZ\}
\]
has the single real level $\operatorname{Re}\lambda_A$.
\end{remark}

The analytic datum used below is the exponential fibre
over the continued multiplier: its branches differ by $2\pi ik$, every
branch has the same real part by GPV uniqueness, and every branch
exponentiates to the same A-specific disk action.

\subsection*{The functor for the total Grothendieck construction:
$\mathcal A_A:\mathcal B\to\mathbf{Grpd}$}

The total Grothendieck construction of the next subsection takes one
input: a functor from the base to $\mathbf{Grpd}$, and every ingredient of
this one is already built. An object of $\mathbf{Grpd}$ is a groupoid; $\mathcal H_1$, the sphere
world, and $\mathcal B$ are three of its objects. An arrow of
$\mathbf{Grpd}$ is a functor between groupoids, and the two functors just
built are two of its arrows: $A^{\mathrm{slice}}:\mathcal B\to\SphereWorld$
records where the inputs are positioned, and
$A_{\OO}:\mathcal H_1\to\mathcal H_1$ records the values
$A^{\mathrm{slice}}$ produces. The functor this subsection builds lands in
$\mathbf{Grpd}$ itself:
\[
  \mathcal A_A:\mathcal B\longrightarrow\mathbf{Grpd}.
\]
To each base object $b$ it assigns a groupoid: its objects are the
A-section's states at $b$, and its arrows are elements of the group $G_2$,
one element $\gamma$ with $\gamma\cdot x=y$ --- $G_2$ acting on those
states. To each projective element $h:a\to b$ it
assigns the transport of those states: each state moves by the concrete
M\"obius matrix action of $A^{\mathrm{slice}}(h)$
(\texttt{projectiveArrowHom}). The states are positioned by the A-section's action, and the hypotheses
built that action once, at the pole $p_A$: C1 continues $g_A$ through the
pole, the value $u_A$ there generates the disk action $D_A$, and C3 gives
it its Weierstrass form. The
total construction glues these groupoids, one per point of the real great
circle, into $\mathcal T_A$ of the next subsection. The source
$\mathcal B=\mathrm{PGL}(2,\RR)\ltimes\mathrm{OnePoint}(\RR)$ is the
projective action groupoid: an object is a point of the compactified real
projective line, and an arrow $h:a\to b$ is one element of
$\mathrm{PGL}(2,\RR)$ with $h\cdot a=b$. The base is this circle for the
A-section's own reason: the compactified real axis is the one line every
slice shares, the finite pole $p_A$ is one of its objects, and every base
arrow is an invertible projective matrix whose
$\operatorname{cayleyProjective}$ action is an invertible M\"obius
transformation on the corresponding slice spheres.  The A-specific disk
action these matrices transport is the exponential of the complete
prime-power GPV lift, determined by its initial value in the exponential
fibre.

Each point of $\OO^{*}$ off the real axis lies on exactly one slice sphere,
so it is two data: the sphere, named by its direction $I\in S^6$, and the
point on it, $q=x+Iy$. Its slice coordinate is the complex number $u=x+iy$,
the point read on the one compactified plane $\Cstar$ that every slice
identifies with --- normalized means read there, the same plane for every
direction. The pair $(I,u)$ is that point in components, $q=(I,u)$, paired
as the normalized value states introduced with the sphere world, and
$\gamma\cdot(I,u)=(\gamma I,u)$.

\begin{definition}[The A-section states and their fibre groupoids]\label{def:states}
\lean{ASection.AsectionState, ASection.AsectionStateInput,
ASection.AsectionStateOutput,
ASection.AsectionState_input_then_equivariant,
ASection.AsectionActionState}
\uses{prop:aslice-functor, def:section-map}
On a base object $b$, the groupoid $\mathcal A_A(b)$ consists of the
A-section states positioned at $b$ by $A^{\mathrm{slice}}$. A state is
one triple,
\[
  \Bigl(\,(I,u),\ \ A^{\mathrm{slice}}(b)\cdot(I,u),\ \
  A_{\OO}\bigl(A^{\mathrm{slice}}(b)\cdot(I,u)\bigr)\,\Bigr):
\]
a slice direction $I\in S^6$ with a normalized input coordinate $u$, the
coordinate positioned at $b$, and the value produced there.
Only the input is free. The positioned entry is determined by the position, and
the value entry is the value $A_{\OO}$ computes at the positioned entry: the
second and third entries are one point of the equivariant realization's
graph. The middle entry is a coordinate on its slice Riemann sphere and the
third entry lies in $\OO^{*}$.  The label $b=p_A$ means that the triple lies
in the fibre positioned at the finite pole; it does not say that its middle
entry is the point $p_A$, and therefore does not force its value entry to be
$N$.  In particular, a residue state in that fibre has a C-residue zero in
its middle entry and zero in its value entry.  The action used here is
constructed from the finite value $u_A=g_A(p_A)\ne0$ of the
pole-cancelled factor that condition \textup{C1} continues through the pole.

For two such triples $x$ and $y$, define
\[
 \operatorname{Hom}_{\mathcal A_A(b)}(x,y)
 :=\{\gamma\in G_2:\gamma\cdot x=y\}.
\]
Thus an arrow is a single element $\gamma\in G_2$ acting on the slice
direction and carrying the rest of the triple with it. The disk action,
conjugated by $\operatorname{cayleyProjective}$, is the slice-preserving
action on each Riemann sphere of the continuum in $\SphereWorld$, the same
complex map on every sphere: in components $D\cdot(I,u)=(I,\,D\cdot u)$,
while $\gamma$ moves the direction, $\gamma\cdot(I,u)=(\gamma I,\,u)$. Each
acts on its own component, so the two commute:
\[
  D\cdot(\gamma\cdot q)=D\cdot(\gamma I,\;u)=(\gamma I,\;D\cdot u)
    =\gamma\cdot(I,\;D\cdot u)=\gamma\cdot(D\cdot q):
\]
the positioned entry moves coherently, and the value entry moves by
equivariance (Definition~\ref{def:section-map}(ii)).  The identity arrow of
$x$ is $1\in G_2$; if $\gamma:x\to y$ and $\delta:y\to z$, their composite
is $\delta\gamma:x\to z$ because
\[
 (\delta\gamma)\cdot x=\delta\cdot(\gamma\cdot x)=z;
\]
and $\gamma^{-1}:y\to x$ is the inverse arrow.  Associativity, the identity
equalities, and the two inverse equalities are those of the group $G_2$.
Consequently $\mathcal A_A(b)$ is, concretely, the action groupoid of $G_2$
on the displayed triples; in particular every arrow is invertible. The fibres assemble the
equivariant realization's graph, and their value entries are where condition
\textup{C3} selects the residue-$\CC$ zero states the concentricity theorem
reads. In the Lean source the pairs are the states \texttt{AsectionState\,A},
whose point and value maps into $\mathcal H_1$ are
\texttt{AsectionStateInput\,A} and \texttt{AsectionStateOutput\,A}, natural
transformations checked with \texttt{AsectionState\_input\_then\_equivariant};
the triples are \texttt{AsectionActionState}, their $G_2$ action groupoid is
\texttt{AsectionActionStateFiber}, and
\texttt{AsectionActionFiber\,A\,X} is the fibre groupoid of the displayed
triples at the base object $X$.
\end{definition}

\begin{definition}[The transport]\label{def:transport}
\lean{ASection.AsectionActionTransport,
ASection.AsectionActionTransport_obj_input,
ASection.AsectionActionTransport_obj_positioned,
ASection.AsectionActionTransport_obj_value}
\notready
\uses{def:states, prop:aslice-functor}
On a base arrow $h:a\to b$, $\mathcal A_A(h)$ is a functor between the
two fibre groupoids,
\[
  \mathcal A_A(h):\mathcal A_A(a)\longrightarrow\mathcal A_A(b),
\]
called \emph{the transport}.  The arrow $h$ is already a concrete
invertible matrix morphism of the groupoid $\mathcal B$, and
$A^{\mathrm{slice}}(h)$ is already its concrete invertible matrix action
in $\SphereWorld$.  Definition~\ref{def:transport} uses that functorial
action directly; it does not reopen the construction that proved
$A^{\mathrm{slice}}$ well defined.

The middle entry of a state at $a$ is
$A^{\mathrm{slice}}(a)\cdot(I,u)$.  Apply the matrix action
$A^{\mathrm{slice}}(h)$ to that entry.  The normalized input of the
resulting state at $b$ is obtained by the inverse of the invertible action
at $b$.  Thus the whole assignment on a triple is
\[
\begin{aligned}
  &\Bigl(\,(I,u),\ \ A^{\mathrm{slice}}(a)\cdot(I,u),\ \
    A_{\OO}\bigl(A^{\mathrm{slice}}(a)\cdot(I,u)\bigr)\,\Bigr)\\
  &\quad\longmapsto
  \Bigl(\,\bigl(I,\
      A^{\mathrm{slice}}(b)^{-1}\mathbin{\cdot}
      \bigl(A^{\mathrm{slice}}(h)\mathbin{\cdot}
      (A^{\mathrm{slice}}(a)\mathbin{\cdot}u)\bigr)\bigr),\\
  &\hspace{3.8em}
    A^{\mathrm{slice}}(h)\mathbin{\cdot}
      \bigl(A^{\mathrm{slice}}(a)\mathbin{\cdot}(I,u)\bigr),\\
  &\hspace{3.8em}
    A_{\OO}\!\left(A^{\mathrm{slice}}(h)\mathbin{\cdot}
      \bigl(A^{\mathrm{slice}}(a)\mathbin{\cdot}(I,u)\bigr)\right)\,\Bigr).
\end{aligned}
\]
This is a state at $b$ because applying the action at $b$ to its displayed
normalized input gives its displayed middle entry:
\[
\begin{aligned}
 &A^{\mathrm{slice}}(b)\mathbin{\cdot}
 \left(I,
 A^{\mathrm{slice}}(b)^{-1}\mathbin{\cdot}
 \bigl(A^{\mathrm{slice}}(h)\mathbin{\cdot}
 (A^{\mathrm{slice}}(a)\mathbin{\cdot}u)\bigr)\right)\\
 &\qquad=
 A^{\mathrm{slice}}(h)\mathbin{\cdot}
 \bigl(A^{\mathrm{slice}}(a)\mathbin{\cdot}(I,u)\bigr).
\end{aligned}
\]
The third entry is then exactly what $A_{\OO}$ computes at that middle
entry.  On a fibre arrow $\gamma$, the transport uses the same matrix
$\gamma\in G_2$: the projective/M\"obius matrices change the complex
coordinate and the $G_2$ matrix changes the slice direction, so their
actions commute componentwise.  These are the object and arrow assignments
\textup{(}\texttt{AsectionActionTransport\_obj\_input},
\texttt{AsectionActionTransport\_obj\_positioned},
\texttt{AsectionActionTransport\_obj\_value}\textup{)}.
\end{definition}

\begin{proposition}[$\mathcal A_A$ is a functor]\label{prop:AA-functor}
\lean{ASection.AsectionActionDiagram,
ASection.AsectionActionTransport_id,
ASection.AsectionActionTransport_comp}
\notready
\uses{def:states, def:transport, prop:aslice-functor}
The assignments $b\mapsto\mathcal A_A(b)$ and $h\mapsto\mathcal A_A(h)$
define a functor $\mathcal A_A:\mathcal B\to\mathbf{Grpd}$, written
\texttt{AsectionActionDiagram\,A} in the Lean source.
\end{proposition}

\begin{proof}
Fix a base arrow $h:a\to b$.  A fibre arrow is a matrix
$\gamma\in G_2$ satisfying $\gamma\cdot(I,u)=(\gamma I,u)$, whereas every
matrix occurring in the object assignment of
Definition~\ref{def:transport} acts on the complex coordinate and leaves
$I$ fixed.  Hence applying $\gamma$ before that displayed assignment and
applying the same $\gamma$ afterward give the same complete triple.  The
middle entries agree by the matrix action on the complex coordinate, and
the third entries agree by the equivariance of $A_{\OO}$.  Therefore
$\mathcal A_A(h)$ sends the fibre arrow $\gamma$ to that same matrix
$\gamma$ between the transported states.  In particular it sends the
identity matrix to the identity matrix and matrix products to the same
products:
\[
 \mathcal A_A(h)(1_x)=1_{\mathcal A_A(h)(x)},\qquad
 \mathcal A_A(h)(\delta\gamma)
   =\mathcal A_A(h)(\delta)\,\mathcal A_A(h)(\gamma).
\]
Thus every $\mathcal A_A(h)$ is a functor between the two action groupoids.

For the identity base arrow $1_a$, Proposition~\ref{prop:aslice-functor}
gives $A^{\mathrm{slice}}(1_a)=1$.  The normalized-input entry of
Definition~\ref{def:transport} is therefore
\[
 A^{\mathrm{slice}}(a)^{-1}\mathbin{\cdot}
 \bigl(A^{\mathrm{slice}}(1_a)\mathbin{\cdot}
 (A^{\mathrm{slice}}(a)\mathbin{\cdot}u)\bigr)=u.
\]
The middle and value entries are likewise unchanged, and the arrow map
retains the same $G_2$ matrix.  Hence
\[
 \mathcal A_A(1_a)=1_{\mathcal A_A(a)}.
\]

For $h_1:a\to b$ followed by $h_2:b\to c$, apply the displayed object
assignment twice.  On the normalized input the two actions at the
intermediate object cancel, and functoriality of $A^{\mathrm{slice}}$
replaces the two consecutive matrix actions by the action of the base
composite:
\[
\begin{aligned}
 &A^{\mathrm{slice}}(c)^{-1}\mathbin{\cdot}
 A^{\mathrm{slice}}(h_2)\mathbin{\cdot}
 A^{\mathrm{slice}}(b)\mathbin{\cdot}
 \left(A^{\mathrm{slice}}(b)^{-1}\mathbin{\cdot}
 A^{\mathrm{slice}}(h_1)\mathbin{\cdot}
 A^{\mathrm{slice}}(a)\mathbin{\cdot}u\right)\\
 &\qquad=
 A^{\mathrm{slice}}(c)^{-1}\mathbin{\cdot}
 A^{\mathrm{slice}}(h_2)\mathbin{\cdot}
 A^{\mathrm{slice}}(h_1)\mathbin{\cdot}
 A^{\mathrm{slice}}(a)\mathbin{\cdot}u\\
 &\qquad=
 A^{\mathrm{slice}}(c)^{-1}\mathbin{\cdot}
 A^{\mathrm{slice}}(h_1\fatsemi h_2)\mathbin{\cdot}
 A^{\mathrm{slice}}(a)\mathbin{\cdot}u.
\end{aligned}
\]
The same equality gives the positioned entry; applying $A_{\OO}$ gives
the value entry; and on arrows both sides retain the same $G_2$ matrix.
Therefore
\[
 \mathcal A_A(h_1\fatsemi h_2)
   =\mathcal A_A(h_1)\fatsemi\mathcal A_A(h_2).
\]
These are exactly
\texttt{AsectionActionTransport\_id} and
\texttt{AsectionActionTransport\_comp}.  Hence the assignments define the
functor $\mathcal A_A:\mathcal B\to\mathbf{Grpd}$.
\end{proof}

$\mathcal A_A$ is the
graph of $A$ from the opening of the section, made categorical: input,
position, value, the three columns of the picture, with the transport squares
as the symmetries that carry graph points to graph points.

\subsection*{The total $\mathcal T_A=\int_{\mathcal B}\mathcal A_A$}

The Grothendieck construction gathers the indexed family built above ---
over each point $b$ of the projective base the fibre $\mathcal A_A(b)$ ---
into a single category, and the construction needs four things of its
input. First, the base is a groupoid. Second, the functor assigns a
groupoid to each base object. Third, it assigns a functor between those
groupoids to each base arrow. Fourth, the assignments are functorial,
identities to identities and composites to composites. The first is
already established: the base
$\mathcal B=PGL(2,\RR)\ltimes\operatorname{OnePoint}(\RR)$ is a groupoid
because its arrows are group elements --- an arrow is one projective
element $h$ with $h\cdot a=b$, composition is the
group's multiplication, and the group inverse inverts every arrow. The remaining three are met by what the previous subsection built. A
groupoid is assigned to each base object: the fibre $\mathcal A_A(b)$, the
$G_2$ groupoid of the state triples, each triple constructed from the two
functors $A^{\mathrm{slice}}$ and $A_{\OO}$
(Definition~\ref{def:states}). A functor is assigned to each base arrow: the
transport $\mathcal A_A(h):\mathcal A_A(a)\to\mathcal A_A(b)$
(Definition~\ref{def:transport}). And the assignments are functorial
(Proposition~\ref{prop:AA-functor}). With all four in hand,
$\mathcal T_A$ is a well-defined total groupoid, and this subsection
describes its objects and morphisms and proves the groupoid structure.

An object of the total $\int_{\mathcal B}\mathcal A_A$ is a pair
$(b,x)$: a point $b$ of the real great circle and a state $x$ of the
fibre $\mathcal A_A(b)$ (Definition~\ref{def:states}), the triple
\[
  x=\Bigl(\,(I,u),\ \ A^{\mathrm{slice}}(b)\cdot(I,u),\ \
  A_{\OO}\bigl(A^{\mathrm{slice}}(b)\cdot(I,u)\bigr)\,\Bigr)
\]
positioned at $b$: one graph point of the A-section. The pairing
carries analytic content, and it is how this construction was first
conceived: by GPV uniqueness, the state over $b$ has one exponential fibre ---
every branch of the lift over it differs by one purely imaginary constant
$2\pi ik$, so all its branches have the same real part, the level the
transport carries, and the exponential action every branch generates is one element
(Theorem~\ref{thm:winding-lift}, Lemma~\ref{lem:exp-degenerate}).

The morphisms of the total are fixed by the transport
(Definition~\ref{def:transport}): for two objects $(a,x)$ and
$(b,x')$,
\[
\begin{aligned}
  \operatorname{Hom}_{\mathcal T_A}\bigl((a,x),(b,x')\bigr)
  =\bigl\{(h,\gamma)\ :\ &h:a\to b\ \text{in }\mathcal B,\\
  &\gamma:\mathcal A_A(h)(x)\to x'\ \text{in }\mathcal A_A(b)\bigr\}.
\end{aligned}
\]
Here the base component $h$ is read through $\mathcal A_A(h)$, whose
transport is the A-specific Euler--Weierstrass action and GPV lift already
constructed for $h$.  The pair $(h,\gamma)$ therefore records the
concrete projective matrix action and the displayed $G_2$ comparison in the
target fibre; the GPV data are already part of the functor applied to $h$,
not an additional component of the pair.  The states $x$ and $x'$ are
triples in the fibres at $a$ and $b$, respectively.  The transport first produces
$\mathcal A_A(h)(x)$ in the fibre at $b$; the element $\gamma$ then compares
that triple with $x'$.  Thus a morphism has two stages, base then fibre:
$h$ carries the state to the fibre at
$b$, and $\gamma$ compares directions there, the first entry $(I,u)$
going to $(\gamma I,u)$ with the positioned and value entries moving by
the commutation and equivariance displayed at $\mathcal A_A$.

Both components are concrete matrix group actions in their corresponding
groupoid categories.  The projective matrix $h$ satisfies $h\cdot a=b$,
and the matrix $\gamma\in G_2$ satisfies
$\gamma\cdot\mathcal A_A(h)(x)=x'$.  Their identity matrices, products,
and inverse matrices are the identity morphisms, composites, and inverse
morphisms used below.

That these sets are the hom-sets of a groupoid --- identities,
composites, inverses --- is Proposition~\ref{prop:total-groupoid} below.

\begin{proposition}[$\mathcal T_A$ is a groupoid]\label{prop:total-groupoid}
\lean{ASection.ambientTotalCategory, ASection.ambientTotalGroupoid,
grothendieckGrpdGroupoid, CategoryTheory.Grothendieck}
\notready
\uses{prop:AA-functor}
$\mathcal T_A=\int_{\mathcal B}\mathcal A_A$ is a groupoid, with composite
\[
  (h_1,\gamma)\fatsemi(h_2,\gamma')
    =\bigl(h_1\fatsemi h_2,\ \gamma'\gamma\bigr),
\]
where $h_1\fatsemi h_2$ is multiplication in $\mathrm{PGL}(2,\RR)$,
$\gamma'\gamma$ is multiplication in $G_2$, and the A-specific slice
transports associated to $h_1$ and $h_2$ multiply in
$\operatorname{M\ddot ob}(\CC^*)=\mathrm{PGL}(2,\CC)$.  This M\"obius
transformation is the one already assigned by
$A^{\mathrm{slice}}(h_1\fatsemi h_2)$, and Proposition~\ref{prop:AA-functor}
identifies it with the composite of the two transports. Every morphism has the inverse
$(h,\gamma)^{-1}=(h^{-1},\gamma^{-1})$; the GPV transport already belonging
to $\mathcal A_A(h^{-1})$ is the reversal of the transport belonging to
$\mathcal A_A(h)$.
\end{proposition}

\begin{proof}
The identity morphism of $(a,x)$ is the pair $(1_a,1_x)$:
$\mathcal A_A(1_a)$ is the identity functor of the fibre
(Proposition~\ref{prop:AA-functor}), so $1_x$ is a fibre arrow from
$\mathcal A_A(1_a)(x)=x$ to $x$
\textup{(}\texttt{AsectionActionTransport\_id}\textup{)}.

For composition take $(h_1,\gamma):(a,x)\to(b,x')$ followed by
$(h_2,\gamma'):(b,x')\to(c,x'')$: so $\gamma$ is an arrow of the fibre
at $b$ with $\gamma\cdot\mathcal A_A(h_1)(x)=x'$, and $\gamma'$ one of
the fibre at $c$ with $\gamma'\cdot\mathcal A_A(h_2)(x')=x''$. The
composite must be a pair over $h_1\fatsemi h_2$: one arrow of the fibre
at $c$ from $\mathcal A_A(h_1\fatsemi h_2)(x)$ to $x''$. It is built in
three steps. First, by the composites-to-composites functoriality of
Proposition~\ref{prop:AA-functor},
\[
  \mathcal A_A(h_1\fatsemi h_2)(x)
  =\mathcal A_A(h_2)\bigl(\mathcal A_A(h_1)(x)\bigr).
\]
Second, apply the functor
$\mathcal A_A(h_2)$ to the arrow $\gamma$ of the fibre at $b$: by the
fibre-arrow computation in the proof of
Proposition~\ref{prop:AA-functor}, the result is the same group element,
now an arrow of the fibre at $c$,
\[
  \gamma\cdot\mathcal A_A(h_2)\bigl(\mathcal A_A(h_1)(x)\bigr)
  =\mathcal A_A(h_2)(x').
\]
Third, compose with $\gamma'$ in the fibre at $c$ --- multiplication in
$G_2$:
\[
  (\gamma'\gamma)\cdot\mathcal A_A(h_1\fatsemi h_2)(x)
  =\gamma'\cdot\mathcal A_A(h_2)(x')
  =x''.
\]
So $(h_1,\gamma)\fatsemi(h_2,\gamma')=(h_1\fatsemi h_2,\ \gamma'\gamma)$:
the base component is multiplied in $\mathrm{PGL}(2,\RR)$; its associated
slice transformation is multiplied in
$\operatorname{M\ddot ob}(\CC^*)=\mathrm{PGL}(2,\CC)$; and the fibre
component is multiplied in $G_2$.  The GPV transport belonging to
$\mathcal A_A(h_1\fatsemi h_2)$ is the composite of the transports already
belonging to $\mathcal A_A(h_1)$ and $\mathcal A_A(h_2)$. Associativity
follows from associativity in these groups together with the functoriality
that binds the slice M\"obius transport to the base composite
\textup{(}\texttt{AsectionActionTransport\_comp}\textup{)}.

For the inverse of $(h,\gamma)$: the base is a groupoid, so
$h\fatsemi h^{-1}=1_a$ and $h^{-1}\fatsemi h=1_b$.  Applying the identity
and composition equalities of Proposition~\ref{prop:AA-functor} gives
\[
\begin{aligned}
 \mathcal A_A(h)\fatsemi\mathcal A_A(h^{-1})
   &=\mathcal A_A(h\fatsemi h^{-1})
    =\mathcal A_A(1_a)=1_{\mathcal A_A(a)},\\
 \mathcal A_A(h^{-1})\fatsemi\mathcal A_A(h)
   &=\mathcal A_A(h^{-1}\fatsemi h)
    =\mathcal A_A(1_b)=1_{\mathcal A_A(b)}.
\end{aligned}
\]
Thus every transport $\mathcal A_A(h)$ is an isomorphism in
$\mathbf{Grpd}$ with the two-sided inverse functor
$\mathcal A_A(h^{-1})$, an
automorphism of the fibre $\mathcal A_A(a)$ when $b=a$ --- and
both functors carry each $G_2$ matrix to that same matrix.  Consequently
$(h,\gamma)^{-1}=(h^{-1},\gamma^{-1})$ composes with $(h,\gamma)$ to the
identity pairs on both sides. The total is a groupoid. In the Lean
source the construction is Mathlib's
\texttt{CategoryTheory.Grothendieck} applied to
\texttt{AsectionActionDiagram\,A}, and the groupoid classification is
\texttt{grothendieckGrpdGroupoid}.
\end{proof}

Every fibre consists of the complete A-section triples at its base object;
every transport is the concrete matrix action already assigned by
$A^{\mathrm{slice}}$; and every value entry is the equivariant realization
computed at the transported positioned point. The total is the A-section's graph made one category: what
moves from fibre to fibre is the A-specific Euler--Weierstrass/GPV action, and the residue-$\CC$
zero-sphere states are the degenerate-fibre outputs of this completed construction,
produced by \textup{C3--C4}. The total is read in two ways, and both are
functors it comes with. The projection to the base, $(b,x)\mapsto b$ and
$(h,\gamma)\mapsto h$, forgets the state
\textup{(}\texttt{CategoryTheory.Grothendieck.forget}\textup{)}; over it
the total holds one fibre per point of the real great circle. This
projection never moves a state into the pole fibre: passage between fibres,
including passage to or from $p_A$, is always an actual base morphism acted
on by the A-specific transport functor. And the
components diagram $\pi_0\circ\mathcal A_A:\mathcal B\to\mathbf{Set}$ sends
each base object to the set of components of its fibre; its colimit enters
the readout through Lemma~\ref{lem:pi0-grothendieck} --- $\pi_0$ of the
total is the colimit of the fibrewise components, which is how the next
section turns connectedness into one real number.

\begin{definition}[The projective base, the A-section functor, and the total category]\label{def:base}
\lean{GreatCircle.Base, GreatCircle.cayleyGL, GreatCircle.diagonalGL,
GreatCircle.diagonalGL_mul, GreatCircle.diagonalGLHom, GreatCircle.expUnit,
GreatCircle.expUnit_add, GreatCircle.diskExpAction, GreatCircle.diskExpAction_add,
GreatCircle.cayleyProjective, GreatCircle.cayleyProjective_mul,
ASection.distinguishedDiskAction, ASection.distinguishedDiskAction_eq_fullMultiplier,
ASection.distinguishedPoleFactor_euler, ASection.distinguishedPoleFactor_weierstrass,
ASection.eulerPrimeSum, ASection.eulerDiskAction_eq_value,
ASection.GpvTransport, ASection.GpvTransport.diskExpAction_eq_value,
ASection.GpvTransport.diskExpAction_endpoint_eq,
ASection.GpvTransport.lift_endpoint_re_eq, ASection.GpvTransport.level,
ASection.GpvTransport.continuous_level, ASection.GpvTransport.lift_unique,
ASection.GpvTransport.level_independent, ASection.GpvTransport.value_at_source,
ASection.GpvTransport.value_at_target,
ASection.GpvTransport.endpoint_log_norm_eq,
ASection.GpvTransport.endpoint_norm_eq,
ASection.canonicalAsectionPresentation_euler_prime_stack,
ASection.canonicalAsectionPresentation_euler_lift_exp,
ASection.canonicalAsectionPresentation_euler_lift_unique,
ASection.canonicalAsectionPresentation_euler_level_continuous,
ASection.canonicalAsectionPresentation_euler_winding,
ASection.distinguishedDiskAction_fixes_cayley_zero,
ASection.AsectionSlice, ASection.AsectionEquivariant,
ASection.AsectionStateInput, ASection.AsectionStateOutput,
ASection.AsectionState_input_then_equivariant,
ASection.ActionTransportSquare.coordinateTransport_commutes,
ASection.ActionTransportSquare.input_commutes,
ASection.ActionTransportSquare.output_commutes,
ASection.AsectionActionState, ASection.AsectionActionTransport,
ASection.AsectionActionState.ofInput,
ASection.AsectionActionState_ofInput_value,
ASection.AsectionActionTransport_obj_input,
ASection.AsectionActionTransport_obj_positioned,
ASection.AsectionActionTransport_obj_value,
ASection.AsectionActionTransport_id, ASection.AsectionActionTransport_comp,
ASection.AsectionActionDiagram, ASection.residueActionState,
ASection.residueActionState_mem, ASection.c4_infinite,
ASection.transportLevel,
ASection.ambientTotalCategory, ASection.ambientTotalGroupoid,
grothendieckGrpdGroupoid,
CategoryTheory.Grothendieck}
\notready
\uses{def:A-section, def:two-worlds, thm:winding-lift, prop:winding-signature}
The base is the projective action groupoid
$\mathcal B=PGL(2,\RR)\ltimes\operatorname{OnePoint}(\RR)$. The A-section functor is the
assignment constructed in this section --- on objects the $G_2$ action groupoid of
state triples positioned at each base object, on arrows transport by the M\"obius transformation
assigned by the already-proved functor $A^{\mathrm{slice}}$ and the
A-specific GPV transport --- written \texttt{AsectionActionDiagram\,A},
$\mathcal A_A:\mathcal B\to\mathbf{Grpd}$. The total value-transport category is
$\mathcal T_A=\int_{\mathcal B}\mathcal A_A$.
\end{definition}

\section{The concentricity theorem}\label{sec:detector}

We now have defined our groupoids and functors, so this section defines the class of functions
precisely and proves the theorem. Two facts about slice-preserving functions come first: the Weierstrass
factorization that gives condition C3 its meaning, and the fact that a residue-$\CC$ zero is a whole
$6$-sphere, a single connected component of the octonionic world. Then the categorical engine,
Riehl's computation of the components of a Grothendieck construction as a colimit, which sorts the
assembled total into the classes linked by the section's real-value transports. With these in hand
we define an A-section by its four conditions C1--C4, state and prove the Concentricity Theorem, and
record the corollary that carries its conclusion back into the language of classical zeros.

\begin{proposition}[Slice-regular Weierstrass factorization; the content of C3]\label{prop:weierstrass}
\lean{weierstrassE, spherePrimary, ASection.c3_multipliable,
ASection.c3_locMajorant, ASection.c3_factorization}
\uses{def:R}
Let $A\in\mathcal R$ have a zero of order $m\ge 0$ at the origin and no non-real isolated zeros
\textup{(}for $A\in\mathcal R$ the latter is automatic: non-real zeros come in whole spheres,
Lemma~\ref{lem:residue-spheres}\textup{)}. Then on $\OO^{*}\setminus\{N\}$ it factors as
\[
  A=q^{m}\,R\,S\,h,
\]
where $q$ denotes the octonionic coordinate function \textup{(}the identity section, so $q^m$
carries the order-$m$ zero at the origin\textup{)}, $R$ is a slice-preserving factor vanishing
exactly at the real zeros of $A$, $S$ is a factor vanishing exactly at its spherical
\textup{(}residue-$\CC$\textup{)} zero-spheres, and $h$ is never vanishing; each of $R,S$ is the
slice-regular Weierstrass product over the corresponding zeros. The factorization holds
\emph{including the case where either family of zeros is infinite}.
\end{proposition}
\begin{proof}
The proposition is general: $A$ here is any member of $\mathcal R$, the
hypotheses \textup{C1--C4} play no role in the argument, and either family of
zeros may be finite, infinite, or empty. Every analytic step happens one
algebra down, at the stem; the slice formalism then lifts the finished
factorization to $\OO^{*}$.

\emph{The reduction to the stem.} By the $I$-independent lift
(Definition~\ref{def:slice-squares}), $A$ is induced by one holomorphic stem
$F=F_1+\iota F_2$ with $F_1,F_2$ $\RR$-valued:
\[
  A(x+Iy)=F_1(x+iy)+I\,F_2(x+iy),\qquad
  \text{the same components for every }I.
\]
Real components give $F(\bar z)=\overline{F(z)}$, so the zero set of $F$ is
conjugation-symmetric: the real zeros $\rho_k$ stand individually, and the
non-real zeros stand in conjugate pairs $\{\rho_j,\bar\rho_j\}$. Under the lift
a real zero is a residue-$\RR$ zero of $A$, and one conjugate pair is one
zero-sphere, its whole $G_2$-orbit (Lemma~\ref{lem:residue-spheres}).

\emph{The classical factorization, symmetrized.} A holomorphic function $g$
has real Taylor coefficients exactly when it satisfies the conjugation
symmetry $\overline{g(\bar z)}=g(z)$; this is the form in which realness is
tracked below, and it is the symmetry the stem already has. The classical
Weierstrass theorem factors $F$ over its zero set: with the primary factors
\[
  E_p(w)=(1-w)\exp\Bigl(w+\tfrac{w^2}{2}+\dots+\tfrac{w^p}{p}\Bigr)
\]
and the integers $p_n$ chosen so that $\sum_n(r/|\rho_n|)^{p_n+1}$ converges
for every $r>0$, the product converges normally on compact sets and
\[
  F(z)=z^{m}\,\prod_{k}E_{p_k}\!\bigl(z/\rho_k\bigr)\,
       \prod_{j}\Bigl(E_{p_j}\!\bigl(z/\rho_j\bigr)\,
                     E_{p_j}\!\bigl(z/\bar\rho_j\bigr)\Bigr)\,h_0(z),
\]
with $h_0$ holomorphic and zero-free, the two factors of each conjugate pair
multiplied as displayed. Each
primary factor is built from real coefficients, so conjugating it conjugates
its argument: $E_p(\bar w)=\overline{E_p(w)}$. For a real zero $\rho_k$ the
factor $E_{p_k}(z/\rho_k)$ inherits the symmetry in $z$ directly. For a
conjugate pair, write $P_j(z)=E_{p_j}(z/\rho_j)\,E_{p_j}(z/\bar\rho_j)$ and
compute, the two factors exchanging under conjugation:
\[
  \overline{P_j(\bar z)}
    =\overline{E_{p_j}\!\bigl(\bar z/\rho_j\bigr)}\;
     \overline{E_{p_j}\!\bigl(\bar z/\bar\rho_j\bigr)}
    =E_{p_j}\!\bigl(z/\bar\rho_j\bigr)\,E_{p_j}\!\bigl(z/\rho_j\bigr)
    =P_j(z),
\]
real coefficients again, visible at the degree-two core:
\[
  \bigl(1-z/\rho\bigr)\bigl(1-z/\bar\rho\bigr)
    =1-2\operatorname{Re}(1/\rho)\,z+|\rho|^{-2}z^{2}.
\]
Finally, $F$ and every factor satisfy the symmetry, and the symmetry
passes to quotients, so the zero-free $h_0$ satisfies it as well. Every factor
of the factorization, $h_0$ included, is therefore a holomorphic stem with real
components.

\emph{The lift.} The stem formalism of Ghiloni--Perotti
\cite[Def.~2.1, Prop.~2.2]{SeriesExp} is stated for real alternative
$\ast$-algebras, hence for $\OO$, and for \emph{arbitrary} stems.  The real
components produced in the previous step give exactly two further properties: the induced section
is intrinsic, and on slice-preserving sections the regular product is the
pointwise product (Theorem~\ref{thm:wang}). So the factorization lifts factor
by factor to
$A=q^{m}RSh$ on $\OO^{*}\setminus\{N\}$: $R$ the lift of the real-zero product,
$S$ the lift of the paired products, each vanishing exactly on its zero-sphere,
and $h$ the lift of $h_0$, never vanishing. The compactified frame in which
this is read, the slice spheres of $\OO^{*}$ with $\varphi_v(\infty)=N$, is
GPV's (Remark~\ref{rmk:slice-pres-compact};
\cite[Prop.~4.1, Thm.~4.2]{VS}): slice preservation itself carries their
facts, before and independently of the transport sections. Convergence
transfers because the
components are direction-blind: for every $I$,
\[
  |A(x+Iy)|^{2}=F_1(x+iy)^{2}+F_2(x+iy)^{2}=|F(x+iy)|^{2},
\]
so a normal bound for the stem products is the same bound on every slice at
once.

\emph{Placement in the literature, and the instance used.} Over $\HH$ the
factorization is the theorem of Gentili--Vignozzi \cite{GV}, with the
slice-preserving case in Altavilla--de Fabritiis
\cite[Prop.~3.1, Thm.~3.2]{AdFslice}. Over $\OO$ we state and prove it by the
reduction above: for an intrinsic section the analysis is the classical
theorem at the stem, and the octonionic content is carried entirely by the
lift. The instance used in this paper is the A-section's: there the pair
$R,S$ is the residue split named in \textup{C3}, $R$ carrying the
residue-$\RR$ \textup{(}classically trivial\textup{)} zeros and $S$ the
residue-$\CC$ \textup{(}classically nontrivial\textup{)} zero-spheres, the
latter infinite in number by \textup{C4}.
\end{proof}

\begin{lemma}[Residue-$\CC$ zeros are $6$-sphere components of $\mathcal H_1$]\label{lem:residue-spheres}
\lean{ASection.sphereZero_mem_CResidueZeroLocus,
ASection.mem_CResidueZeroLocus_iff_exists_sphereZero,
G2.exists_smul_eq_of_mem_unitImaginarySphere}
\uses{def:two-worlds, thm:G2-S6, thm:wang}
For $A\in\mathcal R$, the non-real zeros of $A$ are $6$-spheres, each a single connected
component of $\mathcal H_1$.
\end{lemma}
\begin{proof}
$A$ is slice-preserving, so on a sphere $\sigma+\gamma S^6$ \textup{(}$\gamma>0$\textup{)} it has
the form $A(\sigma+\gamma v)=\alpha(\sigma,\gamma)+v\,\beta(\sigma,\gamma)$ with $\alpha,\beta$
real (the real stem, \cite{Wang}). Suppose $A$ vanishes at one point of the
sphere: $\alpha+v\beta=0$. The real part kills the first component: $v$ is
imaginary, so $\operatorname{Re}(\alpha+v\beta)=\alpha=0$. The modulus of what
remains kills the second: $|v\beta|=|v|\,|\beta|=|\beta|$, so $\beta=0$.
Neither component depends on $v$: the zero occurs at every point of
$\sigma+\gamma S^6$ or at none (the spherical-zero structure, \cite{AdF});
real zeros are isolated real points (residue field $\RR$).

The sphere $\sigma+\gamma S^6$ lies in $\OO^{*}$ off the real axis, and it is
one $\Gtwo$-orbit: $\Gtwo$ acts transitively on $S^6$
(Theorem~\ref{thm:G2-S6}, stabiliser $\mathrm{SU}(3)$), fixing $\sigma$ and
$\gamma$ and sweeping the direction. The connected components of
$\mathcal H_1=\Gtwo\ltimes\OO^{*}$ are its orbits, the computation
$\pi_0(G\ltimes X)=X/G$ (Quillen \cite[\S1]{Quillen73}), so each
zero-sphere is one component of $\mathcal H_1$. Therefore the residue-$\CC$
zeros, the non-real ones, are exactly the $6$-sphere components of
$A^{-1}(0)$.
\end{proof}

\begin{lemma}[$\pi_0$ of a Grothendieck construction]\label{lem:pi0-grothendieck}
\lean{CategoryTheory.Grothendieck, pi0GrothendieckEquiv, pi0_grothendieck}
\uses{def:base}
For a functor $F:\mathcal B\to\Grpd$, the connected-components functor carries the Grothendieck
construction to the colimit of the component diagram:
\[
  \pi_0\Bigl(\int_{\mathcal B}F\Bigr)\ \cong\ \operatorname*{colim}_{\mathcal B}\,(\pi_0\circ F).
\]
\end{lemma}
\begin{proof}
By Riehl's proof of Lemma~8.3.4, for any $X:\mathcal C\to\mathbf{Set}$ each arrow of the
category of elements imposes exactly the condition that the corresponding elements be identified
in every cone, so that $\pi_0(\operatorname{el}X)\cong\operatorname{colim}_{\mathcal C}X$
\cite[p.~102]{Riehl}. Taking $X=\pi_0\circ F$, the colimit performs exactly the identifications
generated by the maps of $F$. In the Lean development this equivalence is
\texttt{pi0GrothendieckEquiv}, built from the canonical cocone \texttt{pi0Cocone} and
\texttt{colimit.desc} over Mathlib's \texttt{colimit\_sound} and \texttt{colimit\_eq}.
\end{proof}

\begin{lemma}[Fibrewise $\pi_0$ descent and colimit readouts]\label{lem:connected-readout}
\lean{CategoryTheory.ConnectedComponents.typeToCatHomEquiv,
CategoryTheory.Limits.colimit.desc,
CategoryTheory.Limits.colimit.ι_desc}
Let $F:\mathcal B\to\Grpd$ be a functor and let $X$ be a set. Suppose that
for every $b\in\mathcal B$ a functor
\[
  \Lambda_b:F(b)\longrightarrow\operatorname{Disc}(X)
                                                        \tag{D}\label{eq:fibre-discrete-read}
\]
has been constructed, and that for every base arrow $h:a\to b$ these
functors satisfy
\[
  \Lambda_a=F(h)\fatsemi\Lambda_b.                    \tag{N}\label{eq:fibre-read-naturality-hyp}
\]
Then each $\Lambda_b$ descends uniquely to a function
\[
  \lambda_b:\pi_0(F(b))\longrightarrow X,
\]
the functions $\lambda_b$ assemble into a natural transformation
\[
  \lambda:\pi_0\circ F\Longrightarrow\Delta(X),       \tag{$\lambda$}\label{eq:general-component-read}
\]
and the universal property of the colimit produces a unique function
\[
  \bar\lambda:
  \operatorname*{colim}_{\mathcal B}(\pi_0\circ F)
  \longrightarrow X                                  \tag{$\bar\lambda$}\label{eq:general-colimit-read}
\]
whose composite with the colimit injection from the $b$th fibre is
$\lambda_b$. Equivalently, there is a canonical bijection
\[
 \operatorname{Nat}\bigl(\pi_0\circ F,\Delta(X)\bigr)
 \;\cong\;
 \operatorname{Hom}_{\mathbf{Set}}
 \!\left(\operatorname*{colim}_{\mathcal B}(\pi_0\circ F),X\right).
                                                        \tag{U}\label{eq:component-colimit-universal}
\]
Here $\operatorname{Disc}(X)$ is the discrete category on the underlying
set $X$, and $\Delta(X)$ is the constant $X$-valued diagram. Thus $X$ may
be any set, including the underlying set of $\RR$; no topology on $X$ is
used.

Independently, a category $J$ is nonempty and connected exactly when
$\pi_0(J)$ is a singleton. Consequently, if
$J=\int_{\mathcal B}F$ is nonempty and connected, then
Lemma~\ref{lem:pi0-grothendieck} identifies the component colimit in
\eqref{eq:general-colimit-read} with a singleton.
\end{lemma}
\begin{proof}
For each $b$, the adjunction
$\pi_0\dashv\operatorname{Disc}$ gives
\[
  \operatorname{Hom}_{\mathbf{Set}}\bigl(\pi_0(F(b)),X\bigr)
  \;\cong\;
  \operatorname{Hom}_{\mathbf{Cat}}
    \bigl(F(b),\operatorname{Disc}(X)\bigr).
                                                        \tag{A}\label{eq:fibre-pi0-adjunction}
\]
This is Riehl's connected-components adjunction
\textup{(}\cite[\S8.3, Rem.~8.3.5]{Riehl}\textup{)}. In Mathlib the
equivalence is
\[
  \texttt{ConnectedComponents.typeToCatHomEquiv}.
\]
Applying it to $\Lambda_b$ gives $\lambda_b$. Equation
\eqref{eq:fibre-read-naturality-hyp} says that for every $h:a\to b$ the
square
\[
\begin{tikzcd}[column sep=huge,row sep=large]
\pi_0(F(a)) \arrow[r,"{\pi_0(F(h))}"] \arrow[d,"{\lambda_a}"'] &
\pi_0(F(b)) \arrow[d,"{\lambda_b}"] \\
X \arrow[r,"1_X"'] & X
\end{tikzcd}
\]
commutes. These are precisely the naturality squares for
\eqref{eq:general-component-read}.

A natural transformation
$\lambda:\pi_0\circ F\Rightarrow\Delta(X)$ is a cocone from
$\pi_0\circ F$ to $X$. A colimit is the universal such cocone: every cocone
factors through its colimit cocone by one unique map
\textup{(}\cite[\S3.1, Defs.~3.1.5--3.1.6]{RiehlContext}\textup{)}.
Applying this universal property to $\lambda$ gives $\bar\lambda$ and the
factorization asserted above; varying $\lambda$ gives the bijection
\eqref{eq:component-colimit-universal}. In Mathlib, the induced map is
\texttt{Limits.colimit.desc}, and its value on each colimit injection is
\texttt{Limits.colimit.}$\iota$\texttt{\_desc} \cite{Mathlib}.

Finally, Riehl's Remark~8.3.5 gives the remaining criterion. A category is
nonempty and connected exactly when its connected-components set is a
singleton. Combining that fact with Lemma~\ref{lem:pi0-grothendieck} gives the
last assertion.
\end{proof}

\begin{definition}[$A$-sections]\label{def:A-section}
\lean{ASection, weierstrassE, spherePrimary}
\uses{def:R, prop:weierstrass, lem:residue-spheres}
An \emph{$A$-section} is a section $A$ of the ring $\mathcal R$ of slice-preserving
slice-regular functions on $\OO^{*}=S^8$ \textup{(Definition~\ref{def:R})} with the following
four properties.
\begin{enumerate}[label=\textup{(C\arabic*)},leftmargin=3.2em]
\item \emph{Meromorphic continuation; single simple pole.} $A$ is meromorphically continued to
  $\OO^{*}$ with exactly one pole; it is simple and lies at a finite real point $p_A$, and
  the pole-cancelled factor $(q-p_A)A(q)$ continues through $p_A$ with a finite nonzero
  value $u_A$.  This is entirely a finite-point condition.
\item \emph{Infinite Euler product.} On a slice right half-space $\Omega_0\subset\OO^{*}$,
  $A$ exponentiates one prime-indexed degenerate logarithmic base --- a family \emph{made
  of} the primes, not merely indexed by them:
  $A=\exp\!\big(\sum_{p}\ell_p\big)$, the sum over the infinite set of primes,
  where each $\ell_p$ is slice-preserving, slice-regular, zero-free on $\Omega_0$, and of
  \emph{prime-power logarithmic form}
  \[
    \ell_p(q)=\sum_{k\ge1}\frac{a_{p,k}}{k}\,p^{-kq},
  \]
  with real coefficients $a_{p,k}$ bounded uniformly in $p$ and $k$, the family summable
  \emph{locally normally} on $\Omega_0$ \textup{(}the sense in which the cited Euler products
  converge; Theorem~\ref{thm:euler}\textup{)}. In particular $A$ is zero-free on $\Omega_0$;
  every term decays as $\re q\to+\infty$, so $A\to1$ there; and the logarithm of $A$ is \emph{one} exponential base
  --- prime-power frequencies, constant-scale coefficients, no translated anchor, no additive
  pole term, the pole of \textup{C1} arising instead from the divergence of the prime sum at
  the boundary of the half-space, as for $\zeta$ at $s=1$ --- the single base the unique tame
  continuous lift carries on the degenerate set (Theorems~\ref{thm:log-manifold},
  \ref{thm:winding-lift}). The prime indexes its term through its own powers, not as a label:
  each $\ell_p$ is a power series in $p^{-q}$ for its own prime $p$, so a family merely
  enumerated by primes, whose terms are not such power series, is not of this form.
\item \emph{Infinite Weierstrass factorization.} Over the \emph{full divisor} of $A$,
  including the single pole of \textup{(C1)}, at the real point $p_0$, multiplicity $-1$,
  carried by the explicit left factor, on $\OO^{*}\setminus\{p_0\}$,
  \[
    (q-p_0)\,A\;=\;q^{m}\,R\;e^{g}\prod_n\mathcal E(\,\cdot\,;q_n),
  \]
  where $q$ is the octonionic coordinate, $m\ge0$, $R$ is the slice-regular Weierstrass product
  over the \emph{nonzero} residue-$\RR$ \textup{(}real\textup{)} zeros of $A$, vanishing only
  on $\RR$, $g$ is slice-preserving entire, $\{q_n\}$ enumerates the residue-$\CC$ zero-spheres
  of $A$, and $\mathcal E(\,\cdot\,;q_n)$ is the slice-regular Weierstrass primary factor of
  $q_n$ \textup{(}Proposition~\ref{prop:weierstrass}; \cite[Prop.~3.1, Thm.~3.2]{AdFslice},
  \cite{GV}\textup{)}; the factorization holds with either family infinite, the product
  converging \emph{locally normally} on the finite domain
  $\OO\setminus\{p_0\}$ \textup{(}the convergence of
  Proposition~\ref{prop:weierstrass}\textup{)}. \textup{(}The left
  factor makes ``full divisor'' literal: the divisor of a meromorphic section includes its
  single pole; classically, Hadamard factors $(s-1)\zeta(s)$.\textup{)}
\item \emph{Infinitely many residue-$\CC$ zeros.} The index set $\{q_n\}$ is infinite. A closed
  zero of such a section is a real point \textup{(}residue field $\RR$\textup{)} or a $6$-sphere
  $S_{(\sigma,\gamma)}$ \textup{(}residue field $\CC$; \cite{AdFslice}\textup{)}; residue-$\CC$
  names the classically nontrivial zeros \textup{(}Part~2\textup{)}.
\end{enumerate}
\end{definition}

\begin{definition}[The residue-$\CC$ subdiagram and its inclusion]\label{def:residue-subdiagram}
\lean{ASection.CResidueZeroLocus, ASection.IsCResidueState,
ASection.AsectionCResidueDiagram,
ASection.AsectionCResidueInclusion}
\notready
\uses{def:base, def:A-section, lem:residue-spheres}
By this point conditions \textup{C1--C4} have been used to construct the
three functors that enter the residue restriction, together with the
Grothendieck total of $\mathcal A_A$:
\[
\begin{gathered}
 A^{\mathrm{slice}}:\mathcal B\longrightarrow\SphereWorld,
 \qquad
 A_{\OO}:\mathcal H_1\longrightarrow\mathcal H_1,\\
 \mathcal A_A:\mathcal B\longrightarrow\mathbf{Grpd},
 \qquad
 \mathcal T_A=\int_{\mathcal B}\mathcal A_A.
\end{gathered}                                      \tag{A}\label{eq:residue-built-objects}
\]
The first three displayed constructions are functors.  The fourth is the
Grothendieck total of $\mathcal A_A$: $\mathcal T_A$ is the category obtained
by assembling its fibres over $\mathcal B$, and
Proposition~\ref{prop:total-groupoid} proves that this category is a
groupoid.  It comes with the canonical projection to $\mathcal B$, which
sends $(b,x)$ to $b$ and $(h,\gamma)$ to $h$: the forgetful functor of the
total construction
(\texttt{CategoryTheory.Grothendieck.forget};
Proposition~\ref{prop:total-groupoid}).

The functor $A^{\mathrm{slice}}$ supplies the positioned slice state at each
object of the projective base, and $A_{\OO}$ realizes that state by the
slice-preserving A-section. Thus an object of the fibre $\mathcal A_A(b)$ is
the already constructed triple
\[
 x=\Bigl((I,u),\ A^{\mathrm{slice}}(b)\mathbin{\cdot}(I,u),\
 A_{\OO}\bigl(A^{\mathrm{slice}}(b)\mathbin{\cdot}(I,u)\bigr)\Bigr).
                                                        \tag{S}\label{eq:residue-state-recap}
\]

The analytic work used to construct these triples has already been done.
On the right half-space of \textup{C2}, where the Euler product has no
zeros,
\[
 \Gamma_A(z)=\sum_p\ell_p(z),\qquad
 M(\Gamma_A(z))
 =\operatorname{diag}\!\left(\exp\!\left(\sum_p\ell_p(z)\right),1\right)
 =\operatorname{diag}(A(z),1),
                                                        \tag{E}\label{eq:residue-euler-matrix}
\]
each $\ell_p$ the prime-power series of \textup{C2},
$\ell_p(z)=\sum_{k\ge1}\tfrac{a_{p,k}}{k}\,p^{-kz}$: the one degenerate
base is made of the primes, not merely indexed by them.
Condition \textup{C1} continues the pole-cancelled factor through its finite
real pole $p_A$, and \textup{C3} presents that same factor in Weierstrass
form. Evaluation at $p_A$ gives the invertible multiplier from which the
A-specific disk action was constructed:
\[
\begin{aligned}
 g_A(z)&=(z-p_A)\exp\!\left(\sum_p\ell_p(z)\right),\\
 u_A&=g_A(p_A),\qquad
 D_A=\operatorname{cayleyProjective}
       \bigl(\operatorname{diag}(u_A,1)\bigr),\\
 \exp^{-1}(u_A)=\{\lambda_A+2\pi ik:k\in\ZZ\}.
\end{aligned}                                         \tag{EW}\label{eq:residue-weierstrass-recap}
\]
The last set has the one real level $\operatorname{Re}\lambda_A$, already
built into the action total (Remark~\ref{rmk:gpv-real-fibre}). These analytic
facts belong to the construction of $A^{\mathrm{slice}}$ and
$\mathcal A_A$. The residue restriction therefore begins with the semantic
zero-sphere states that these functors already contain.

The residue subdiagram is defined by two data, both already standing: the
semantic C-residue zero locus, which selects the objects, and the natural
inclusion into $\mathcal A_A$, whose components are the identities of the
$A$-specific fibre groupoids.  The hypotheses' one action supplies the
arrows and transports in $\mathcal A_A$; the residue construction retains
them on the selected objects.  The naturality square is verified on those
object and arrow assignments.  This definition states those data.  That they define a functor
$\mathcal R_A:\mathcal B\to\mathbf{Grpd}$ and that the inclusion is full
and faithful --- fibrewise and on the totals, an inclusion of a subgroupoid
--- is proved below, at Lemma~\ref{lem:iota-full-faithful}.

\emph{The selected objects.}
The semantic C-residue zero locus has already been constructed from the
divisor of $A$:
\[
 \{z\in\CC:A(z)=0,\ \operatorname{Im}z>0\}.
                                                        \tag{Z}\label{eq:semantic-c-residue-locus}
\]
By Lemma~\ref{lem:residue-spheres}, each point of
\eqref{eq:semantic-c-residue-locus} determines its whole $6$-sphere as the
$G_2$-orbit of its slice direction. Condition \textup{C3} constructs these
zero-sphere states, and condition \textup{C4} says that their family is
infinite. In the triple \eqref{eq:residue-state-recap}, the residue
condition is the following one statement: the positioned entry
\[
 A^{\mathrm{slice}}(b)\mathbin{\cdot}(I,u)
\]
is one of the C-residue zero-sphere states already constructed from
\eqref{eq:semantic-c-residue-locus}. Its determined realization by
$A_{\OO}$ is the third entry of the same triple.

For a base object $a$, select a state $x\in\mathcal A_A(a)$ precisely when
there are an already constructed C-residue zero-sphere state in the fibre at the pole,
\[
  x_0\in\mathcal A_A(p_A),
\]
and a base morphism $h_1:p_A\to a$ satisfy the concrete action-image equation
\[
  \mathcal A_A(h_1)(x_0)=x.                       \tag{W}\label{eq:residue-action-image}
\]
The fibre at $p_A$ is where the hypotheses constructed these states:
\textup{C1} continues the factor through the pole, and \textup{C3} presents
it there in Weierstrass form. The morphism $h_1$ is evaluated by the already
constructed functor $\mathcal A_A(h_1)$, so its projective/M\"obius action
and its GPV lift are the A-specific ones already present in
$\mathcal T_A$. This is the action-category inverse image: the semantic
residue condition supplies the zero-sphere states at the pole, and the
action of $\mathcal A_A$ supplies all of their images over the projective
base.

\emph{The objects of the residue fibre.}
Fix a base object $a$.  The action-image equation \textup{(W)} gives the
object collection
\[
\begin{aligned}
 \operatorname{Ob}\mathcal R_A(a):=
 \bigl\{x\in\operatorname{Ob}\mathcal A_A(a):{}&
   \text{there are a C-residue zero-sphere state}\\
 &x_0\in\operatorname{Ob}\mathcal A_A(p_A)
   \text{ and an arrow }h_1:p_A\to a,\\
 &\mathcal A_A(h_1)(x_0)=x\bigr\}.
\end{aligned}                                           \tag{O}\label{eq:residue-fibre-objects}
\]
Thus \eqref{eq:residue-fibre-objects} says precisely that
$\mathcal R_A(a)$ contains the A-specific transports to $a$ of the
already-constructed pole-fibre residue states.

One distinction organizes the restriction.  The initial zero-sphere states
in \textup{(W)} lie in the one pole fibre; their A-specific transports lie
in the other fibres.  What distinguishes
two initial residue states is their fibre data $(I,u)$: the direction and the
normalized input whose positioned image is the state's zero-sphere
coordinate, non-real by
\eqref{eq:semantic-c-residue-locus}.  The prime-power sum varies along the
domain path that built $D_A$ (Definition~\ref{def:gpv-transport}), and the
zero-sphere coordinates stand as the fixed inputs on which the positioned
action acts.

\emph{The retained arrows.}  Let $x$ and $y$ be
two objects of the displayed collection.  Define
\[
 \operatorname{Hom}_{\mathcal R_A(a)}(x,y)
 :=\bigl\{\text{$A$-specific arrows }\gamma:x\to y
          \text{ in }\mathcal A_A(a)\bigr\}.            \tag{H0}\label{eq:residue-hom-definition}
\]
This definition places the residue restriction on objects: once the source
and target objects are selected, every $G_2$ arrow between them in the
already constructed fibre $\mathcal A_A(a)$ is retained.

The identity of a selected object $x$ is the identity $1_x$ in
$\mathcal A_A(a)$.  Let
$\gamma:x\to y$ and $\delta:y\to z$ be retained arrows --- $\gamma\cdot
x=y$ and $\delta\cdot y=z$ --- their composite is the composite in
$\mathcal A_A(a)$:
composition in the action groupoid $\mathcal A_A(a)$ is multiplication in
$G_2$ (Definition~\ref{def:states}), so
\[
  \gamma\fatsemi\delta=\delta\gamma:\ x\longrightarrow z,
  \qquad
  (\delta\gamma)\cdot x=\delta\cdot(\gamma\cdot x)=\delta\cdot y=z.
\]
Its source and target objects $x$ and $z$ are selected, so
\eqref{eq:residue-hom-definition} places this composite in
$\mathcal R_A(a)$.  Likewise, the inverse of $\gamma:x\to y$ is the
group inverse $\gamma^{-1}:y\to x$ in $\mathcal A_A(a)$ --- inverting the defining
equation, $\gamma^{-1}\cdot y=x$ --- whose source and target objects are again
selected.  The category and groupoid equations are the group's own, already
proved in $\mathcal A_A(a)$:
\[
 (\gamma\fatsemi\delta)\fatsemi\varepsilon
   =\gamma\fatsemi(\delta\fatsemi\varepsilon),\qquad
 \gamma\fatsemi\gamma^{-1}=1_x,\qquad
 \gamma^{-1}\fatsemi\gamma=1_y,
\]
associativity of multiplication in $G_2$ and the two inverse equalities at
the relevant objects.  Therefore the displayed objects, hom-sets, identities,
composites, and inverses construct a groupoid $\mathcal R_A(a)$.  This
constructed groupoid is what is meant by the \emph{full subgroupoid} of
$\mathcal A_A(a)$ on the selected objects.

\emph{The restricted transports.}
Fix a base arrow $h_2:a\to b$.  On a selected object $x$ of
$\mathcal R_A(a)$, use the $A$-specific transport:
\[
  \mathcal R_A(h_2)(x):=\mathcal A_A(h_2)(x).           \tag{Obj}\label{eq:residue-object-restriction}
\]
The right side is again a selected object.  Indeed, write the action-image
equation for $x$ using $x_0\in\mathcal A_A(p_A)$ and
$h_1:p_A\to a$.  Then the composite base morphism
$h_1\fatsemi h_2:p_A\to b$ supplies the transported state by functoriality:
\[
\begin{aligned}
  \mathcal A_A(h_1\fatsemi h_2)(x_0)
    &=\mathcal A_A(h_2)\bigl(\mathcal A_A(h_1)(x_0)\bigr)
      &&\text{because $\mathcal A_A$ preserves composition,}\\
    &=\mathcal A_A(h_2)(x)
      &&\text{by \eqref{eq:residue-action-image}.}
\end{aligned}
\]
On a retained arrow $\gamma:x\to y$, use the same $A$-specific
transport:
\[
 \mathcal R_A(h_2)(\gamma):=\mathcal A_A(h_2)(\gamma),
                                                        \tag{A}\label{eq:residue-arrow-restriction}
\]
an arrow in $\mathcal A_A(b)$ between the two transported objects just shown to be
selected, hence retained by \eqref{eq:residue-hom-definition}.  Equations
\eqref{eq:residue-object-restriction} and
\eqref{eq:residue-arrow-restriction} are the $A$-specific assignments read on
the selected states.

\emph{The natural inclusion.}
Define the fibre inclusion
$\iota_{A,a}:\mathcal R_A(a)\hookrightarrow\mathcal A_A(a)$ by
\[
  \iota_{A,a}(x):=x,\qquad
  \iota_{A,a}(\gamma):=\gamma.                         \tag{Inc}\label{eq:residue-fibre-inclusion}
\]
Because the operations in $\mathcal R_A(a)$ are those of $\mathcal A_A(a)$,
these two assignments preserve identities and composition and hence define
a functor.  For a base arrow $h_2:a\to b$, the naturality square is
\[
\begin{tikzcd}[column sep=huge,row sep=large]
\mathcal R_A(a) \arrow[r,"\mathcal R_A(h_2)"] \arrow[d,hook,"\iota_{A,a}"'] &
\mathcal R_A(b) \arrow[d,hook,"\iota_{A,b}"]\\
\mathcal A_A(a) \arrow[r,"\mathcal A_A(h_2)"'] & \mathcal A_A(b),
\end{tikzcd}
\]
Evaluated at an object, this square is equation
\eqref{eq:residue-object-restriction}; evaluated at an arrow, it is equation
\eqref{eq:residue-arrow-restriction}.  On objects both routes apply the
same projective/M\"obius matrix action $\mathcal A_A(h_2)$, and on arrows
both routes retain the same $G_2$ matrix.  Thus the two routes agree on every
object and arrow.  These component squares, one for every $h_2$, are the
natural inclusion $\iota_A:\mathcal R_A\Rightarrow\mathcal A_A$.

The definition is complete: the locus selects the objects
\eqref{eq:residue-fibre-objects}, the selected source and target objects retain the arrows
\eqref{eq:residue-hom-definition}, the transports are the restrictions
\eqref{eq:residue-object-restriction}--\eqref{eq:residue-arrow-restriction},
and the inclusion is componentwise the identity
\eqref{eq:residue-fibre-inclusion}.  That these assignments define a
functor $\mathcal R_A:\mathcal B\to\mathbf{Grpd}$, and that $\iota_A$ is a
full and faithful inclusion of a subgroupoid --- each fibre inclusion a
bijection on hom-sets, and the induced inclusion of totals full and
faithful --- is Lemma~\ref{lem:iota-full-faithful}.
\end{definition}


\begin{lemma}[The residue-$\CC$ inclusion]\label{lem:iota-full-faithful}
\lean{ASection.AsectionCResidueInclusionTotal,
ASection.AsectionCResidueInclusionTotal_full,
ASection.AsectionCResidueInclusionTotal_faithful}
\notready
\uses{def:A-section, def:base, def:residue-subdiagram}
For an A-section $A$, the assignments of
Definition~\ref{def:residue-subdiagram} define a functor
$\mathcal R_A:\mathcal B\to\mathbf{Grpd}$, the fibre inclusions
$\iota_{A,a}$ are full and faithful and assemble into the natural
inclusion $\iota_A:\mathcal R_A\Rightarrow\mathcal A_A$, and the total
inclusion
\[
\begin{tikzcd}[column sep=huge]
\int_{\mathcal B}\mathcal R_A \arrow[r, hook, "\int\iota_A"] & \int_{\mathcal B}\mathcal A_A=\mathcal T_A
\end{tikzcd}
\]
is full and faithful: an isomorphism onto its image, which consists exactly of the
selected residue states together with the same projective/M\"obius transports and
$G_2$ fibre arrows between them.
\end{lemma}
\begin{proof}
Everything is inherited from the already-constructed action functor
$\mathcal A_A$, after restricting its objects.  The closure calculation in
Definition~\ref{def:residue-subdiagram} shows that every
$\mathcal A_A(h)$ carries selected objects to selected objects, so its
restriction $\mathcal R_A(h)$ is defined.  Since it uses exactly the same
object map and the same $G_2$ matrix on each arrow, the identity calculation is
\[
 \mathcal R_A(1_a)(x)=\mathcal A_A(1_a)(x)=x,\qquad
 \mathcal R_A(1_a)(\gamma)=\mathcal A_A(1_a)(\gamma)=\gamma,
\]
and for $h_2:a\to b$ and $h_3:b\to c$ the composite calculation, on an
object $x$, is
\[
 \mathcal R_A(h_2\fatsemi h_3)(x)
  =\mathcal A_A(h_2\fatsemi h_3)(x)
  =\mathcal A_A(h_3)\bigl(\mathcal A_A(h_2)(x)\bigr)
  =\mathcal R_A(h_3)\bigl(\mathcal R_A(h_2)(x)\bigr),
\]
with the identical chain on a $G_2$ matrix $\gamma$.  Thus identities go to
identities, composites go to composites, and
$\mathcal R_A:\mathcal B\to\mathbf{Grpd}$ is a functor.

For two selected objects $x,y$ of $\mathcal R_A(a)$, the map induced by
$\iota_{A,a}$ is
\[
\begin{aligned}
 \operatorname{Hom}_{\mathcal R_A(a)}(x,y)
 &\xrightarrow{\ \sim\ }
 \operatorname{Hom}_{\mathcal A_A(a)}
   \bigl(\iota_{A,a}x,\iota_{A,a}y\bigr),\\
 \gamma&\longmapsto\iota_{A,a}(\gamma).
\end{aligned}
                                                        \tag{H}\label{eq:residue-fibre-hom}
\]
The map is injective because it sends each arrow to that same arrow in
$\mathcal A_A(a)$.  It is surjective because
\eqref{eq:residue-hom-definition} retained every arrow in $\mathcal A_A(a)$
between the selected source and target objects.  This proves the
bijection \textup{(H)}: each $\iota_{A,a}$ is full and faithful.

For a base arrow $h:a\to b$, the square displayed in
Definition~\ref{def:residue-subdiagram} commutes for the same concrete
reason.  On an object, both routes apply the projective/M\"obius matrix
action $\mathcal A_A(h)$; on a fibre arrow, both retain its $G_2$ matrix.
Therefore the fibre inclusions are the components of the natural
transformation $\iota_A:\mathcal R_A\Rightarrow\mathcal A_A$.  On a
selected object $x$, this says
\[
  \mathcal A_A(h)(\iota_{A,a}x)
   =\iota_{A,b}\bigl(\mathcal R_A(h)(x)\bigr).
                                                        \tag{Src}
\]

Now fix $(a,x)$ and $(b,y)$ in the residue total.  A morphism between them
is exactly a pair consisting of a projective matrix $h:a\to b$ and a
$G_2$ matrix
$\gamma:\mathcal R_A(h)(x)\to y$ in the fibre at $b$.  By \textup{(Src)},
the induced total inclusion sends this pair to
\[
  (h,\gamma)\longmapsto
  \bigl(h,\iota_{A,b}(\gamma)\bigr).                    \tag{Tot}
\]
This is faithful: equality after \textup{(Tot)} first equates the projective
matrices $h$, then the injectivity in \textup{(H)} equates the $G_2$
matrices.  It is full: given a target pair $(h,\alpha)$, keep its concrete
base matrix $h$; equation \textup{(Src)} identifies the source of $\alpha$
with the included transported source, and surjectivity in \textup{(H)} says
that this same $G_2$ matrix is retained in $\mathcal R_A(b)$.  Hence
$(h,\alpha)$ has a preimage.  The total inclusion is full and faithful, and
its image is exactly the selected residue states with the same A-specific
projective/M\"obius and $G_2$ matrix actions between them.
\end{proof}

\begin{lemma}[Transitivity of the residue-$\CC$ total]\label{lem:c-residue-transitive}
\notready
\uses{def:A-section, def:base, def:residue-subdiagram, prop:weierstrass,
lem:residue-spheres, lem:normalized-pole-transitive, thm:G2-S6}
For an A-section $A$, the total inverse-image groupoid
\[
  \int_{\mathcal B}\mathcal R_A
\]
is transitive: any two of its objects are joined by a morphism.
\end{lemma}
\begin{proof}
Let $P=(a,x)$ and $Q=(b,y)$ be arbitrary objects of
$\int_{\mathcal B}\mathcal R_A$.  Their fibre entries are the complete
A-section triples
\[
\begin{aligned}
x&=\Bigl((I_P,u_P),\
 A^{\mathrm{slice}}(a)\mathbin{\cdot}(I_P,u_P),\
 A_{\OO}\bigl(A^{\mathrm{slice}}(a)\mathbin{\cdot}(I_P,u_P)\bigr)\Bigr),\\
y&=\Bigl((I_Q,u_Q),\
 A^{\mathrm{slice}}(b)\mathbin{\cdot}(I_Q,u_Q),\
 A_{\OO}\bigl(A^{\mathrm{slice}}(b)\mathbin{\cdot}(I_Q,u_Q)\bigr)\Bigr).
\end{aligned}
                                    \tag{T}\label{eq:arbitrary-residue-triples}
\]
We now use the membership data of these two particular objects. By
Definition~\ref{def:residue-subdiagram}, there are semantic C-residue
zero-sphere states $x_0,y_0\in\mathcal A_A(p_A)$, both in the fibre at the
pole --- the fibre in which the locus constructs them --- and A-specific
base morphisms
\[
  h_1:p_A\longrightarrow a,
  \qquad h_2:p_A\longrightarrow b,
  \qquad
  \mathcal A_A(h_1)(x_0)=x,
  \quad
  \mathcal A_A(h_2)(y_0)=y.                    \tag{W}\label{eq:PQ-action-images}
\]
The states $x_0$ and $y_0$ are the semantic C-residue zero-sphere states
whose transports are the two given objects. Write their complete triples
as
\[
\begin{aligned}
x_0&=\Bigl((I_1,u_1),\
 A^{\mathrm{slice}}(p_A)\mathbin{\cdot}(I_1,u_1),\
 A_{\OO}\bigl(A^{\mathrm{slice}}(p_A)\mathbin{\cdot}(I_1,u_1)\bigr)\Bigr),\\
y_0&=\Bigl((I_2,u_2),\
 A^{\mathrm{slice}}(p_A)\mathbin{\cdot}(I_2,u_2),\
 A_{\OO}\bigl(A^{\mathrm{slice}}(p_A)\mathbin{\cdot}(I_2,u_2)\bigr)\Bigr).
\end{aligned}                                      \tag{R}\label{eq:residue-seed-triples}
\]
The middle entries in this display are the positioned states of two members
of the already constructed C-residue zero-sphere family.  Thus
$I_1,I_2\in S^6$, and their normalized complex coordinates lie in the
semantic locus \eqref{eq:semantic-c-residue-locus}.  These displayed entries
are exactly their object data in the action groupoid.

The inverse transports in \eqref{eq:PQ-action-images} reduce the problem to
constructing a morphism inside the one fibre total position at the pole,
\[
   (p_A,x_0)\longrightarrow(p_A,y_0).                    \tag{C}\label{eq:residue-central-comparison}
\]
We construct its base component from the two state-specific evaluations of
the one A-specific action.  The coordinates $u_1$ and $u_2$ lie in the
same residue-$\CC$ half-plane.  Lemma~\ref{lem:normalized-pole-transitive}
therefore gives an invertible matrix morphism $k:p_A\to p_A$ whose action
on the normalized coordinate is
\[
 A^{\mathrm{slice}}(p_A)^{-1}\mathbin{\cdot}
 \bigl(A^{\mathrm{slice}}(k)\mathbin{\cdot}
 (A^{\mathrm{slice}}(p_A)\mathbin{\cdot}u_1)\bigr)=u_2.
                                                               \tag{I}\label{eq:residue-input-comparison}
\]
Every term in this equation is the matrix action of a morphism in the
corresponding groupoid category.  Equation
\eqref{eq:residue-input-comparison} is exactly the normalized-input entry
of Definition~\ref{def:transport}, specialized to the actual arrow
$k:p_A\to p_A$.  The corresponding square is therefore
\[
\begin{tikzcd}[column sep=huge,row sep=large]
(I_1,u_1)
  \arrow[r,"{A^{\mathrm{slice}}(p_A)^{-1}
    A^{\mathrm{slice}}(k)A^{\mathrm{slice}}(p_A)}"]
  \arrow[d,"{A^{\mathrm{slice}}(p_A)}"'] &
(I_1,u_2) \arrow[d,"{A^{\mathrm{slice}}(p_A)}"] \\
A^{\mathrm{slice}}(p_A)\mathbin{\cdot}(I_1,u_1)
  \arrow[r,"{A^{\mathrm{slice}}(k)}"'] &
A^{\mathrm{slice}}(p_A)\mathbin{\cdot}(I_1,u_2).
\end{tikzcd}                                                   \tag{Sq}\label{eq:residue-comparison-square}
\]
Its commutativity is the cancellation already displayed in
Definition~\ref{def:transport}, written here on these two particular
objects:
\[
\begin{aligned}
 A^{\mathrm{slice}}(k)\cdot
   \bigl(A^{\mathrm{slice}}(p_A)\cdot u_1\bigr)
 &=A^{\mathrm{slice}}(p_A)\cdot
   \left(A^{\mathrm{slice}}(p_A)^{-1}\cdot
   \bigl(A^{\mathrm{slice}}(k)\cdot
   (A^{\mathrm{slice}}(p_A)\cdot u_1)\bigr)\right)\\
 &=A^{\mathrm{slice}}(p_A)\cdot u_2,
\end{aligned}
\]
the first equality being inverse cancellation and the second the
concrete input equality \eqref{eq:residue-input-comparison}.  Hence the bottom arrow in
\eqref{eq:residue-comparison-square} sends the positioned entry of $x_0$ to
the positioned entry of $y_0$.  This is the
point at which the proof consumes the Euler--Weierstrass/GPV action already
bound into $\mathcal A_A$.

Thus $\mathcal A_A(k)(x_0)$ and $y_0$ have the same normalized complex
coordinate and the same positioned C-residue coordinate. The remaining
comparison in their triples is the slice direction. Since
$I_1,I_2\in S^6$, Theorem~\ref{thm:G2-S6} supplies
$\gamma\in G_2$ with $\gamma I_1=I_2$. It sends the positioned zero-sphere
state in direction $I_1$ to the state with the same complex coordinate in
direction $I_2$. Equivariance of $A_{\OO}$ sends the realized third entry
by that same element $\gamma$. Hence $\gamma$ is the fibre arrow
\[
  \gamma:\mathcal A_A(k)(x_0)\longrightarrow y_0.
\]
The pair $(k,\gamma)$ is therefore the morphism in
\eqref{eq:residue-central-comparison}.

The three stages now compose in the Grothendieck total:
\[
  P\longrightarrow(p_A,x_0)
   \xrightarrow{\;(k,\gamma)\;}(p_A,y_0)
   \longrightarrow Q,
\]
with base component $(h_1^{-1}\fatsemi k)\fatsemi h_2$.  Thus the proof of
transitivity uses exactly the projective/M\"obius arrow $k$ whose action is
$\mathcal A_A(k)$ and the $G_2$ arrow $\gamma$ between its two residue
states. Since $P$ and $Q$ were arbitrary, the residue total is transitive.
This is where the residue restriction is essential: its realized states are
the C-residue zero spheres, one $G_2$-orbit in each fixed complex
coordinate, whereas the ambient total also contains arbitrary realizations
in $\OO^*$ lying in different $G_2$-orbits.
\end{proof}

\begin{theorem}[Concentricity]\label{thm:concentricity}
\lean{ASection.concentricity, ASection.residueTotalObject,
ASection.residueTotalTransportRead_certified, ASection.transportLevel}
\notready
\uses{def:A-section, def:base, def:residue-subdiagram, lem:iota-full-faithful,
lem:c-residue-transitive,
lem:residue-spheres, lem:pi0-grothendieck, lem:connected-readout,
thm:G2-S6, thm:winding-lift, prop:winding-signature, thm:identity}
Let $A$ be an element of the ring $\mathcal R$ of slice-preserving
functions on the compactified octonions $\OO^{*}$, satisfying conditions
\textup{(C1)--(C4)}. The infinitely many residue-$\CC$ zero $6$-spheres of
$A$ form one connected-component class in the residue action total
$\int_{\mathcal B}\mathcal R_A$; equivalently, their component colimit is a
singleton. That single class has one real readout
\[
  c\in\RR.
\]
Every residue-$\CC$ zero $6$-sphere therefore has centre $c$, so the
infinitely many residue-$\CC$ zero spheres are concentric.
\end{theorem}
\begin{proof}
By conditions \textup{C1--C4}, the following are already constructed:
\[
\begin{gathered}
 A^{\mathrm{slice}}:\mathcal B\longrightarrow\SphereWorld,
 \qquad
 A_{\OO}:\mathcal H_1\longrightarrow\mathcal H_1,\\
 \mathcal A_A:\mathcal B\longrightarrow\mathbf{Grpd},
 \qquad
 \mathcal T_A=\displaystyle\int_{\mathcal B}\mathcal A_A,\\
 \mathcal R_A:\mathcal B\longrightarrow\mathbf{Grpd},
 \qquad
 \displaystyle\int_{\mathcal B}\mathcal R_A
   \xrightarrow{\ \int\iota_A\ }\mathcal T_A.
\end{gathered}                                         \tag{C}\label{eq:concentricity-built-chain}
\]
The analytic data have already entered these constructions. Condition
\textup{C2} supplies the prime-power lift and its matrix,
\[
 \Gamma_A(z)=\sum_p\ell_p(z),\qquad
 M(\Gamma_A(z))=\operatorname{diag}(A(z),1).
\]
Condition \textup{C1} continues the pole-cancelled function $g_A$ through
$p_A$, and \textup{C3} presents that same function in Weierstrass form.
At the finite point $p_A$ the resulting multiplier and disk action are
\[
 g_A(z)=(z-p_A)\exp\!\left(\sum_p\ell_p(z)\right),\qquad
 u_A=g_A(p_A),\qquad
 D_A=\operatorname{cayleyProjective}
       \bigl(\operatorname{diag}(u_A,1)\bigr).          \tag{D}\label{eq:concentricity-action-recap}
\]
The already-proved functor $A^{\mathrm{slice}}$ carries this disk action
through the projective base: its object assignment gives the action at each
base object and its arrow assignment gives the invertible M\"obius matrix
action of each base morphism. The functor $A_{\OO}$ realizes each
positioned state equivariantly, and the action diagram $\mathcal A_A$
combines these two functors in the triples
\[
 \Bigl((I,u),\ A^{\mathrm{slice}}(b)\mathbin{\cdot}(I,u),\
 A_{\OO}\bigl(A^{\mathrm{slice}}(b)\mathbin{\cdot}(I,u)\bigr)\Bigr),
\]
whose Grothendieck total is $\mathcal T_A$. The GPV lift used in these
transports has the exponential fibre
\[
 \exp^{-1}(u_A)=\{\lambda_A+2\pi ik:k\in\ZZ\},
\]
and every member of that fibre has the same real part
$\operatorname{Re}\lambda_A$ (Remark~\ref{rmk:gpv-real-fibre}). Thus the
Euler--Weierstrass continuation, the multiplier, and its one real fibre have
already been consumed in $A^{\mathrm{slice}}$, $\mathcal A_A$, and
$\mathcal T_A$ before the residue restriction is made.

Condition \textup{C3} supplies the semantic C-residue zero-sphere states,
and condition \textup{C4} makes their family infinite.
Definition~\ref{def:residue-subdiagram} forms the full A-specific
inverse-image subdiagram on those states. Lemma~\ref{lem:iota-full-faithful}
gives the literal full and faithful inclusion
\[
 \int_{\mathcal B}\mathcal R_A
   \xrightarrow[\text{full and faithful}]{\ \int\iota_A\ }
 \mathcal T_A .                                      \tag{I}\label{eq:concentricity-total-inclusion}
\]
Its objects are precisely the A-specific images of the displayed
zero-sphere triples, and its morphisms are the projective/M\"obius and
$G_2$ arrows already constructed in $\mathcal A_A$.

Lemma~\ref{lem:c-residue-transitive} constructs a morphism between any two
objects of this exact total. Its base component is the locus-specific
projective matrix $k:p_A\to p_A$ supplied by the concrete real-projective
action \eqref{eq:real-projective-transitivity}, and the commutative square
for the already-constructed functor $A^{\mathrm{slice}}$ maps the two
positioned coordinates. Its fibre component is the matrix in $G_2$ that
maps the remaining slice direction and, by equivariance of $A_{\OO}$, the
realized third entry. Hence
\[
  \int_{\mathcal B}\mathcal R_A
\]
is a transitive groupoid and therefore a connected category.

Connectedness now has two consequences, both for this same residue total.
First, Lemma~\ref{lem:connected-readout} says that a nonempty connected
category has one connected component. Condition \textup{C4} supplies
objects, and connectedness was just established; hence
\[
  \pi_0\Bigl(\int_{\mathcal B}\mathcal R_A\Bigr)
\]
is a singleton. Now apply Lemma~\ref{lem:pi0-grothendieck} to the specific
functor $\mathcal R_A:\mathcal B\to\mathbf{Grpd}$. The resulting
equivalence is
\[
  \pi_0\Bigl(\int_{\mathcal B}\mathcal R_A\Bigr)
  \;\cong\;
  \operatorname*{colim}_{\mathcal B}(\pi_0\circ\mathcal R_A),
\]
so the colimit on the right is the same singleton, which we denote by
$\{\kappa\}$.

Only after this collapse do we apply connected components to the residue
functor itself:
\[
  \pi_0\circ\mathcal R_A:\mathcal B\longrightarrow\mathbf{Set}.
                                                        \tag{$\pi$}\label{eq:residue-pi0-diagram}
\]
For each base object $b$, the intrinsic real face $\operatorname{Re}\Gamma$
already bound into the A-specific states of $\mathcal R_A(b)$ is a functor
from that fibre to the discrete category on $\RR$.  Lemma~\ref{lem:connected-readout},
in the form of Mathlib's
\texttt{ConnectedComponents.typeToCatHomEquiv}, descends this fibrewise
functor to a function
\[
  \lambda_b:\pi_0\bigl(\mathcal R_A(b)\bigr)\longrightarrow\RR.
\]
On a class represented by the action-state triple with normalized input
$(I,u)$, this is the real coordinate of its positioned entry:
\[
 \lambda_b([x])
  =\operatorname{Re}\Gamma(x)
  =\operatorname{Re}\!\left(A^{\mathrm{slice}}(b)\mathbin{\cdot}u\right).
                                                        \tag{V}\label{eq:residue-representative-read}
\]
In particular, for the canonical object belonging to the $n$th C-residue
zero sphere, its state equation identifies the positioned entry with
$A.\operatorname{sphereZero}(n)$, and hence
\[
 \lambda_b([x_n])
   =\operatorname{Re}\bigl(A.\operatorname{sphereZero}(n)\bigr).
                                                        \tag{V$_n$}\label{eq:residue-certified-read}
\]
For a base arrow $h:a\to b$, the A-specific action square used by
$\mathcal R_A(h)$ carries this same bound real face.  Consequently the
square
\[
\begin{tikzcd}[column sep=huge,row sep=large]
\pi_0\bigl(\mathcal R_A(a)\bigr)
  \arrow[r,"{\pi_0(\mathcal R_A(h))}"]
  \arrow[d,"{\lambda_a}"'] &
\pi_0\bigl(\mathcal R_A(b)\bigr)
  \arrow[d,"{\lambda_b}"] \\
\RR \arrow[r,"1_{\RR}"'] & \RR
\end{tikzcd}                                            \tag{N}\label{eq:residue-read-naturality}
\]
commutes. These squares assemble the natural transformation
\[
  \lambda:\pi_0\circ\mathcal R_A\Longrightarrow\Delta(\RR),
                                                        \tag{$\lambda$}\label{eq:residue-component-read}
\]
where $\Delta(\RR)$ is the constant diagram on the set $\RR$.

The universal property of the colimit now applies to the natural
transformation \eqref{eq:residue-component-read}.  It produces the one
function
\[
  \bar\lambda:
  \operatorname*{colim}_{\mathcal B}(\pi_0\circ\mathcal R_A)
  \longrightarrow\RR.                                \tag{$\bar\lambda$}\label{eq:residue-colimit-read}
\]
This is the readout of the component diagram
\eqref{eq:residue-pi0-diagram}.  Its relation to the transitivity argument
is therefore the explicit chain
\[
\begin{tikzcd}[column sep=large,row sep=large]
\pi_0\!\left(\int_{\mathcal B}\mathcal R_A\right)
  \arrow[r,"\cong"] &
\operatorname*{colim}_{\mathcal B}(\pi_0\circ\mathcal R_A)
  =\{\kappa\}
  \arrow[r,"{\bar\lambda}"] & \RR .
\end{tikzcd}                                           \tag{Read}\label{eq:residue-readout-chain}
\]

The colimit injection sends every certified class $[x_n]$ to the sole
element $\kappa$.  Apply $\bar\lambda$ in \eqref{eq:residue-readout-chain}
and use the defining calculation \eqref{eq:residue-certified-read}.  The
$n$th and zeroth representatives therefore give directly
\[
 \operatorname{Re}\bigl(A.\operatorname{sphereZero}(n)\bigr)
 =\bar\lambda(\kappa)
 =\operatorname{Re}\bigl(A.\operatorname{sphereZero}(0)\bigr).
                                                               \tag{C$'$}\label{eq:residue-constant-read}
\]
Set $c=\operatorname{Re}\bigl(A.\operatorname{sphereZero}(0)\bigr)$.
Then \eqref{eq:residue-constant-read} gives
\[
 \operatorname{Re}\bigl(A.\operatorname{sphereZero}(n)\bigr)=c
 \qquad\text{for every }n.                            \tag{R}\label{eq:residue-common-centre}
\]
The real coordinate of a residue-$\CC$ zero sphere is its centre.
Therefore every residue-$\CC$ zero $6$-sphere has centre $c$, and the
infinitely many such spheres are concentric.
\end{proof}

\begin{corollary}[Translation to the classical framework]\label{cor:nontrivial}
\lean{ASection.nontrivial_one_centre}
\uses{thm:concentricity, thm:zero-equivalence, thm:zero-spheres, lem:residue-spheres}
Under the identification of residue-$\CC$ zeros with classically nontrivial zeros, the already-proved
residue-$\CC$ zero spheres are the classically nontrivial zero spheres and retain the
common centre $c$ supplied by Theorem~\ref{thm:concentricity}.
\end{corollary}
\begin{proof}
This is the translation of the residue-$\CC$ family through
Theorems~\ref{thm:zero-equivalence} and~\ref{thm:zero-spheres}. The centre comes directly from
Theorem~\ref{thm:concentricity}, while \textup{C4} and Lemma~\ref{lem:residue-spheres} supply
infinitude and disjointness.
\end{proof}

% The superseded great-circle/F-based proof frame was removed from the active
% master on 2026-07-17. It remains recoverable from git history and is not an
% operative proof or expository alternative.

\section{Corollaries}\label{sec:verification}

Two corollaries close the argument. The first checks that the octonionic zeta is an A-section,
verifying its four conditions C1--C4 one by one against the classical facts of Part~1; the
Concentricity Theorem then applies to it directly, and its residue-$\CC$ zero-spheres share one real
centre $c$. The second corollary is the Riemann Hypothesis. The functional equation of $\xi$ now
pins that centre: it sends a zero of real part $c$ to one of real part $1-c$, and applied to the
common centre it forces $c=1-c$, so $c=\tfrac12$. Every nontrivial zero lies on the critical line.

\begin{corollary}[$\zetaO$ is an $A$-section]\label{cor:zeta-section}
\lean{zetaSection}
\uses{def:A-section, thm:zeta-in-R, thm:riemann, thm:euler, cor:hadamard-infinitude, prop:weierstrass, thm:zero-spheres}
$\zetaO$ is an $A$-section \textup{(Definition~\ref{def:A-section})}: a section of
$\mathcal R$ satisfying \textup{(C1)--(C4)}.
\end{corollary}
\begin{proof}
$\zetaO\in\mathcal R$ by slice-preservation (Theorem~\ref{thm:zeta-in-R}). The verification of
\textup{(C1)--(C4)} cites the classical facts of the Riemann zeta function (Part~1) and the
slice-regular zero-geometry that translates them:
\begin{itemize}[leftmargin=2em]
\item \textbf{(C1)} \emph{Meromorphic continuation; simple pole.} Classical fact: $\zeta$ is
  meromorphically continued with a single simple pole at $s=1$ (Theorem~\ref{thm:riemann});
  its pole-cancelled factor $(s-1)\zeta(s)$ continues through $s=1$ with the finite nonzero
  value $1$.  The slice-preserving continuation gives the same statement for
  $(q-1)\zetaO(q)$.
\item \textbf{(C2)} \emph{Infinite Euler product.} Classical fact: $\zeta=\prod_p(1-p^{-s})^{-1}
  =\exp(\sum_p\ell_p)$ on $\Re s>1$ with $\ell_p(s)=\sum_{k\ge1}p^{-ks}/k$ --- prime-power
  logarithmic form with every coefficient $a_{p,k}=1$, one base, decaying to $\zeta\to1$ as
  $\Re s\to+\infty$ (Theorem~\ref{thm:euler}).
\item \textbf{(C3)} \emph{Infinite Weierstrass factorization.} Supplied by the
  \emph{slice-regular} Weierstrass factorization (Proposition~\ref{prop:weierstrass};
  \cite[Prop.~3.1, Thm.~3.2]{AdFslice}, \cite{GV}):
  $(q-1)\,\zetaO=q^{m}Re^{g}\prod_n\mathcal E(\cdot;q_n)$ \textup{(}pole factor at $p_0=1$;
  classically, Hadamard factors $(s-1)\zeta(s)$\textup{)}, where $m=0$
  \textup{(}$\zeta(0)=-\tfrac12\neq0$\textup{)}, $R$ is the product over the nonzero
  residue-$\RR$ zeros, the classically trivial zeros $-2,-4,\dots$, honest real zeros of
  $\zetaO$, and $\{q_n\}$ enumerates its residue-$\CC$ zero-spheres.
\item \textbf{(C4)} \emph{Infinitely many residue-$\CC$ zeros.} Classical fact: $\zeta$ has
  infinitely many nontrivial zeros (Corollary~\ref{cor:hadamard-infinitude}; also Hardy,
  Theorem~\ref{thm:hardy}). By the residue dictionary below these are exactly the residue-$\CC$
  zeros, so $\{q_n\}$ is infinite.
\end{itemize}
After the \textup{(C1)--(C4)} verification, the downstream centre-pinning uses the
\emph{functional equation} $\xi(s)=\xi(1-s)$,
$\xi(\bar s)=\overline{\xi(s)}$ (Theorem~\ref{thm:riemann}), used in Corollary~\ref{cor:rh}. The
\emph{residue dictionary}, a closed zero of $\zetaO$ is a real point (residue field $\RR$, the
classically \emph{trivial} zeros) or a $6$-sphere $S_{(\sigma,\gamma)}$ (residue field $\CC$, the
classically \emph{nontrivial} zeros), is the slice-regular zero-geometry of
Theorem~\ref{thm:zero-spheres} and \cite{AdF}; it is what makes ``residue-$\CC$'' and
``nontrivial'' the same notion. Thus \textup{(C1)--(C4)} hold.
\end{proof}

\begin{corollary}[Riemann Hypothesis]\label{cor:rh}
\lean{zeta_riemannHypothesis, zeta_criticalLine_zeros_infinite}
\leanok
\uses{cor:zeta-section, thm:concentricity, cor:nontrivial, thm:rh-equiv}
Every nontrivial zero of the classical Riemann zeta function has real part $\tfrac12$; in
particular infinitely many zeros lie on the line $\Re s=\tfrac12$.
\end{corollary}
\begin{proof}
By Corollary~\ref{cor:zeta-section}, $\zetaO$ satisfies \textup{(C1)--(C4)}. By
Theorem~\ref{thm:concentricity}, its residue-$\CC$ value-transport colimit is the real
singleton $\{c\}$ and its residue-$\CC$ zero spheres are already concentric about $c$. By
Corollary~\ref{cor:nontrivial}, these are
the classically nontrivial zeros, so every nontrivial zero has real part $c$. The functional
equation $\xi(s)=\xi(1-s)$ (Theorem~\ref{thm:riemann}) sends a zero of real part $c$ to one of
real part $1-c$; applied to the \emph{common} centre it forces $c=1-c$, whence $c=\tfrac12$ (this
is Theorem~\ref{thm:rh-equiv} read on $\zetaO$). Therefore every nontrivial zero has real part
$\tfrac12$, and by \textup{(C4)} infinitely many lie on $\Re s=\tfrac12$.
\end{proof}

A word on why the argument takes the shape it does. No one can hold an $8$-dimensional graph in the
mind's eye, or draw the continuum of Riemann spheres the section lives on. The epigraph names the
way through: no single viewpoint sees the whole, but when the viewpoints are multiplied and
conjugated, a vision appears that none of them held alone. Each slice direction is one such
viewpoint; the value-bearing A-section functor gathers all of them over the base, the total object holds them
together, and the colimit conjugates them into a single answer. The concentricity we could not
picture becomes one component we can compute.

\begin{remark}[The image span of a fully faithful inclusion]\label{rmk:image-span}
For any fully faithful functor $F:\mathcal C\to\mathcal D$ there is a canonical roof through
the image,
\[
\begin{tikzcd}[column sep=large]
& \operatorname{Im}(F) \arrow[dl, "F^{-1}"'] \arrow[dr, hook, "j"] & \\
\mathcal C & & \mathcal D
\end{tikzcd}
\]
the backward leg being an isomorphism of categories onto the image and the forward leg the
canonical inclusion. It is a general fact, available here for $\iota_A$ and for
$\int\iota_A$, and it is the shape a length-two zigzag argument would use.
Theorem~\ref{thm:concentricity} obtains connectedness directly from transitivity at the total,
by a single element. The image span records the corresponding length-two alternative.
\end{remark}

\begin{remark}[Two alternative routes]\label{rmk:thomason}
Two other proofs of Theorem~\ref{thm:concentricity} are available. The first runs over the
base. The restriction sits over $\mathcal B$ through its two projections,
\[
\begin{tikzcd}[column sep=large]
\int_{\mathcal B}\mathcal R_A \arrow[rr, hook, "\int\iota_A"] \arrow[dr, "\pi_{\mathcal R}"'] & &
\int_{\mathcal B}\mathcal A_A \arrow[dl, "\pi_{\mathcal A}"] \\
& \mathcal B &
\end{tikzcd}
\]
and one may connect residue states footpoint by footpoint: zigzags in the connected base,
lifted through the fibres. That route travels through $\mathcal B$ and reassembles the value data
leg by leg. The second is Thomason's homotopy colimit theorem \cite{Thomason79}: the
classifying space of a Grothendieck construction is homotopy equivalent to the homotopy
colimit of the classifying spaces of its fibres, and reading connected components through that
equivalence yields the same collapse of $\pi_0\bigl(\int_{\mathcal B}\mathcal R_A\bigr)$. The
proof given here stays at the level of categories, where Lemma~\ref{lem:pi0-grothendieck}
suffices: the base route assembles what the action already supplies whole, the Thomason route
is the topological shadow of the same computation, and both remain expository routes outside
the Lean development.
\end{remark}

%=====================================================================
\begin{thebibliography}{99}
%=====================================================================
\bibitem{AdF} A.~Altavilla and C.~De Fabritiis, \emph{$\ast$-logarithm for slice regular
functions}, Rend.\ Lincei Mat.\ Appl.\ (arXiv:2106.04227); the slice-regular exponential
$\exp_{*}$ (Def.~2.14, and Rem.~2.23 for the slice-preserving formula
$\exp_{*}(f)=e^{f_0}(\cos f_1+I\sin f_1)$) and the $\ast$-logarithm on never-vanishing
functions (Def.~2.24, Thm.~6.17).
\bibitem{AdFslice} A.~Altavilla and C.~De Fabritiis, \emph{$s$-regular functions which
preserve a complex slice}, Ann.\ Mat.\ Pura Appl.\ (arXiv:1801.01318); Prop.~3.1 and
Thm.~3.2 factor a slice-regular function into a factor carrying exactly its (real and
spherical) zeros and a never-vanishing factor, including the case of infinitely many zeros,
the citable form of C3.
\bibitem{Baez02} J.~C.~Baez, \emph{The octonions}, Bull.\ Amer.\ Math.\ Soc.\
\textbf{39} (2002), 145--205.
\bibitem{BisiWinkelmann} C.~Bisi and J.~Winkelmann, \emph{Invariants and Automorphisms for slice
regular functions: the octonionic case}, arXiv:2411.16762 (2024); companion arXiv:2411.15896 in
J.~Noncommutative Geometry (EMS). Supplies the $I$-independent lift $\phi_I\circ F=A\circ\phi_I$
(the slice projection) and, in \S3.2/\S3.7/\S14, the axially-symmetric-domain, slice-preserving, and
central-divisor notions.
\bibitem{CSS12} F.~Colombo, I.~Sabadini, and D.~C.~Struppa, \emph{Noncommutative
Functional Calculus: Theory and Applications of Slice Hyperholomorphic Functions},
Progress in Mathematics \textbf{289}, Birkh\"auser, 2011.
\bibitem{GP11} R.~Ghiloni and A.~Perotti, \emph{Slice regular functions on real
alternative algebras}, Adv.\ Math.\ \textbf{226} (2011), 1662--1691.
\bibitem{SeriesExp} R.~Ghiloni, A.~Perotti, and C.~Stoppato, \emph{Singularities of slice regular
functions over real alternative $\ast$-algebras}, Adv.\ Math.\ \textbf{305} (2017), 1085--1130
(slice-function and stem-function formalism, \S1--\S2; the slice-preserving
$\Leftrightarrow$ intrinsic $\Leftrightarrow$ $\RR$-valued-stem characterization used in
Definition~\ref{def:slice-squares}; \S11: slice semiregular functions, Def.~11.1, Thm.~11.3).
\bibitem{GS07} G.~Gentili and D.~C.~Struppa, \emph{A new theory of regular functions
of a quaternionic variable}, Adv.\ Math.\ \textbf{216} (2007), 279--301.
\bibitem{GSS13} G.~Gentili, C.~Stoppato, and D.~C.~Struppa, \emph{Regular Functions of
a Quaternionic Variable}, Springer Monographs in Mathematics, Springer, 2013.
\bibitem{GV} G.~Gentili and I.~Vignozzi, \emph{The Weierstrass factorization theorem
for slice regular functions over the quaternions}, Ann.\ Global Anal.\ Geom.\
\textbf{40} (2011), 435--466.
\bibitem{Hardy14} G.~H.~Hardy, \emph{Sur les z\'eros de la fonction $\zeta(s)$ de
Riemann}, C.\ R.\ Acad.\ Sci.\ Paris \textbf{158} (1914), 1012--1014.
\bibitem{Munkres00} J.~R.~Munkres, \emph{Topology}, 2nd ed., Prentice Hall, 2000.
\bibitem{Riemann1859} B.~Riemann, \emph{\"Uber die Anzahl der Primzahlen unter einer
gegebenen Gr\"osse}, Monatsber.\ Berlin.\ Akad.\ (1859).
\bibitem{Schafer66} R.~D.~Schafer, \emph{An Introduction to Nonassociative Algebras},
Academic Press, 1966.
\bibitem{Titchmarsh86} E.~C.~Titchmarsh (rev.\ D.~R.~Heath-Brown), \emph{The Theory of
the Riemann Zeta-Function}, 2nd ed., Oxford University Press, 1986.
\bibitem{Wang} X.~Wang, \emph{On geometric aspects of quaternionic and octonionic
slice regular functions}, arXiv:1512.01414v3.
\bibitem{VS} G.~Gentili, J.~Prezelj, and F.~Vlacci, \emph{Slice conformality and
Riemann manifolds on quaternions and octonions}, Math.\ Z.\ \textbf{302} (2022),
971--994 (arXiv:2107.07892v2).
\bibitem{Stacks} The Stacks Project Authors, \emph{The Stacks Project},
\url{https://stacks.math.columbia.edu}, 2024.
\bibitem{Mathlib} The mathlib Community, \emph{The Lean mathematical library},
\texttt{CategoryTheory.ConnectedComponents}
(\texttt{typeToCatHomEquiv}) and \texttt{CategoryTheory.Limits.HasLimits}
(\texttt{colimit.desc} and \texttt{colimit.}$\iota$\texttt{\_desc}); see also
\texttt{CategoryTheory.ActionCategory}, \texttt{CategoryTheory.Grothendieck},
and \texttt{CategoryTheory.Limits.Types.Colimits},
\url{https://leanprover-community.github.io/mathlib4_docs/}.
\bibitem{GJ} P.~G.~Goerss and J.~F.~Jardine, \emph{Simplicial Homotopy Theory}, Modern
Birkh\"auser Classics, Birkh\"auser, 2009 (Ch.~I: simplicial sets and the nerve; Ch.~IV: the
bisimplicial engine for homotopy colimits).
\bibitem{GPVwind} G.~Gentili, J.~Prezelj, and F.~Vlacci, \emph{On a continuation of
quaternionic and octonionic logarithm along curves and the winding number},
J.~Math.\ Anal.\ Appl.\ \textbf{536} (2024), no.~1, Paper No.~128219,
DOI 10.1016/j.jmaa.2024.128219 (arXiv:2307.14047) (Rem.~2.1: the direction $I(q)$ has no
continuous extension to $\RR$; Cor.~4.13: every branch of the logarithm along a companioned
path is $\log|\gamma|+\mathrm{Arg}_{2k}$ --- the real part is branch-independent, the one
real transport; Def.~4.7: tame path = unique companion (Def.~4.20 for maps;
Def.~5.2: tame/semi-tame at an obstruction parameter); Def.~5.11: the loop lift
$\mathrm{pr}_1\circ\Gamma=\gamma\circ\exp$; Cor.~5.13: a lift of a loop exists iff the
signature lies in $\{0,-1\}$ on each obstruction interval, and is then a loop; Cor.~5.21:
winding $=|\sigma^c|/2$).
\bibitem{Quillen73} D.~Quillen, \emph{Higher algebraic $K$-theory: I}, in \emph{Algebraic
$K$-theory, I}, Lecture Notes in Math.\ \textbf{341}, Springer, 1973, 85--147 (Theorem~A:
a functor whose comma categories are contractible is final and induces an equivalence of
classifying spaces).
\bibitem{Riehl} E.~Riehl, \emph{Categorical Homotopy Theory}, New Mathematical Monographs \textbf{24},
Cambridge University Press, 2014, DOI 10.1017/CBO9781107261457
(\S8.3, Lem.~8.3.4 and Rem.~8.3.5: the category-of-elements computation,
$\pi_0\dashv\operatorname{Disc}$, finite zigzags, and the singleton characterization of a
nonempty connected category; see also Ch.~1 and 4--6 for the nerve, the two-sided bar construction,
classifying spaces, and homotopy colimits).
\bibitem{RiehlContext} E.~Riehl, \emph{Category Theory in Context}, Aurora: Dover Modern Math
Originals, Dover Publications, 2016 (Ex.~1.5.19: action groupoids and their orbits;
\S3.1, Defs.~3.1.1 and 3.1.5--3.1.6: constant diagrams and the colimit as a universal cocone).
\bibitem{Thomason79} R.~W.~Thomason, \emph{Homotopy colimits in the category of small
categories}, Math.\ Proc.\ Cambridge Philos.\ Soc.\ \textbf{85} (1979), 91--109 (Thm.~1.2:
the classifying space of the Grothendieck construction as a homotopy colimit).
\bibitem{Vistoli} A.~Vistoli, \emph{Notes on Grothendieck topologies, fibered categories and descent
theory}, in \emph{Fundamental Algebraic Geometry}, Math.\ Surveys Monogr.\ \textbf{123}, AMS, 2005
(arXiv:math/0412512); \S3.1 for the Grothendieck construction and categories fibered in groupoids.
\end{thebibliography}

\end{document}
